% Contraction de texte : le résumé (reformuler les idées d’un texte à résumer) — Français 1re Tech — Cours signé OnBuch+
% Programme officiel MINESEC. Compiler avec : tectonic N13-contraction-texte-resume-reformuler-idees-texte-resumer.tex
\newcommand{\DOCMATIERE}{Français}
\newcommand{\DOCNIVEAU}{1re Tech}
\newcommand{\DOCMODULE}{Français – classe de Première technique}
\newcommand{\DOCLECON}{N13}
\newcommand{\DOCTITRE}{Contraction de texte : le résumé (reformuler les idées d’un texte à résumer)}
% OnBuch+ — préambule « cours signés » ENRICHI (compilé avec tectonic / XeLaTeX)
% Système visuel « Pop » d'OnBuch+ : fond crème (jamais blanc), bordures encre
% épaisses, ombres « dures » décalées, titres Archivo Black en capitales.
% Version MATHÉMATIQUES : graphes (pgfplots/tikz), boîtes pédagogiques
% (définition, propriété, méthode, attention, le savais-tu, exercice résolu,
% exercice, TP, bilan), page de garde et signature OnBuch+ ancrée en bas.
%
% Macros attendues dans main.tex AVANT \input{preamble} :
%   \DOCMATIERE  \DOCNIVEAU  \DOCMODULE  \DOCLECON (numéro)  \DOCTITRE
\documentclass[11pt]{article}

\usepackage{fontspec}
\usepackage[french]{babel}
\usepackage[a4paper,margin=2cm,top=3.4cm,bottom=2.8cm,headheight=1.4cm,footskip=1.3cm]{geometry}
\usepackage{amsmath,amssymb,amsfonts}
\usepackage{mathtools} % \xmapsto, \coloneqq... souvent utilisés par les modèles
\usepackage{cancel} % simplifications barrées (bilans de forces)
% \abs{z} (module, valeur absolue) et \norm{u} (norme), utilisés par les modèles
\providecommand{\abs}{}\let\abs\relax\DeclarePairedDelimiter{\abs}{\lvert}{\rvert}
\providecommand{\norm}{}\let\norm\relax\DeclarePairedDelimiter{\norm}{\lVert}{\rVert}
\usepackage[table]{xcolor} % [table] : fournit aussi \rowcolors (alternance de lignes)
\usepackage{pagecolor}
\usepackage{tikz}
\usetikzlibrary{arrows.meta,positioning,calc,shapes.geometric,shapes.misc,shapes.arrows,decorations.pathmorphing,decorations.pathreplacing,decorations.markings,patterns,shadows,mindmap,trees,fit,backgrounds,matrix,intersections,angles,quotes}
\usepackage{pgfplots}
\usepackage[version=4]{mhchem} % équations nucléaires \ce{^{235}_{92}U}
\usepackage[european,siunitx]{circuitikz} % schémas électriques
\pgfplotsset{compat=1.18}
\usepgfplotslibrary{fillbetween,groupplots} % aires entre courbes (calcul intégral)
\usepackage{enumitem}
\usepackage{fancyhdr}
% [most] charge toutes les bibliothèques tcolorbox (listings, minted, raster,
% documentation...) et gonfle inutilement la mémoire TeX sur un document long
% (~300 boîtes) : on ne charge que ce qui est réellement utilisé.
\usepackage[skins,breakable]{tcolorbox}
\usepackage{graphicx}
\usepackage{colortbl}
\usepackage{booktabs}
\usepackage{tabularx}
\usepackage{multirow}
\usepackage{diagbox} % case d'en-tête diagonale des tableaux à double entrée
% Colonnes extensibles centrée / gauche / droite (C, L, R), souvent utilisées par les modèles
\newcolumntype{C}{>{\centering\arraybackslash}X}
\newcolumntype{L}{>{\raggedright\arraybackslash}X}
\newcolumntype{R}{>{\raggedleft\arraybackslash}X}
\usepackage{array}
\usepackage{multicol}
\usepackage{siunitx}
% parse-numbers=false : accepte aussi \SI{2,5 \times 10^{-3}}{...}
\sisetup{locale=FR,output-decimal-marker={,},group-minimum-digits=4,parse-numbers=false}
\DeclareSIUnit\ohm{\ensuremath{\symup{\Omega}}} % U+2126 absent de la police
\DeclareSIUnit\division{div} % réglages d'oscilloscope (ms/div, V/div)
\DeclareSIUnit\tr{tr} % tours (vitesse de rotation, ex. \SI{1500}{\tr\per\minute})
\DeclareSIUnit\molar{M} % concentration molaire (ex. \SI{3}{\molar}, \SI{1}{\micro\molar})
\DeclareSIUnit\an{an} % année (le modèle confond parfois avec \year, un compteur TeX) — ex. \SI{5}{\kilo\meter\per\an}
\DeclareSIUnit\FCFA{FCFA} % franc CFA (ex. \SI{500}{\FCFA\per\kilo\gram})
\DeclareSIUnit\degreeBrix{\textdegree Brix} % teneur en sucre d'un sirop/jus (réfractométrie)
\DeclareSIUnit\month{mois} % le modèle utilise parfois \month, un compteur TeX, comme unité de durée
\DeclareSIUnit\microliter{\micro\liter} % le modèle écrit parfois \microliter en un seul token
\DeclareSIUnit\million{millions} % dénombrement (ex. \SI{15}{\million\per\mL} spermatozoïdes)
\usepackage{qrcode}
% \degree : commande fréquemment utilisée par les modèles pour un angle ou une
% température en dehors de \SI{}{} (non fournie par les paquets chargés).
\providecommand{\degree}{^{\circ}}
% \qed (fin de démonstration), utilisé par les modèles sans amsthm
\providecommand{\qed}{\hfill$\square$}
% \barre : double barre des valeurs interdites dans un tableau de variations (tkz-tab)
\providecommand{\barre}{\ensuremath{\|}}
% environnement proof (amsthm non chargé) utilisé par les modèles
\ifdefined\proof\else\newenvironment{proof}[1][Démonstration]{\par\noindent\textit{#1.}\ }{\qed\par}\fi
\usepackage{lastpage}
\usepackage{hyperref}
\hypersetup{hidelinks}

% ============================================================
% Palette « Pop » OnBuch+ — mêmes hex que la classe P (plus_ui.dart)
% ============================================================
\definecolor{poporange}{HTML}{F59321}
\definecolor{popdark}{HTML}{E07A0C}
\definecolor{popcream}{HTML}{FFF6E8}
\definecolor{popink}{HTML}{1C1714}
\definecolor{poppurple}{HTML}{7A5AE0}
\definecolor{popgreen}{HTML}{2E9E6A}
\definecolor{poppink}{HTML}{E0457B}
\definecolor{popblue}{HTML}{2F7DE1}
\definecolor{popgold}{HTML}{C98A12}
\definecolor{popmuted}{HTML}{8A7F76}
\definecolor{popline}{HTML}{EADFD2}
\definecolor{popgray}{HTML}{8A7F76}
\colorlet{popgrey}{popgray}
% Alias tolérants (le modèle nomme parfois les couleurs différemment)
\colorlet{popwhite}{white}
\colorlet{popblack}{popink}
\colorlet{popred}{poppink}
\colorlet{popyellow}{popgold}
% Teintes claires (fonds de tableaux/encadrés) — le modèle nomme parfois
% les couleurs avec un suffixe L (ex. popblueL) sans les avoir définies.
\colorlet{poporangeL}{poporange!20}
\colorlet{popdarkL}{popdark!20}
\colorlet{poppurpleL}{poppurple!20}
\colorlet{popgreenL}{popgreen!20}
\colorlet{poppinkL}{poppink!20}
\colorlet{popblueL}{popblue!20}
\colorlet{popgoldL}{popgold!20}
\colorlet{popredL}{popred!20}
\colorlet{popgrayL}{popgray!20}
% Teintes claires pour fonds de tableaux / zones
\colorlet{popblueL}{popblue!10!white}
\colorlet{poporangeL}{poporange!14!white}
\colorlet{popgreenL}{popgreen!12!white}
\colorlet{poppurpleL}{poppurple!12!white}
\colorlet{poppinkL}{poppink!12!white}
\colorlet{popgoldL}{popgold!14!white}

\pagecolor{popcream}

% ============================================================
% Typographie
% ============================================================
\defaultfontfeatures{Ligatures=TeX}
\setmainfont{PlusJakartaSans-Medium.ttf}[Path=../../../fonts/,BoldFont=PlusJakartaSans-Bold.ttf]
% Maths en sans-serif (Fira Math) avec chiffres et lettres droites de Plus
% Jakarta, pour que les formules se fondent dans le texte.
\usepackage[math-style=ISO,bold-style=ISO]{unicode-math}
\setmathfont{FiraMath-Regular.otf}[Path=../../../fonts/]
\setmathfont{PlusJakartaSans-Medium.ttf}[Path=../../../fonts/,range={up/num,up/{Latin,latin},\mathup}]
\usepackage{icomma} % virgule décimale sans espace en mode math
% Caractères mathématiques Unicode absents des polices de texte (Plus Jakarta,
% Archivo Black) : rendus par leur équivalent mathématique, en texte comme en
% formule — sinon ils s'affichent en carré vide (ex. « Arithmétique dans ℤ »).
\usepackage{newunicodechar}
\newunicodechar{ℝ}{\ensuremath{\mathbb{R}}}
\newunicodechar{ℤ}{\ensuremath{\mathbb{Z}}}
\newunicodechar{ℂ}{\ensuremath{\mathbb{C}}}
\newunicodechar{ℕ}{\ensuremath{\mathbb{N}}}
\newunicodechar{ℚ}{\ensuremath{\mathbb{Q}}}
\newunicodechar{𝔻}{\ensuremath{\mathbb{D}}}
\newunicodechar{α}{\ensuremath{\alpha}}
\newunicodechar{β}{\ensuremath{\beta}}
\newunicodechar{γ}{\ensuremath{\gamma}}
\newunicodechar{δ}{\ensuremath{\delta}}
\newunicodechar{ε}{\ensuremath{\varepsilon}}
\newunicodechar{θ}{\ensuremath{\theta}}
\newunicodechar{λ}{\ensuremath{\lambda}}
\newunicodechar{μ}{\ensuremath{\mu}}
\newunicodechar{σ}{\ensuremath{\sigma}}
\newunicodechar{τ}{\ensuremath{\tau}}
\newunicodechar{φ}{\ensuremath{\varphi}}
\newunicodechar{ψ}{\ensuremath{\psi}}
\newunicodechar{ω}{\ensuremath{\omega}}
\newunicodechar{Σ}{\ensuremath{\Sigma}}
% majuscules produites par \MakeUppercase dans les titres (α-aminés -> Α)
\newunicodechar{Α}{\ensuremath{\alpha}}
\newunicodechar{Β}{\ensuremath{\beta}}
\newunicodechar{Γ}{\ensuremath{\Gamma}}
\newunicodechar{Δ}{\ensuremath{\Delta}}
\newunicodechar{Ε}{\ensuremath{\varepsilon}}
\newunicodechar{Θ}{\ensuremath{\Theta}}
\newunicodechar{Λ}{\ensuremath{\Lambda}}
\newunicodechar{Μ}{\ensuremath{\mu}}
\newunicodechar{Τ}{\ensuremath{\tau}}
\newunicodechar{Φ}{\ensuremath{\Phi}}
\newunicodechar{Ψ}{\ensuremath{\Psi}}
\newunicodechar{Ω}{\ensuremath{\Omega}}
\newunicodechar{↦}{\ensuremath{\mapsto}}
\newunicodechar{∈}{\ensuremath{\in}}
\newunicodechar{∉}{\ensuremath{\notin}}
\newunicodechar{∧}{\ensuremath{\wedge}}
\newunicodechar{⇔}{\ensuremath{\Leftrightarrow}}
\newunicodechar{⟺}{\ensuremath{\Longleftrightarrow}}
\newunicodechar{⇒}{\ensuremath{\Rightarrow}}
\newunicodechar{⟹}{\ensuremath{\Longrightarrow}}
\newunicodechar{∘}{\ensuremath{\circ}}
\newunicodechar{∪}{\ensuremath{\cup}}
\newunicodechar{∩}{\ensuremath{\cap}}
\newunicodechar{≡}{\ensuremath{\equiv}}
\newunicodechar{⊂}{\ensuremath{\subset}}
\newunicodechar{⊥}{\ensuremath{\perp}}
\newunicodechar{∃}{\ensuremath{\exists}}
\newunicodechar{∀}{\ensuremath{\forall}}
\newunicodechar{∛}{\ensuremath{\sqrt[3]{\,}}}
\newunicodechar{∜}{\ensuremath{\sqrt[4]{\,}}}
\newunicodechar{⃗}{\ensuremath{{}^{\to}}}
\newunicodechar{ⁿ}{\ifmmode^{n}\else\textsuperscript{\ensuremath{n}}\fi}
\newunicodechar{ˣ}{\ifmmode^{x}\else\textsuperscript{\ensuremath{x}}\fi}
\newunicodechar{ᵇ}{\ifmmode^{b}\else\textsuperscript{\ensuremath{b}}\fi}
\newunicodechar{ʳ}{\ifmmode^{r}\else\textsuperscript{\ensuremath{r}}\fi}
\newunicodechar{ᵘ}{\ifmmode^{u}\else\textsuperscript{\ensuremath{u}}\fi}
\newunicodechar{ʸ}{\ifmmode^{y}\else\textsuperscript{\ensuremath{y}}\fi}
\newunicodechar{ᵃ}{\ifmmode^{a}\else\textsuperscript{\ensuremath{a}}\fi}
\newunicodechar{ᵉ}{\ifmmode^{e}\else\textsuperscript{\ensuremath{e}}\fi}
\newunicodechar{ᵏ}{\ifmmode^{k}\else\textsuperscript{\ensuremath{k}}\fi}
\newunicodechar{ᵖ}{\ifmmode^{p}\else\textsuperscript{\ensuremath{p}}\fi}
\newunicodechar{ᴺ}{\ifmmode^{N}\else\textsuperscript{\ensuremath{N}}\fi}
\newunicodechar{ᶜ}{\ifmmode^{c}\else\textsuperscript{\ensuremath{c}}\fi}
\newunicodechar{ʰ}{\ifmmode^{h}\else\textsuperscript{\ensuremath{h}}\fi}
\newunicodechar{ᵠ}{\ifmmode^{q}\else\textsuperscript{\ensuremath{q}}\fi}
\newunicodechar{ₙ}{\ifmmode_{n}\else\textsubscript{\ensuremath{n}}\fi}
\newunicodechar{ₐ}{\ifmmode_{a}\else\textsubscript{\ensuremath{a}}\fi}
\newunicodechar{ᵢ}{\ifmmode_{i}\else\textsubscript{\ensuremath{i}}\fi}
\newunicodechar{ᵣ}{\ifmmode_{r}\else\textsubscript{\ensuremath{r}}\fi}
\newunicodechar{ₖ}{\ifmmode_{k}\else\textsubscript{\ensuremath{k}}\fi}
\newunicodechar{ᵦ}{\ifmmode_{\beta}\else\textsubscript{\ensuremath{\beta}}\fi}
\newunicodechar{ₓ}{\ifmmode_{x}\else\textsubscript{\ensuremath{x}}\fi}
\newfontfamily\popdisplay{ArchivoBlack-Regular.ttf}[Path=../../../fonts/]
\newfontfamily\poplogo{SpaceGrotesk-Bold.ttf}[Path=../../../fonts/]
\newfontfamily\popsemi{PlusJakartaSans-SemiBold.ttf}[Path=../../../fonts/]
\newfontfamily\popxbold{PlusJakartaSans-ExtraBold.ttf}[Path=../../../fonts/]
\newfontfamily\popxboldls{PlusJakartaSans-ExtraBold.ttf}[Path=../../../fonts/,LetterSpace=3.0]

\setlist[enumerate]{leftmargin=1.7em,itemsep=5pt,topsep=4pt}
\setlist[itemize]{leftmargin=1.7em,itemsep=4pt,topsep=4pt,label={\color{poporange}\textbullet}}
\setlength{\parindent}{0pt}
\setlength{\parskip}{5pt}
\renewcommand{\arraystretch}{1.25}

% pgfplots : style Pop par défaut
\pgfplotsset{
  every axis/.append style={axis line style={popink,line width=1pt},tick style={popink},
    grid style={popline,line width=0.5pt},label style={font=\small},tick label style={font=\footnotesize}},
  popcurve/.style={poporange,line width=1.8pt,no markers,smooth},
}
% Même style disponible dans un \draw TikZ simple (hors pgfplots)
\tikzset{popcurve/.style={poporange,line width=1.8pt}}
\tikzset{popgrid/.style={popline,very thin}}
\colorlet{popcurve}{poporange} % le nom du style est parfois utilisé comme couleur

% ============================================================
% Pill / squiggle
% ============================================================
\newcommand{\poppill}[3]{%
  \tikz[baseline=-0.55ex]\node[draw=popink,line width=1.2pt,rounded corners=2.6pt,%
    fill=#2,inner xsep=6.5pt,inner ysep=3pt]{%
    \popxboldls\fontsize{7}{7}\selectfont\color{#3}\MakeUppercase{#1}};%
}
\newcommand{\popsquiggle}[1]{%
  \tikz[baseline=-0.4ex]{
    \draw[#1,line width=2.6pt,line cap=round]
      plot [smooth] coordinates {(0,0.09) (0.34,0.01) (0.68,0.21) (1.02,0.01) (1.36,0.10)};
  }%
}
% Mot-clé surligné (marqueur orange sous le texte)
\newcommand{\cle}[1]{{\popxbold\color{popdark}#1}}

% ============================================================
% En-tête / pied
% ============================================================
\pagestyle{fancy}
\fancyhf{}
\renewcommand{\headrulewidth}{0pt}
\renewcommand{\footrulewidth}{0pt}
\lhead{\raisebox{-0.15ex}{\poplogo\fontsize{13}{13}\selectfont\color{popink}OnBuch\color{poporange}+}}
\rhead{\poppill{\DOCMATIERE\ \textbullet\ \DOCNIVEAU\ \textbullet\ Leçon \DOCLECON}{white}{popink}}
\lfoot{\popxbold\fontsize{7.6}{7.6}\selectfont\color{popmuted}\MakeUppercase{Cours signé OnBuch+}\ \textbullet\ {\popsemi\DOCTITRE}}
\rfoot{\tikz[baseline=-0.6ex]\node[rounded corners=8.5pt,fill=popink,minimum height=17pt,minimum width=17pt,inner xsep=5pt,inner sep=0]{\popxbold\fontsize{8}{8}\selectfont\color{popcream}\thepage\,/\,\pageref*{LastPage}};}

% ============================================================
% Titres
% ============================================================
\newcommand{\chaptitre}[2]{%
  \begin{tcolorbox}[enhanced,colback=poporange,colframe=popink,boxrule=2.6pt,arc=11pt,
      drop shadow=popink,left=15pt,right=15pt,top=12pt,bottom=11pt,before skip=3pt,after skip=18pt]
    \poppill{#2}{popink}{popcream}\par\vspace{9pt}
    {\raggedright\hyphenpenalty=10000\exhyphenpenalty=10000\popdisplay\fontsize{22}{24}\selectfont\color{popcream}\MakeUppercase{#1}\par}
  \end{tcolorbox}
}
% Section numérotée : pastille orange avec le numéro + titre Archivo
\newcounter{popsec}
\newcommand{\coursec}[1]{%
  \stepcounter{popsec}\setcounter{popsub}{0}%
  \par\vspace{16pt}\noindent
  \begin{minipage}{\linewidth}
  \tikz[baseline=-0.7ex]\node[draw=popink,line width=1.6pt,fill=poporange,rounded corners=4pt,minimum size=22pt,inner sep=0]{\popdisplay\fontsize{12}{12}\selectfont\color{popcream}\thepopsec};%
  \hspace{8pt}{\popdisplay\fontsize{15}{17}\selectfont\color{popink}\MakeUppercase{#1}}\par
  \vspace{4pt}\noindent{\color{poporange}\rule{36pt}{3.6pt}}
  \end{minipage}\par\nopagebreak\vspace{8pt}
}
\newcounter{popsub}
\newcommand{\courssub}[1]{%
  \stepcounter{popsub}%
  \par\vspace{9pt}\noindent{\popxbold\fontsize{11.5}{13}\selectfont\color{popink}{\color{poporange}\thepopsec.\thepopsub}\hspace{6pt}#1}\par\nopagebreak\vspace{4pt}
}

% ============================================================
% Boîtes pédagogiques — même recette : fond blanc, bordure épaisse colorée,
% ombre dure encre, badge plein. Titre optionnel [#1].
% ============================================================
\tcbset{popbase/.style={breakable,enhanced,colback=white,boxrule=2.3pt,arc=9pt,
  shadow={3pt}{-3pt}{0pt}{popink},left=14pt,right=14pt,top=11pt,bottom=11pt,before skip=11pt,after skip=15pt,
  fontupper=\popsemi\fontsize{10.6}{14.5}\selectfont\color{popink},
  fontlower=\popsemi\fontsize{10.6}{14.5}\selectfont\color{popink}}}
% Coche dessinée (le glyphe ✓ n'existe pas dans les polices du document)
\newcommand{\popcheck}[1]{\tikz[baseline=-0.5ex]\draw[#1,line width=1.6pt,line cap=round,line join=round](-0.13,0.0)--(-0.04,-0.09)--(0.13,0.1);}
\providecommand{\coche}{\popcheck{popgreen}}
\providecommand{\resultat}[1]{\par\smallskip\noindent\fcolorbox{popgreen}{popgreenL}{\parbox{\dimexpr\linewidth-2\fboxsep-2\fboxrule\relax}{\textbf{Résultat.} #1}}\par\smallskip}
\newcommand{\popboxhead}[3]{\poppill{#1}{#2}{white}\ifx\relax#3\relax\else\hspace{6pt}{\popxbold\fontsize{10.6}{12}\selectfont\color{popink}#3}\fi\par\vspace{7pt}}

\newtcolorbox{aretenir}[1][]{popbase,colframe=popgreen,before upper={\popboxhead{À retenir}{popgreen}{#1}}}
\newtcolorbox{exemplebox}[1][]{popbase,colframe=poporange,before upper={\popboxhead{Exemple}{poporange}{#1}}}
\newtcolorbox{definition}[1][]{popbase,colframe=popblue,before upper={\popboxhead{Définition}{popblue}{#1}}}
\newtcolorbox{propriete}[1][]{popbase,colframe=poppurple,before upper={\popboxhead{Propriété}{poppurple}{#1}}}
\newtcolorbox{methode}[1][]{popbase,colframe=popgold,before upper={\popboxhead{Méthode}{popgold}{#1}}}
\newtcolorbox{attention}[1][]{popbase,colframe=poppink,before upper={\popboxhead{Attention}{poppink}{#1}}}
\newtcolorbox{experience}[1][]{popbase,colframe=popblue,colback=popblueL,before upper={\popboxhead{Expérience}{popblue}{#1}}}
% « Le savais-tu ? » : carte violette pleine (contexte, culture, Cameroun)
\newtcolorbox{savaistu}[1][]{popbase,colback=poppurple,colframe=popink,
  fontupper=\popsemi\fontsize{10.4}{14.2}\selectfont\color{white},
  before upper={\poppill{Le savais-tu\,?}{white}{poppurple}\ifx\relax#1\relax\else\hspace{6pt}{\popxbold\color{white}#1}\fi\par\vspace{7pt}}}
% Exercice résolu : énoncé en haut, \tcblower puis la solution sur fond crème
\newcounter{popexo}\newcounter{popexr}
\newtcolorbox{exoresolu}[1][]{popbase,colframe=poporange,colbacklower=popcream,
  segmentation style={popink,line width=1.2pt,dashed},
  before upper={\stepcounter{popexr}\popboxhead{Exercice résolu \thepopexr}{poporange}{#1}},
  before lower={\poppill{Solution}{popink}{popcream}\par\vspace{6pt}}}
% Exercice (non résolu en place) — la correction est dans la partie Corrigés
\newtcolorbox{exercice}[1][]{popbase,colframe=popink,
  before upper={\stepcounter{popexo}\popboxhead{Exercice \thepopexo}{popink}{#1}}}
% Corrigé
\newtcolorbox{corrige}[1][]{popbase,colframe=popgreen,colback=popgreenL,
  before upper={\popboxhead{Corrigé}{popgreen}{#1}}}
% Objectifs / prérequis / bilan
\newtcolorbox{objectifs}{popbase,colframe=popink,colback=poporangeL,
  before upper={\popboxhead{Objectifs de la leçon}{popink}{}}}
\newtcolorbox{prerequis}{popbase,colframe=popink,colback=popblueL,
  before upper={\popboxhead{Prérequis}{popblue}{}}}
\newtcolorbox{bilan}{popbase,colframe=popink,colback=white,boxrule=2.8pt,
  before upper={\popboxhead{Fiche bilan}{poporange}{L'essentiel à connaître pour l'examen}}}
% Figure encadrée avec légende
\newtcolorbox{popfigure}[1][]{enhanced,breakable=false,colback=white,colframe=popline,boxrule=1.4pt,arc=8pt,
  left=10pt,right=10pt,top=10pt,bottom=8pt,before skip=10pt,after skip=12pt,halign=center,
  lower separated=false,halign lower=center,
  fontlower=\popsemi\fontsize{9}{11.5}\selectfont\color{popmuted},
  #1}
\newcommand{\legende}[1]{\par\vspace{6pt}{\popsemi\fontsize{9}{11.5}\selectfont\color{popmuted}#1}}

% ============================================================
% Signature OnBuch+
% ============================================================
\newcommand{\poplogoinline}[1]{{\poplogo\fontsize{#1}{#1}\selectfont\color{popink}OnBuch\color{poporange}+}}
\newcommand{\signatureonbuch}{%
  \begin{tcolorbox}[enhanced,colback=popink,colframe=popink,boxrule=2.2pt,arc=10pt,
      drop shadow=poporange,left=14pt,right=14pt,top=10pt,bottom=10pt,before skip=0pt,after skip=0pt]
    \tikz[baseline=-0.7ex]\node[circle,fill=popgreen,draw=popcream,line width=1.3pt,minimum size=22pt,inner sep=0]{\popcheck{white}};%
    \hspace{9pt}\begin{minipage}[c]{0.62\linewidth}
      {\popxbold\fontsize{12}{13}\selectfont\color{popcream}Signé\ \textbullet\ {\poplogo OnBuch\color{poporange}+}}\par\vspace{2pt}
      {\raggedright\popsemi\fontsize{8}{10.5}\selectfont\color{popline}\MakeUppercase{Équipe pédagogique OnBuch+}\\\MakeUppercase{Conforme au programme officiel MINESEC}\par}
    \end{minipage}\hfill
    {\poplogo\fontsize{22}{22}\selectfont\color{popcream}OnBuch\color{poporange}+}
  \end{tcolorbox}%
}

% ============================================================
% Page de garde
% #1 titre · #2 matière · #3 niveau · #4 module · #5 n° leçon · #6 durée ·
% #7 description / objectifs (texte ou itemize)
% ============================================================
\newcommand{\pagegarde}[7]{%
  \thispagestyle{empty}
  \newgeometry{a4paper,margin=1.7cm,top=1.6cm,bottom=1.4cm}%
  \begin{tikzpicture}[remember picture,overlay]
    \fill[poporange!18] ([xshift=-3.2cm,yshift=-3.6cm]current page.north east) circle (4.6cm);
    \fill[poppurple!14] ([xshift=2.4cm,yshift=7.6cm]current page.south west) circle (3.8cm);
  \end{tikzpicture}%
  {\poplogo\fontsize{30}{30}\selectfont\color{popink}OnBuch\color{poporange}+}\hfill
  \raisebox{0.6ex}{\poppill{Cours signé \textbullet\ Premium}{popink}{popcream}}\par
  \vspace{1pt}\popsquiggle{poppurple}\par
  \vspace{8pt}
  \begin{tcolorbox}[enhanced,colback=poporange,colframe=popink,boxrule=2.8pt,arc=13pt,
      drop shadow=popink,left=18pt,right=18pt,top=12pt,bottom=13pt,after skip=10pt]
    \poppill{#2 \textbullet\ #3}{popink}{popcream}\hspace{5pt}\poppill{#4}{white}{popink}\par\vspace{8pt}
    {\popdisplay\fontsize{14}{14}\selectfont\color{popink}LEÇON #5}\par\vspace{4pt}
    {\raggedright\hyphenpenalty=10000\exhyphenpenalty=10000\popdisplay\fontsize{25}{27}\selectfont\color{popcream}\MakeUppercase{#1}\par}
    \vspace{8pt}
    {\popxbold\fontsize{9.5}{9.5}\selectfont\color{popink}\MakeUppercase{Durée conseillée : #6}}
  \end{tcolorbox}
  \begin{tcolorbox}[enhanced,colback=white,colframe=popink,boxrule=2.2pt,arc=10pt,
      drop shadow=popink,left=16pt,right=16pt,top=10pt,bottom=10pt,after skip=8pt]
    \poppill{Dans cette leçon}{poppurple}{white}\par\vspace{5pt}
    \popsemi\fontsize{9.3}{12.6}\selectfont\color{popink}#7
  \end{tcolorbox}
  \noindent\begin{minipage}[t]{0.487\linewidth}
    \begin{tcolorbox}[enhanced,colback=popgreen,colframe=popink,boxrule=2.2pt,arc=10pt,
        drop shadow=popink,left=11pt,right=11pt,top=10pt,bottom=10pt,height=3.1cm,valign=center]
      \tcbox[colback=white,colframe=popink,boxrule=1.3pt,arc=5pt,boxsep=0pt,nobeforeafter,
        left=3pt,right=3pt,top=3pt,bottom=3pt,valign=center]{\qrcode[height=1.9cm]{https://play.google.com/store/apps/details?id=com.luvvix.onbuch&hl=fr}}%
      \hspace{8pt}\begin{minipage}[c]{0.5\linewidth}
        {\popdisplay\fontsize{11}{12.5}\selectfont\color{white}\MakeUppercase{Télécharge OnBuch}}\par\vspace{4pt}
        {\raggedright\popsemi\fontsize{8.4}{11}\selectfont\color{white}Gratuit \textbullet\ Google Play\\Tuteur IA, annales, concours, résultats\par}
      \end{minipage}
    \end{tcolorbox}
  \end{minipage}\hfill
  \begin{minipage}[t]{0.487\linewidth}
    \begin{tcolorbox}[enhanced,colback=poppurple,colframe=popink,boxrule=2.2pt,arc=10pt,
        drop shadow=popink,left=11pt,right=11pt,top=10pt,bottom=10pt,height=3.1cm,valign=center]
      \tikz[baseline=-0.7ex]\node[circle,fill=white,draw=popink,line width=1.3pt,minimum size=1.6cm,inner sep=0]{\popdisplay\fontsize{24}{24}\selectfont\color{poppurple}+};%
      \hspace{8pt}\begin{minipage}[c]{0.6\linewidth}
        {\popdisplay\fontsize{11}{12.5}\selectfont\color{white}\MakeUppercase{Passe à OnBuch+}}\par\vspace{4pt}
        {\raggedright\popsemi\fontsize{8.4}{11}\selectfont\color{white}Tous les cours signés, Tuteur IA illimité, fiches PDF — onglet OnBuch+ de l'app\par}
      \end{minipage}
    \end{tcolorbox}
  \end{minipage}
  \vspace{8pt}
  \popcontenu
  \vfill
  \signatureonbuch
  \restoregeometry
  \newpage
}

% Rangée « Ce cours contient » (page de garde)
\newcommand{\poptile}[3]{% #1 couleur #2 gros texte #3 légende
  \begin{tcolorbox}[enhanced,colback=#1,colframe=popink,boxrule=2pt,arc=9pt,shadow={2.5pt}{-2.5pt}{0pt}{popink},
      left=6pt,right=6pt,top=9pt,bottom=9pt,halign=center,valign=center,height=1.75cm,width=\linewidth,nobeforeafter]
    {\popdisplay\fontsize{12.5}{13}\selectfont\color{white}\MakeUppercase{#2}}\par\vspace{3pt}
    {\centering\popsemi\fontsize{7.8}{9.5}\selectfont\color{white}#3\par}
  \end{tcolorbox}}
\newcommand{\popcontenu}{%
  \par\noindent{\popxboldls\fontsize{7.5}{7.5}\selectfont\color{popmuted}CE COURS SIGNÉ CONTIENT}\par\vspace{7pt}
  \noindent\begin{minipage}[t]{0.235\linewidth}\poptile{popblue}{Cours}{complet, illustré, pas à pas}\end{minipage}\hfill
  \begin{minipage}[t]{0.235\linewidth}\poptile{poporange}{Méthodes}{et exercices résolus}\end{minipage}\hfill
  \begin{minipage}[t]{0.235\linewidth}\poptile{poppink}{Exercices}{type examen + corrigés}\end{minipage}\hfill
  \begin{minipage}[t]{0.235\linewidth}\poptile{popgreen}{Fiche bilan}{l'essentiel à retenir}\end{minipage}\par}

% Sommaire
\newtcolorbox{sommaire}{enhanced,colback=white,colframe=popink,boxrule=2.2pt,arc=10pt,
  drop shadow=popink,left=18pt,right=18pt,top=13pt,bottom=13pt,before skip=6pt,after skip=20pt,
  before upper={\poppill{Sommaire}{poppurple}{white}\par\vspace{9pt}\popsemi\fontsize{10.6}{14.5}\selectfont\color{popink}}}

% Signature de fin de cours — poussée tout en bas de la dernière page
\newcommand{\findecours}{%
  \par\vfill\nopagebreak
  \begin{center}\popsquiggle{poporange}\end{center}\vspace{4pt}
  \signatureonbuch
}
\AtBeginDocument{\def\llbracket{\mathopen{[\mkern-3.5mu[}}\def\rrbracket{\mathclose{]\mkern-3.5mu]}}} % intervalles d'entiers (glyphe ⟦ absent de la police)
% Glyphes absents de Fira Math : redéfinis à partir de glyphes disponibles
\AtBeginDocument{%
  \def\setminus{\mathbin{\backslash}}\let\smallsetminus\setminus
  \def\vdots{\mathord{\vbox{\baselineskip3.2pt\lineskiplimit0pt\kern4pt\hbox{.}\hbox{.}\hbox{.}}}}%
  \def\ddots{\mathinner{\mkern1mu\raise6.5pt\hbox{.}\mkern2mu\raise3.6pt\hbox{.}\mkern2mu\raise0.7pt\hbox{.}\mkern1mu}}%
  \def\triangle{\mathord{\scalebox{0.9}{$\symup{\Delta}$}}}%
  \def\nabla{\mathord{\raisebox{\depth}{\scalebox{1}[-1]{$\symup{\Delta}$}}}}%
  \def\checkmark{\popcheck{popdark}}%
  \def\star{\mathbin{\ast}}\def\bigstar{\mathord{\ast}}%
  \def\diamond{\mathbin{\scalebox{0.6}{$\Diamond$}}}%
  \def\bigcup{\mathop{\mathchoice{\raisebox{-0.3em}{\scalebox{1.7}{$\cup$}}}{\scalebox{1.25}{$\cup$}}{\cup}{\cup}}\displaylimits}%
  \def\bigcap{\mathop{\mathchoice{\raisebox{-0.3em}{\scalebox{1.7}{$\cap$}}}{\scalebox{1.25}{$\cap$}}{\cap}{\cap}}\displaylimits}%
  \def\lesseqqgtr{\mathrel{\lessgtr}}\def\bigoplus{\mathop{\oplus}\displaylimits}%
}
\providecommand{\ssi}{si et seulement si} % abréviation fréquente
\AtBeginDocument{\providecommand{\wideparen}{\overparen}} % arc de cercle

%% FIGFIX-BEGIN v5 — correctifs globaux des figures (docs/onbuch-generation/tools/figfix)
\usepackage{adjustbox}\usepackage{etoolbox}\tcbuselibrary{raster}
% toute figure TikZ/pgfplots plus large que la ligne est réduite à la largeur disponible
\BeforeBeginEnvironment{tikzpicture}{\begin{adjustbox}{max width=\linewidth}}
\AfterEndEnvironment{tikzpicture}{\end{adjustbox}}
% légendes pgfplots lisibles par-dessus les courbes
\pgfplotsset{every axis/.append style={legend style={fill=white,fill opacity=0.92,text opacity=1,draw=black!15,inner sep=3pt}}}
% étiquettes d'arêtes (syntaxe edge["..."]) avec fond blanc
\tikzset{every edge quotes/.append style={fill=white,inner sep=1.2pt}}
%% FIGFIX-END
\begin{document}

\pagegarde{Contraction de texte : le résumé (reformuler les idées d’un texte à résumer)}{Français}{1re Tech}{Français – classe de Première technique}{N13}{12 séances (fiche de progression harmonisée)}{Cette leçon outille l'élève de Première technique pour maîtriser la contraction de texte, épreuve reine du Baccalauréat : repérage de la structure argumentative, techniques de réduction et reformulation fidèle des idées essentielles. Elle s'appuie sur l'étude méthodique du Lion et la perle de Wole Soyinka et sur des acquis de langue (figures d'amplification, lexique, néologismes) mobilisés dans des situations concrètes camerounaises.
\vspace{4pt}
\begin{multicols}{2}\small
\begin{itemize}
\item À la fin de cette leçon, je sais distinguer l'idée essentielle de l'idée secondaire dans un texte argumentatif.
\item À la fin de cette leçon, je sais dégager la structure argumentative d'un texte (thèse, arguments, exemples).
\item À la fin de cette leçon, je sais appliquer les techniques de réduction : suppression, généralisation, condensation, substitution.
\item À la fin de cette leçon, je sais reformuler les idées d'un texte avec mes propres mots sans les déformer ni ajouter mon opinion.
\item À la fin de cette leçon, je sais identifier et analyser les figures d'amplification (hyperbole, gradation) dans un texte littéraire.
\item À la fin de cette leçon, je sais distinguer le lexique commun du lexique spécialisé et reconnaître les néologismes dans un texte.
\end{itemize}\end{multicols}}

\begin{sommaire}
\begin{multicols}{2}
\begin{enumerate}
\item Lecture méthodique 4 du Lion et la perle et figures d'amplification
\item La structure argumentative du texte
\item Lecture méthodique 5 et lexique commun / lexique spécialisé
\item Techniques de réduction et de rédaction du résumé
\item Les néologismes
\item Reformuler les idées essentielles et exposés sur l'œuvre
\item Bilan, compte-rendu et remédiation
\item Activité d'intégration
\item Méthodes et exercices résolus
\item Exercices
\item Corrigés des exercices
\item Fiche bilan
\end{enumerate}
\end{multicols}
\end{sommaire}

\begin{objectifs}
\begin{itemize}
\item À la fin de cette leçon, je sais distinguer l'idée essentielle de l'idée secondaire dans un texte argumentatif.
\item À la fin de cette leçon, je sais dégager la structure argumentative d'un texte (thèse, arguments, exemples).
\item À la fin de cette leçon, je sais appliquer les techniques de réduction : suppression, généralisation, condensation, substitution.
\item À la fin de cette leçon, je sais reformuler les idées d'un texte avec mes propres mots sans les déformer ni ajouter mon opinion.
\item À la fin de cette leçon, je sais identifier et analyser les figures d'amplification (hyperbole, gradation) dans un texte littéraire.
\item À la fin de cette leçon, je sais distinguer le lexique commun du lexique spécialisé et reconnaître les néologismes dans un texte.
\item À la fin de cette leçon, je sais lire méthodiquement un extrait du Lion et la perle et en rendre compte à l'oral (exposé).
\item À la fin de cette leçon, je sais rédiger et justifier une démarche de résumé respectant la consigne de longueur (1/4 du texte).
\end{itemize}
\end{objectifs}

\begin{prerequis}
\begin{itemize}
\item Notion de résumé vue en classe de Seconde (reformulation simple, respect de la consigne de longueur).
\item Grammaire de la phrase : groupes nominaux, expansions du nom, subordonnées (classes de 4e et 3e).
\item Notions de base sur le texte argumentatif : thèse, argument, exemple (Seconde).
\item Champs lexicaux et relations de sens : synonymes, hyperonymes/hyponymes (Seconde).
\item Lecture suivie de l'œuvre intégrale Le Lion et la perle engagée depuis le début de l'année.
\end{itemize}
\end{prerequis}

\newpage

\coursec{Lecture méthodique 4 du Lion et la perle et figures d'amplification}

\courssub{Situation d'accroche et problématique}

À la veille du concours d'entrée à l'ENAM, Amina, élève de Première technique à Yaoundé, reçoit un article de trois pages du quotidien \emph{Cameroon Tribune} sur la transformation numérique de l'économie camerounaise. Elle dispose de 20 minutes pour en produire un résumé de 10 lignes sans trahir la pensée de l'auteur. Comment repérer les idées essentielles ? Comment les reformuler avec ses propres mots sans les déformer ? Cette compétence, exigée au Baccalauréat technique, est aussi celle du cadre qui synthétise un rapport pour son ministre.

Mais avant de résumer, il faut savoir lire. Et lire, c'est d'abord comprendre ce que le texte fait à son lecteur. Or un texte littéraire, comme la pièce de théâtre que nous étudions cette année, est truffé de procédés expressifs : exagérations, montées en puissance, images éclatantes. Ces procédés donnent du relief aux personnages, mais ils ne sont pas des \emph{idées} : ce sont des \emph{manières de dire}. Le bon résumeur doit donc apprendre à les reconnaître pour les neutraliser. C'est précisément l'objet de cette séance.

\begin{aretenir}[Problématique de la séance]
Comment une lecture méthodique de l'extrait permet-elle de repérer les figures d'amplification (hyperbole et gradation) dans les répliques des personnages, et pourquoi le futur résumeur doit-il savoir les reconnaître pour les écarter de sa reformulation ?
\end{aretenir}

\courssub{Rappel sur l'œuvre : Le Lion et la perle de Wole Soyinka}

Avant d'entrer dans la lecture méthodique n° 4, rappelons les repères indispensables sur l'œuvre étudiée.

\emph{Le Lion et la perle} (\emph{The Lion and the Jewel}, 1959) est une pièce de théâtre du dramaturge nigérian \cle{Wole Soyinka}. L'action se déroule dans le village imaginaire d'Ilujinle, au Nigeria. La pièce met en scène un conflit entre deux conceptions du monde :

\begin{itemize}
\item \cle{Baroka}, le baalè (chef) du village, surnommé « le Lion » : vieil homme rusé, polygame, il incarne la \cle{tradition}, la sagesse coutumière, mais aussi la ruse et la vitalité ;
\item \cle{Lakunle}, le jeune instituteur du village : il incarne la \cle{modernité}, mais une modernité superficielle, faite de mots anglais mal compris, de prétention et de refus de payer la dot ;
\item \cle{Sidi}, la plus belle jeune fille du village, surnommée « la Perle » : c'est autour d'elle que s'affrontent les deux hommes, car tous deux veulent l'épouser.
\end{itemize}

L'enjeu dramatique est donc double : un enjeu amoureux (qui épousera Sidi ?) et un enjeu culturel (la tradition ou la modernité l'emportera-t-elle ?). Soyinka, loin de prendre parti de façon simpliste, montre les forces et les faiblesses des deux camps : la ruse de Baroka triomphe, mais la tradition qu'il défend est aussi ambiguë que la modernité caricaturale de Lakunle.

\begin{savaistu}[Wole Soyinka, un géant des lettres africaines]
Wole Soyinka, né en 1934 à Abeokuta (Nigeria), est le premier Africain à avoir reçu le prix Nobel de littérature, en 1986. Dramaturge, poète, romancier et militant politique, il a écrit \emph{Le Lion et la perle} en 1959, à la veille de l'indépendance du Nigeria (1960). La pièce pose une question qui reste d'actualité dans toute l'Afrique, y compris au Cameroun : comment concilier les valeurs traditionnelles et les apports de la modernité sans perdre son âme ?
\end{savaistu}

\courssub{Lecture méthodique n° 4 : démarche et repérage}

Une lecture méthodique n'est pas une lecture rapide : c'est une lecture \emph{organisée}, qui suit des étapes précises. Pour l'extrait étudié (une scène d'affrontement verbal entre Baroka et Lakunle autour de Sidi), nous appliquons la démarche en trois temps.

\begin{methode}[Les trois étapes du repérage préalable]
\begin{enumerate}
\item \textbf{La situation d'énonciation} : qui parle ? à qui ? où ? quand ? dans quelles circonstances ? Dans notre extrait, la scène se déroule dans le village ; les répliques s'échangent devant témoins, ce qui pousse chaque personnage à « en faire trop » pour impressionner l'assistance.
\item \textbf{Les personnages en présence} : Baroka, Lakunle, Sidi, et éventuellement les chorèges ou villageois. Il faut noter qui domine l'échange, qui réplique, qui se tait.
\item \textbf{L'enjeu de l'extrait} : que cherche chaque personnage ? Baroka veut séduire et impressionner Sidi par l'éloquence ; Lakunle veut la convaincre par des arguments « modernes » et rabaisser le vieux chef. L'extrait est donc un \emph{duel de paroles}.
\end{enumerate}
\end{methode}

Ce repérage permet une première observation essentielle : dans un duel de paroles, les personnages ne parlent pas pour informer, mais pour \emph{gagner}. D'où l'usage massif de procédés qui amplifient : c'est là que naissent le comique et la poésie de la scène.

\courssub{Comique de scène et langage poétique chez Soyinka}

L'extrait étudié fait rire, et ce rire a des causes précises que la lecture méthodique doit identifier.

\textbf{Le comique de scène} naît d'abord du contraste entre les personnages : Lakunle, l'instituteur, multiplie les grands mots et les gestes emphatiques ; Baroka, lui, répond par des images tirées de la nature et de la sagesse populaire, avec une tranquille assurance. Le décalage entre la prétention de Lakunle et son impuissance réelle (il ne veut pas payer la dot, il ne fait pas le poids face au Lion) produit un comique de caractère. S'ajoute le comique de mots : Lakunle abuse de termes anglais savants que personne ne comprend au village, ce qui le ridiculise.

\textbf{Le langage poétique} de Soyinka, quant à lui, se reconnaît à la densité des images : comparaisons avec les animaux de la brousse, métaphores de la force et de la beauté, rythme des répliques qui imite parfois le chant et le proverbe yoruba. Baroka ne dit jamais simplement les choses : il les \emph{exagère}, il les \emph{grossit}, il les fait \emph{monter en puissance}. C'est exactement le domaine des figures d'amplification.

\begin{attention}[Piège de lecture]
Ne pas confondre \emph{comique} et \emph{poétique} : le comique vise à faire rire (par le ridicule, le décalage, la répétition), tandis que le langage poétique vise à émouvoir ou à admirer (par les images et le rythme). Chez Soyinka, les deux se combinent : une même réplique de Baroka peut être à la fois poétique par ses images et comique par son exagération calculée.
\end{attention}

\courssub{Les figures d'amplification : hyperbole et gradation}

\textbf{1. L'hyperbole}

Intuition d'abord : quand un élève affamé déclare « je meurs de faim », personne ne croit qu'il va mourir. Tout le monde comprend qu'il exagère pour marquer l'intensité de sa faim. Cette exagération volontaire et visible est une figure de style très ancienne, que les rhéteurs appellent \emph{hyperbole}.

\begin{definition}[L'hyperbole]
L'\cle{hyperbole} est une figure de style qui consiste à \textbf{exagérer la réalité} au-delà du vraisemblable, afin de frapper l'esprit du lecteur ou de l'auditeur et de produire un effet de surprise, d'admiration, d'émotion ou de rire.
\end{definition}

\begin{exemplebox}[Exemples d'hyperboles]
\begin{itemize}
\item « Sidi possède une beauté à faire pâlir le soleil » : aucune beauté ne peut littéralement pâlir le soleil ; l'exagération magnifie la jeune fille.
\item « Un océan de misère » : la misère n'est pas un océan ; l'image hyperbolique en souligne l'immensité.
\item « Je te l'ai dit mille fois » : exagération familière qui exprime l'insistance.
\end{itemize}
\end{exemplebox}

\textbf{2. La gradation}

Intuition d'abord : un conteur qui décrit un orage ne dit pas « il y eut du vent, de la pluie, du tonnerre » dans le désordre. Il dit : « le vent se leva, la pluie tomba, le tonnerre éclata, et le ciel entier sembla s'effondrer ». Chaque terme est plus fort que le précédent : l'auditeur sent la tension monter. Cette montée en puissance organisée est la gradation.

\begin{definition}[La gradation]
La \cle{gradation} est une figure de style qui consiste à énumérer des mots ou des propositions \textbf{ordonnés selon une intensité croissante} (gradation ascendante) ou \textbf{décroissante} (gradation descendante). Elle crée un effet de montée en puissance ou d'affaiblissement progressif.
\end{definition}

\begin{exemplebox}[Gradation ascendante et descendante]
\begin{itemize}
\item \textbf{Ascendante} : « Il cria, hurla, rugit » : chaque verbe est plus intense que le précédent ; la colère monte jusqu'au paroxysme.
\item \textbf{Descendante} : « Il avait perdu sa richesse, sa famille, son honneur, sa vie » : l'énumération s'enfonce vers le pire, ou bien, selon l'ordre choisi, s'apaise progressivement.
\end{itemize}
\end{exemplebox}

\begin{popfigure}
\begin{center}
\begin{tikzpicture}[font=\small]
% Gradation ascendante
\node[font=\bfseries, popblue] at (3,3.4) {Gradation ascendante (faible $\rightarrow$ fort)};
\draw[->, thick, popblue] (0.3,0.6) -- (5.7,2.9);
\node[draw, fill=popblueL, rounded corners, minimum width=1.5cm] at (0.8,0.3) {cria};
\node[draw, fill=popblue!50, rounded corners, minimum width=1.5cm] at (2.7,1.4) {hurla};
\node[draw, fill=popblue, text=white, rounded corners, minimum width=1.5cm] at (4.6,2.5) {rugit};
% Gradation descendante
\node[font=\bfseries, poporange] at (9.4,3.4) {Gradation descendante (fort $\rightarrow$ faible)};
\draw[->, thick, poporange] (6.8,2.9) -- (12.2,0.6);
\node[draw, fill=poporange, text=white, rounded corners, minimum width=1.7cm] at (7.3,2.5) {tonnerre};
\node[draw, fill=poporange!60, rounded corners, minimum width=1.7cm] at (9.2,1.4) {pluie};
\node[draw, fill=poporangeL, rounded corners, minimum width=1.7cm] at (11.1,0.3) {bruine};
\end{tikzpicture}
\legende{La gradation ordonne les termes selon une échelle d'intensité : montée en puissance (à gauche) ou affaiblissement progressif (à droite).}
\end{center}
\end{popfigure}

\begin{attention}[Erreur fréquente : gradation ou simple énumération ?]
Une énumération n'est pas forcément une gradation. « Il acheta du pain, du sel et de l'huile » est une simple liste : aucun terme n'est plus intense que l'autre. La gradation suppose une \textbf{hiérarchie d'intensité} perceptible : « il grogna, cria, tempêta ». Au Baccalauréat technique, nommer « gradation » une simple liste est une erreur d'analyse qui coûte des points.
\end{attention}

\textbf{3. Tableau comparatif}

\begin{popfigure}
\begin{center}
\begin{tikzpicture}[font=\small]
\fill[popblueL] (0,4) rectangle (12.6,4.9);
\node[font=\bfseries] at (1.6,4.45) {Critère};
\node[font=\bfseries] at (5.6,4.45) {Hyperbole};
\node[font=\bfseries] at (10.1,4.45) {Gradation};
\draw[thick] (0,0) rectangle (12.6,4.9);
\draw[thick] (3.2,0) -- (3.2,4.9);
\draw[thick] (8,0) -- (8,4.9);
\foreach \y in {1,2,3,4} {\draw (0,\y) -- (12.6,\y);};
\node[font=\bfseries, align=left] at (1.6,3.5) {Définition};
\node[align=left, text width=4.4cm] at (5.6,3.5) {Exagération au-delà\\ du vraisemblable};
\node[align=left, text width=4.2cm] at (10.1,3.5) {Suite ordonnée selon\\ une intensité croissante\\ ou décroissante};
\node[font=\bfseries, align=left] at (1.6,2.5) {Marqueurs};
\node[align=left, text width=4.4cm] at (5.6,2.5) {Terme unique excessif,\\ image démesurée};
\node[align=left, text width=4.2cm] at (10.1,2.5) {Au moins trois termes\\ en ordre d'intensité};
\node[font=\bfseries, align=left] at (1.6,1.5) {Effet de sens};
\node[align=left, text width=4.4cm] at (5.6,1.5) {Surprendre, admirer,\\ dramatiser, ridiculiser};
\node[align=left, text width=4.2cm] at (10.1,1.5) {Monter en puissance,\\ créer une tension ou\\ un relâchement};
\node[font=\bfseries, align=left] at (1.6,0.5) {Exemple};
\node[align=left, text width=4.4cm] at (5.6,0.5) {« une beauté à faire\\ pâlir le soleil »};
\node[align=left, text width=4.2cm] at (10.1,0.5) {« Il cria, hurla, rugit »};
\end{tikzpicture}
\legende{Tableau comparatif des deux figures d'amplification : hyperbole et gradation.}
\end{center}
\end{popfigure}

\courssub{Effets de sens : à quoi servent ces figures dans la pièce ?}

Une figure de style n'est jamais gratuite. Dans \emph{Le Lion et la perle}, l'hyperbole et la gradation produisent trois grands effets :

\begin{itemize}
\item \textbf{Valoriser} : quand Baroka ou le poète vante Sidi par des hyperboles (« la perle du village », une beauté qui éclipse toutes les autres), la jeune fille est magnifiée, élevée au rang de trésor.
\item \textbf{Dramatiser} : la gradation ascendante dans les répliques de colère ou de défi fait monter la tension dramatique ; le spectateur sent que l'affrontement approche de son point de rupture.
\item \textbf{Ridiculiser} : quand Lakunle enchaîne les grands mots dans une gradation pompeuse, l'excès même de son éloquence le rend comique : Soyinka utilise l'amplification pour dénoncer la vanité d'une modernité mal digérée.
\end{itemize}

\begin{methode}[Analyser une figure d'amplification en quatre étapes]
\begin{enumerate}
\item \textbf{Repérer} : identifier dans la réplique le terme excessif (hyperbole) ou la suite ordonnée (gradation).
\item \textbf{Nommer} : donner le nom exact de la figure, en précisant, pour la gradation, si elle est ascendante ou descendante.
\item \textbf{Citer} : recopier entre guillemets le passage exact concerné, sans le paraphraser.
\item \textbf{Interpréter} : expliquer l'\emph{effet de sens} produit dans ce contexte précis (valorisation, dramatisation, ridicule), en lien avec le personnage qui parle et son but.
\end{enumerate}
\end{methode}

\begin{exoresolu}[Analyse d'une réplique]
\textbf{Énoncé.} Dans une réplique, un admirateur de Sidi s'écrie : « Sa beauté illumine le village, éblouit le marché, fait pâlir le soleil lui-même ! » Repérez et analysez les figures d'amplification contenues dans cette phrase.
\tcblower
\textbf{Solution détaillée.}
\begin{enumerate}
\item \emph{Repérage} : on relève trois propositions successives (« illumine le village », « éblouit le marché », « fait pâlir le soleil ») et un terme excessif final.
\item \emph{Nomination} : la suite des trois propositions forme une \textbf{gradation ascendante} (l'espace touché par la beauté s'élargit : village, marché, soleil) ; la dernière proposition est en outre une \textbf{hyperbole} (aucune beauté ne peut faire pâlir le soleil).
\item \emph{Citation} : « illumine le village, éblouit le marché, fait pâlir le soleil lui-même ».
\item \emph{Interprétation} : la gradation fait monter l'admiration jusqu'au paroxysme, et l'hyperbole finale couronne le mouvement ; Sidi est ainsi magnifiée, élevée au rang de merveille. Cet éloge excessif confirme son surnom de « Perle » et justifie la rivalité entre Baroka et Lakunle.
\end{enumerate}
\end{exoresolu}

\courssub{Du repérage des figures au résumé : neutraliser l'amplification}

Venons-en au lien décisif avec l'objet de notre leçon : la contraction de texte. Le résumé consiste à \emph{reformuler les idées essentielles} d'un texte avec ses propres mots, en un nombre de lignes imposé. Or les figures d'amplification ne sont pas des idées : ce sont des \emph{habillages expressifs} des idées.

La règle est donc la suivante : dans un résumé, on \textbf{neutralise} les hyperboles et les gradations, c'est-à-dire qu'on rend l'idée sans l'image, sans l'exagération, sans la montée en puissance.

\begin{center}
\begin{tabularx}{\linewidth}{|X|X|}
\hline
\rowcolor{popblueL} \textbf{Texte original (expressif)} & \textbf{Reformulation pour le résumé} \\
\hline
« Sa beauté illumine le village, éblouit le marché, fait pâlir le soleil lui-même ! » & Sidi est d'une beauté exceptionnelle, admirée de tous. \\
\hline
« Il cria, hurla, rugit sa colère contre le sort. » & Il exprima violemment sa colère. \\
\hline
« Lakunle déversa un torrent, un fleuve, un océan de grands mots. » & Lakunle parla abondamment et avec emphase. \\
\hline
\end{tabularx}
\end{center}

\begin{attention}[Deux pièges du résumé]
\begin{itemize}
\item \textbf{Recopier les images} : garder « faire pâlir le soleil » dans un résumé, c'est échouer à reformuler ; l'examinateur sanctionne la paraphrase déguisée.
\item \textbf{Perdre l'idée en neutralisant} : attention à ne pas supprimer l'information portée par la figure. L'hyperbole « un océan de misère » signifie que la misère est \emph{extrême} : le résumé doit conserver cette intensité par un terme sobre (« une misère extrême »), et non la réduire à « une misère ».
\end{itemize}
\end{attention}

\begin{exercice}[À vous de jouer]
Dans les répliques de Baroka et de Lakunle étudiées en classe, relevez deux hyperboles et une gradation. Pour chacune, appliquez la méthode en quatre étapes (repérer, nommer, citer, interpréter), puis proposez une reformulation neutre qui pourrait figurer dans un résumé de la scène.
\end{exercice}

\begin{corrige}[Exercice 1]
Éléments de correction (les citations exactes dépendent de l'édition utilisée en classe).
\begin{enumerate}
\item \emph{Hyperbole chez Baroka} : lorsque le baalè se vante d'une vigueur ou d'une sagesse sans limites, on repère l'exagération, on la nomme hyperbole, on cite la réplique, et on montre qu'elle sert à impressionner Sidi et à ridiculiser la jeunesse de Lakunle. Reformulation neutre : « Baroka affirme rester vigoureux et expérimenté malgré son âge. »
\item \emph{Hyperbole chez Lakunle} : ses déclarations grandiloquentes sur le progrès (« le village entrera d'un bond dans le siècle des machines et des lumières ») exagèrent la transformation à venir. Effet : comique de prétention. Reformulation : « Lakunle promet au village une modernisation rapide. »
\item \emph{Gradation} : toute suite de trois termes d'intensité croissante dans les répliques de colère ou de séduction est recevable, à condition de justifier la hiérarchie d'intensité et l'effet de montée en puissance. Reformulation : on rend l'idée par un seul terme fort, sans la liste.
\end{enumerate}
\end{corrige}

\begin{aretenir}[Bilan de la séance]
\begin{itemize}
\item \emph{Le Lion et la perle} (1959) de Wole Soyinka met en scène l'affrontement entre Baroka (tradition) et Lakunle (modernité) autour de Sidi.
\item La lecture méthodique commence par le repérage de la situation d'énonciation, des personnages et de l'enjeu de l'extrait.
\item L'\cle{hyperbole} exagère la réalité pour frapper l'esprit ; la \cle{gradation} ordonne des termes selon une intensité croissante ou décroissante.
\item Ces figures d'amplification valorisent, dramatisent ou ridiculisent ; on les analyse en quatre étapes : repérer, nommer, citer, interpréter.
\item Dans un résumé, on neutralise ces procédés : on rend l'idée, pas l'image, tout en conservant le degré d'intensité de la pensée de l'auteur.
\end{itemize}
\end{aretenir}

\coursec{La structure argumentative du texte}

Dans la section précédente, nous avons poursuivi la lecture méthodique du \emph{Lion et la perle} de Wole Soyinka et découvert les figures d'amplification (hyperbole, gradation) qui donnent de la force aux paroles des personnages. Ces figures servent un objectif précis : convaincre et persuader. Baroka, le Lion, ne parle jamais pour ne rien dire : chacune de ses tirades défend une idée, s'appuie sur des raisons et s'orne d'exemples. C'est exactement ce fonctionnement que nous allons maintenant étudier de manière générale, car c'est lui qui rend possible la contraction de texte. Pour résumer un texte argumentatif sans le trahir, il faut d'abord savoir comment il est construit : repérer sa \cle{thèse}, ses \cle{arguments}, ses \cle{exemples} et les \cle{connecteurs logiques} qui relient le tout. C'est l'objet de cette section.

\courssub{Qu'est-ce qu'un texte argumentatif ?}

\begin{definition}[Texte argumentatif]
Un \cle{texte argumentatif} est un texte dans lequel un auteur défend un point de vue (une thèse) en avançant des raisons (des arguments) illustrées par des faits concrets (des exemples). Il vise deux objectifs : \textbf{convaincre}, c'est-à-dire amener le lecteur à adhérer à une idée par la raison, et/ou \textbf{persuader}, c'est-à-dire l'amener à adhérer en touchant ses émotions.
\end{definition}

L'éditorial de presse, le discours politique, la plaidoirie d'avocat, la lettre de réclamation sont des textes argumentatifs. Au Probatoire et au Baccalauréat technique, les textes soumis à la contraction sont presque toujours argumentatifs : articles de presse, extraits d'essais, discours. Savoir en dégager la structure est donc la compétence de base du résumé.

\begin{attention}[Convaincre et persuader ne se confondent pas]
Convaincre s'adresse à l'intelligence (preuves, raisonnement, logique) ; persuader s'adresse aux sentiments (émotions, images, valeurs). Un même texte mêle souvent les deux. Dans un résumé, on conserve les raisons (le convaincre) et on élimine la plupart des effets émotionnels (le persuader), car ce sont eux qui prennent le plus de place.
\end{attention}

\courssub{La thèse : l'idée principale défendue}

\begin{definition}[Thèse]
La \cle{thèse} est le point de vue que l'auteur défend et que tout le texte sert à établir. Elle répond à la question : « Que veut me faire admettre l'auteur ? »
\end{definition}

Intuitivement, la thèse est ce qui resterait si l'on ne pouvait garder qu'\emph{une seule phrase} du texte. Tout le reste — arguments, exemples, anecdotes — n'existe que pour la rendre acceptable au lecteur.

\textbf{Où repérer la thèse ?}
\begin{itemize}
\item Elle se trouve le plus souvent au \textbf{début} du texte (l'auteur annonce sa position, puis la justifie) ou à la \textbf{fin} (l'auteur accumule les raisons, puis conclut par sa position).
\item Elle est fréquemment signalée par des formulations comme : \emph{il faut admettre que}, \emph{nous pensons que}, \emph{force est de constater que}, \emph{en définitive}, \emph{c'est pourquoi}.
\item Elle est \textbf{générale} : elle ne concerne pas un cas particulier, mais une situation d'ensemble.
\end{itemize}

\begin{exemplebox}[Thèse dans un éditorial]
Dans la phrase : « La jeunesse camerounaise n'est pas une génération perdue : elle est une génération en attente d'opportunités », la thèse est claire — l'auteur refuse l'idée reçue d'une jeunesse « perdue » et défend l'idée que le problème vient du manque d'opportunités, non des jeunes eux-mêmes.
\end{exemplebox}

\courssub{Les arguments : les raisons qui justifient la thèse}

\begin{definition}[Argument]
Un \cle{argument} est une raison avancée pour justifier la thèse. Il répond à la question « \textbf{pourquoi ?} » : pourquoi faut-il admettre la thèse ?
\end{definition}

Un texte argumentatif comporte en général plusieurs arguments, et ils ne se valent pas tous. Il faut donc les \textbf{hiérarchiser} :
\begin{itemize}
\item l'\textbf{argument principal} est celui sans lequel la démonstration s'effondre ; il porte directement la thèse ;
\item les \textbf{arguments secondaires} renforcent, précisent ou complètent l'argument principal ; ils peuvent parfois être fusionnés ou fortement réduits dans un résumé.
\end{itemize}

Pour distinguer un argument d'une simple affirmation, on applique le test du « parce que » : si l'on peut dire « la thèse est vraie \emph{parce que}... », ce qui suit est un argument.

\begin{exemplebox}[Thèse et arguments]
Thèse : « Il faut généraliser l'enseignement technique au Cameroun. »
\begin{itemize}
\item Argument 1 (principal) : il répond aux besoins du marché de l'emploi, qui manque de techniciens qualifiés.
\item Argument 2 (secondaire) : il valorise les talents pratiques des élèves peu à l'aise dans les filières générales.
\item Argument 3 (secondaire) : il favorise l'entrepreneuriat et la création d'emplois.
\end{itemize}
\end{exemplebox}

\courssub{Les exemples : les illustrations concrètes}

\begin{definition}[Exemple]
L'\cle{exemple} est un cas particulier, un fait, une anecdote, une statistique ou une citation qui illustre un argument et le rend concret. Il répond à la question « par exemple ? » ou « quelle preuve visible ? ».
\end{definition}

L'exemple donne de la chair au raisonnement, mais il est \textbf{particulier} là où l'argument est \textbf{général}. C'est pourquoi, dans un résumé, les exemples sont \emph{les premiers candidats à la suppression} : on garde la règle générale, on abandonne le cas particulier.

\begin{exemplebox}[Argument et exemple]
Argument : « L'enseignement technique répond aux besoins du marché de l'emploi. »
Exemple : « Ainsi, une grande entreprise pétrolière de Douala a recruté en 2024 plus de deux cents techniciens issus des lycées techniques de la ville. »
Dans le résumé, on conserve l'argument et on supprime l'exemple (ou on le réduit à « comme le montrent les recrutements récents » si l'auteur insiste dessus).
\end{exemplebox}

\begin{attention}[Erreur fréquente]
Le piège classique consiste à prendre un \textbf{exemple frappant} (une anecdote émouvante, un chiffre spectaculaire) pour l'\textbf{idée principale} du texte. L'exemple est mémorable parce qu'il est concret ; mais l'idée essentielle est toujours la plus générale. Avant de retenir une phrase comme « idée principale », demandez-vous : « Cette phrase s'applique-t-elle à un cas particulier ou à l'ensemble du sujet ? »
\end{attention}

\courssub{Les connecteurs logiques : les balises de la structure}

Les \cle{connecteurs logiques} sont les mots et expressions qui articulent les idées entre elles. Ce sont de véritables panneaux indicateurs qui signalent où l'on se trouve dans le raisonnement.

\begin{center}
\begin{tabularx}{\linewidth}{|l|X|X|}
\hline
\rowcolor{popblueL} \textbf{Fonction} & \textbf{Connecteurs} & \textbf{Rôle dans le texte} \\
\hline
Énumération & d'abord, ensuite, enfin, de plus, en outre & annonce les arguments successifs \\
\hline
Explication & en effet, car, parce que, c'est que & introduit une justification \\
\hline
Illustration & par exemple, ainsi, c'est le cas de & introduit un exemple (à supprimer au résumé) \\
\hline
Opposition & cependant, mais, pourtant, néanmoins & signale une objection ou une nuance \\
\hline
Conséquence & donc, par conséquent, c'est pourquoi, ainsi & annonce la conclusion ou la thèse \\
\hline
\end{tabularx}
\end{center}

Relever les connecteurs, c'est donc reconstituer le plan du texte sans effort : chaque « ensuite » annonce un nouvel argument, chaque « par exemple » annonce un élément à éliminer, chaque « donc » final annonce la thèse ou la conclusion.

\courssub{Schéma de la structure argumentative}

Un texte argumentatif simple obéit au schéma suivant : thèse $\rightarrow$ argument 1 + exemple $\rightarrow$ argument 2 + exemple $\rightarrow$ conclusion (qui reprend souvent la thèse).

\begin{popfigure}
\begin{center}
\begin{tikzpicture}[
  grow=down,
  level distance=1.6cm,
  these/.style={rectangle, rounded corners, draw=popdark, fill=popblueL, text width=6.5cm, align=center, font=\bfseries},
  arg/.style={rectangle, rounded corners, draw=popblue, fill=popblue!15, text width=4.2cm, align=center},
  ex/.style={rectangle, rounded corners, draw=popmuted, fill=popcream, text width=3.6cm, align=center, font=\small},
  lien/.style={draw=popdark, thick}
]
\node[these, align=center]{THÈSE\\ (idée principale défendue)}
  child { node[arg] {Argument 1 (principal)}
    child { node[ex] {Exemple 1} }
    child { node[ex] {Exemple 2} }
  }
  child { node[arg] {Argument 2}
    child { node[ex] {Exemple 3} }
  }
  child { node[arg] {Argument 3}
    child { node[ex] {Exemple 4} }
  };
\node[rectangle, rounded corners, draw=poporange, fill=poporangeL, text width=6.5cm, align=center, font=\bfseries] at (0,-6.4) {CONCLUSION (reprise de la thèse)};
\end{tikzpicture}
\end{center}
\legende{Arbre de la structure argumentative : la thèse est au sommet, les arguments forment les branches, les exemples sont les feuilles. Au résumé, on coupe les feuilles (exemples) et on garde le tronc (thèse) et les branches (arguments).}
\end{popfigure}

\begin{popfigure}
\begin{center}
\begin{tikzpicture}
% Entonnoir : 4 niveaux
\fill[popblue!35]   (0,4.5) -- (8,4.5) -- (6.8,3.3) -- (1.2,3.3) -- cycle;
\fill[popblue!25]   (1.2,3.3) -- (6.8,3.3) -- (5.6,2.1) -- (2.4,2.1) -- cycle;
\fill[popgreen!25]  (2.4,2.1) -- (5.6,2.1) -- (4.6,0.9) -- (3.4,0.9) -- cycle;
\fill[poporange!35] (3.4,0.9) -- (4.6,0.9) -- (4.3,-0.6) -- (3.7,-0.6) -- cycle;
\draw[popdark, thick] (0,4.5) -- (8,4.5) -- (4.3,-0.6) -- (3.7,-0.6) -- cycle;
\node[align=left, anchor=west, font=\small] at (8.3,3.9) {\textbf{Texte complet : 100\,\%}\\ thèse + arguments + exemples\\ + répétitions + digressions};
\node[align=left, anchor=west, font=\small] at (7.1,2.7) {Suppression des exemples,\\ anecdotes, citations};
\node[align=left, anchor=west, font=\small] at (5.9,1.5) {Fusion des arguments\\ secondaires, suppression\\ des répétitions};
\node[align=left, anchor=west, font=\small] at (4.9,0.15) {\textbf{Résumé : 25\,\%}\\ thèse + arguments\\ essentiels reformulés};
\draw[-{Stealth}, popdark, thick] (8.2,3.9) -- (7.0,3.9);
\draw[-{Stealth}, popdark, thick] (7.0,2.7) -- (5.8,2.7);
\draw[-{Stealth}, popdark, thick] (5.8,1.5) -- (4.8,1.5);
\draw[-{Stealth}, popdark, thick] (4.8,0.15) -- (4.35,0.15);
\end{tikzpicture}
\end{center}
\legende{Schéma en entonnoir de la contraction : du texte complet (100\,\%) au résumé (environ 25\,\%), par élimination progressive des éléments secondaires. La structure argumentative sert de filtre : ce qui n'est ni thèse ni argument sort de l'entonnoir.}
\end{popfigure}

\courssub{Idées essentielles et idées secondaires}

Toute la contraction de texte repose sur une distinction fondamentale :

\begin{aretenir}[Idée essentielle / idée secondaire]
\begin{itemize}
\item Sont \textbf{essentielles} : la \cle{thèse} et les \cle{arguments}. Elles constituent le squelette du texte ; les supprimer détruit le raisonnement.
\item Sont \textbf{secondaires} : les \cle{exemples}, les anecdotes, les citations, les chiffres détaillés, les répétitions, les paraphrases, les digressions. Elles ornent ou éclairent le raisonnement ; les supprimer l'appauvrit mais ne le détruit pas.
\end{itemize}
Le résumé conserve les idées essentielles (reformulées) et élimine les idées secondaires.
\end{aretenir}

\begin{exemplebox}[Un éditorial en chiffres]
Dans un éditorial de 400 mots, on compte en général \textbf{1 thèse}, \textbf{3 à 4 arguments} et \textbf{5 à 8 exemples}, auxquels s'ajoutent répétitions et formules de transition. Le résumé de 100 mots (le quart) ne garde que la thèse et les arguments, reformulés avec les mots du résumé et reliés par quelques connecteurs.
\end{exemplebox}

\courssub{Méthode : dégager la structure argumentative d'un texte}

\begin{methode}[Les quatre gestes du repérage]
\begin{enumerate}
\item \textbf{Souligner la thèse} : après une première lecture complète, repérer la phrase (ou les deux phrases) qui exprime le point de vue de l'auteur. Vérifier qu'elle est générale et que tout le texte la sert.
\item \textbf{Numéroter les arguments} : dans la marge, inscrire A1, A2, A3... devant chaque raison avancée. Appliquer le test du « parce que » pour vérifier qu'il s'agit bien d'arguments.
\item \textbf{Barrer les exemples} : rayer (au crayon) anecdotes, statistiques détaillées, citations, cas particuliers — tout ce qui suit « par exemple », « ainsi », « c'est le cas de ».
\item \textbf{Relever les connecteurs} : entourer d'abord, ensuite, en effet, cependant, donc... et reconstituer le plan : thèse $\rightarrow$ A1 $\rightarrow$ A2 $\rightarrow$ A3 $\rightarrow$ conclusion.
\end{enumerate}
\end{methode}

\begin{exoresolu}[Structure argumentative d'un éditorial camerounais]
Énoncé. Voici un extrait d'éditorial (fictif) paru dans un quotidien de Douala :

\emph{« Nos dirigeants répètent que la jeunesse est l'avenir du Cameroun. Il est temps de passer des discours aux actes : l'avenir se construit aujourd'hui, en donnant aux jeunes un emploi. D'abord, un jeune qui travaille participe à la production nationale et enrichit le pays au lieu de grossir les files des candidats à l'émigration. On a vu, l'an dernier, plus de trois mille diplômés quitter le pays par le port de Douala. Ensuite, l'emploi des jeunes garantit la paix sociale : un oisif, dit le proverbe, est un danger pour le village. Enfin, investir dans la jeunesse, c'est préparer les cadres de demain. Certains objectent que l'État manque de moyens ; mais chaque franc placé dans la formation professionnelle rapporte demain des impôts et des entreprises. Donc, financer l'emploi des jeunes n'est pas une dépense : c'est le meilleur investissement national. »}

Dégager la structure argumentative de ce texte : thèse, arguments, exemples, connecteurs.

\tcblower
Solution détaillée.
\begin{itemize}
\item \textbf{Thèse} : « Financer l'emploi des jeunes n'est pas une dépense : c'est le meilleur investissement national » (annoncée dès la deuxième phrase : « l'avenir se construit aujourd'hui, en donnant aux jeunes un emploi », et reprise dans la conclusion introduite par « donc »).
\item \textbf{Arguments} : A1 — un jeune qui travaille enrichit le pays au lieu d'émigrer (argument principal) ; A2 — l'emploi des jeunes garantit la paix sociale ; A3 — investir dans la jeunesse prépare les cadres de demain. L'objection « l'État manque de moyens » est réfutée par un quatrième raisonnement : la formation rapporte plus qu'elle ne coûte.
\item \textbf{Exemples} (à supprimer au résumé) : les trois mille diplômés partis par le port de Douala ; le proverbe sur l'oisif.
\item \textbf{Connecteurs relevés} : « il est temps de » (annonce de la thèse), « d'abord », « ensuite », « enfin » (énumération des arguments), « certains objectent que... mais » (objection réfutée), « donc » (conclusion).
\end{itemize}
Plan dégagé : Thèse $\rightarrow$ A1 (emploi = richesse nationale) $\rightarrow$ A2 (emploi = paix sociale) $\rightarrow$ A3 (emploi = cadres de demain) $\rightarrow$ réfutation de l'objection $\rightarrow$ conclusion. Ce plan suffira, dans la section suivante, à rédiger le résumé.
\end{exoresolu}

\begin{exercice}[À vous de jouer]
Relevez dans un article de presse camerounais de votre choix (ou dans l'éditorial ci-dessus) : 1) la phrase qui exprime la thèse ; 2) trois arguments, numérotés et hiérarchisés ; 3) tous les exemples, que vous barrerez ; 4) cinq connecteurs logiques, en précisant leur fonction. Rédigez ensuite en trois lignes le plan du texte.
\end{exercice}

\begin{corrige}[Exercice]
Éléments de correction (appliqués à l'éditorial de l'exercice résolu) : 1) Thèse : financer l'emploi des jeunes est le meilleur investissement national. 2) A1 (principal) : le jeune employé produit des richesses ; A2 : l'emploi assure la paix sociale ; A3 : il forme les cadres de demain. 3) Exemples barrés : les trois mille diplômés émigrés ; le proverbe de l'oisif. 4) Connecteurs : « d'abord », « ensuite », « enfin » (énumération), « mais » (opposition), « donc » (conclusion). Plan : thèse $\rightarrow$ A1 $\rightarrow$ A2 $\rightarrow$ A3 $\rightarrow$ réfutation $\rightarrow$ conclusion.
\end{corrige}

\begin{savaistu}[Le résumé au Baccalauréat technique]
Au Probatoire et au Baccalauréat technique, le résumé (contraction de texte) est noté selon trois critères principaux : la \textbf{fidélité aux idées} de l'auteur (on ne doit ni ajouter son avis, ni déformer la pensée), la \textbf{reformulation} (on n'a pas le droit de recopier les phrases du texte : il faut les dire autrement, avec ses propres mots) et le \textbf{respect de la longueur imposée} (en général le quart du texte, avec une marge de 10\,\%). Dégager correctement la structure argumentative, comme nous venons de l'apprendre, est la condition du premier critère : impossible d'être fidèle à des idées qu'on n'a pas identifiées.
\end{savaistu}

\begin{aretenir}[L'essentiel de la section]
\begin{itemize}
\item Un texte argumentatif défend une \cle{thèse} (idée principale) au moyen d'\cle{arguments} (raisons, répondant à « pourquoi ? ») illustrés par des \cle{exemples} (cas particuliers, répondant à « par exemple ? »).
\item Les \cle{connecteurs logiques} (d'abord, ensuite, en effet, cependant, donc) balisent la structure et permettent de reconstituer le plan.
\item Le résumé conserve la thèse et les arguments (idées essentielles) et élimine les exemples, répétitions et digressions (idées secondaires).
\item Méthode : souligner la thèse, numéroter les arguments, barrer les exemples, relever les connecteurs.
\end{itemize}
\end{aretenir}

\coursec{Lecture méthodique 5 et lexique commun / lexique spécialisé}

Dans la section précédente, nous avons appris à repérer la \cle{structure argumentative} d'un texte : la thèse, les arguments, les exemples. Mais pour résumer un texte, il ne suffit pas de connaître son architecture : il faut aussi comprendre les \cle{mots} qui la construisent. Or tout texte mélange deux grandes familles de vocabulaire : les mots que tout le monde comprend, et les mots réservés à un domaine précis. C'est cette distinction — \cle{lexique commun} et \cle{lexique spécialisé} — que nous allons étudier maintenant, en l'appliquant d'abord à notre œuvre au programme, \emph{Le Lion et la perle} de Wole Soyinka.

\courssub{Lecture méthodique n° 5 du \emph{Lion et la perle}}

\textbf{Rappel du contexte et situation de l'extrait.}\par

Rappelons la situation de la pièce : dans le village d'Ilujinlè, le vieux chef Baroka (le « Lion ») rivalise avec le jeune instituteur Lakunlé pour conquérir Sidi, la « perle » du village, devenue célèbre depuis que ses photographies ont paru dans un magazine. Lors des lectures précédentes, nous avons vu Lakunlé tenter de convaincre Sidi de l'épouser sans verser la dot, au nom de la « modernité », et Sidi refuser avec vivacité.

Dans cette cinquième lecture méthodique, nous poursuivons l'analyse de l'extrait en nous concentrant sur trois points : le \cle{thème}, le \cle{ton} et la \cle{progression dramatique}.

\textbf{1. Le thème.} L'extrait met en scène l'affrontement entre deux conceptions du mariage : la \emph{tradition} (la dot, les coutumes du village, l'autorité du chef) et la \emph{modernité} (l'amour libre, l'école, le rejet des coutumes jugées « arriérées »). Ce thème central — le conflit entre tradition et modernité — traverse toute l'œuvre.

\textbf{2. Le ton.} Soyinka choisit un ton essentiellement \emph{comique} et \emph{satirique} : Lakunlé est ridiculisé par ses grands mots maladroits, Sidi se moque de lui avec esprit, et le dialogue théâtral fait éclater le rire. Mais derrière le rire se cache une réflexion sérieuse : le dramaturge se moque autant du faux modernisme de Lakunlé que de la ruse du vieux Baroka.

\textbf{3. La progression dramatique.} On observe une montée de la tension : l'échange commence par une discussion presque amicale, s'envenime progressivement (reproches, sarcasmes), et aboutit à un éclat — le refus net de Sidi. Cette progression en trois temps (amorce, escalade, chute) est typique de la scène de dispute au théâtre.

\begin{aretenir}[Repères pour la lecture méthodique]
Analyser un extrait de théâtre, c'est toujours répondre à trois questions : \textbf{de quoi parle-t-on ?} (le thème), \textbf{comment en parle-t-on ?} (le ton : comique, pathétique, polémique, satirique), \textbf{comment la scène avance-t-elle ?} (la progression dramatique : exposition, tension, éclat, chute).
\end{aretenir}

\textbf{Observation : deux vocabulaires dans l'extrait.}\par

Lisons attentivement le dialogue. Certains mots sont compris immédiatement par n'importe quel lecteur francophone : \emph{village}, \emph{femme}, \emph{mariage}, \emph{épouser}, \emph{aimer}. D'autres renvoient à une réalité culturelle précise : la \emph{dot}, le \emph{chef coutumier}, la \emph{danse rituelle}, les \emph{masques}, la \emph{baie} (le palais du chef). Cette observation simple est le point de départ de notre leçon de langue : le vocabulaire d'un texte se partage entre un fonds commun et des vocabulaires spécialisés.

\courssub{Le lexique commun}

Commençons par l'intuition. Si vous dites à n'importe quel francophone — un enfant de Yaoundé, un commerçant de Douala, un agriculteur de Bafoussam — les phrases « Je \emph{mange} », « Il \emph{marche} vers le \emph{village} », « L'\emph{eau} est \emph{chaude} », tout le monde comprend sans dictionnaire. Ces mots forment le socle de la langue.

\begin{definition}[Le lexique commun]
Le \cle{lexique commun} (ou vocabulaire courant) est l'ensemble des mots compris et employés par tous les locuteurs d'une langue dans les situations ordinaires de la vie quotidienne : \emph{manger}, \emph{marcher}, \emph{maison}, \emph{village}, \emph{enfant}, \emph{parler}, \emph{beau}, \emph{grand}.
\end{definition}

Le lexique commun présente trois caractéristiques :

\begin{itemize}
\item il est \textbf{universel} dans la communauté linguistique : aucune formation particulière n'est nécessaire pour le comprendre ;
\item il est \textbf{fréquent} : ce sont les mots qui reviennent le plus souvent dans les conversations, les journaux, les discours ;
\item il est \textbf{stable} : ces mots traversent les générations sans grande modification.
\end{itemize}

\begin{exemplebox}[Lexique commun dans \emph{Le Lion et la perle}]
Dans la bouche de Sidi : « Tu veux m'\emph{épouser} sans \emph{payer} ? » — les mots \emph{épouser}, \emph{payer}, \emph{femme}, \emph{honte} appartiennent au lexique commun : tout locuteur les comprend, quelle que soit sa profession.
\end{exemplebox}

\courssub{Le lexique spécialisé}

Imaginons maintenant une autre scène. À l'hôpital, le médecin annonce : « Le patient présente une \emph{anémie} ; nous prescrivons une \emph{transfusion}. » Au tribunal, l'avocat déclare : « Mon client fait \emph{appel} du \emph{jugement} rendu en première \emph{instance}. » Dans une salle informatique : « La \emph{connexion} est coupée, vérifie le \emph{routeur}. » Chaque fois, un profane peut rester perplexe : ces mots appartiennent à un domaine précis.

\begin{definition}[Le lexique spécialisé]
Le \cle{lexique spécialisé} est l'ensemble des termes propres à un domaine d'activité, à une discipline ou à une profession : médecine (\emph{diagnostic}, \emph{ordonnance}), droit (\emph{arrêté}, \emph{décret}, \emph{tribunal}), informatique (\emph{logiciel}, \emph{forfait}, \emph{serveur}), agriculture (\emph{semis}, \emph{récolte}, \emph{cacao}), économie (\emph{subvention}, \emph{déficit}).
\end{definition}

Le lexique spécialisé se reconnaît à trois critères :

\begin{itemize}
\item il est \textbf{lié à un domaine} : c'est le critère essentiel — un mot est spécialisé parce qu'il appartient au vocabulaire d'une activité précise ;
\item il est \textbf{précis} : chaque terme désigne exactement une réalité du domaine, sans ambiguïté ;
\item il est \textbf{compris pleinement par les spécialistes} : le profane peut l'ignorer ou le comprendre vaguement.
\end{itemize}

\begin{attention}[Erreur fréquente]
Un mot long, savant ou rare n'est pas forcément un terme spécialisé ! \emph{Magnifique}, \emph{extraordinaire}, \emph{incommensurable} sont des mots longs mais ils appartiennent au lexique commun : tout le monde les emploie. Inversement, un mot court peut être spécialisé : \emph{semis} (agriculture), \emph{greffe} (médecine), \emph{bail} (droit). \textbf{Le seul critère est l'appartenance à un domaine d'activité}, pas la longueur ni la rareté du mot.
\end{attention}

\begin{exemplebox}[Lexique spécialisé de l'économie]
Dans un article sur le budget du Cameroun, les mots \emph{subvention}, \emph{recette budgétaire}, \emph{déficit}, \emph{dette publique}, \emph{croissance} appartiennent au lexique spécialisé de l'économie et des finances. Ils désignent des notions précises que l'économiste maîtrise exactement.
\end{exemplebox}

\courssub{Lexique commun et lexique spécialisé dans \emph{Le Lion et la perle}}

L'œuvre de Soyinka offre un excellent terrain d'observation, car elle mêle constamment les deux registres de vocabulaire.

\textbf{Le lexique spécialisé de la tradition yoruba.}\par

Le village d'Ilujinlè vit selon les coutumes yoruba (peuple du Nigeria, voisin culturel des peuples de l'ouest camerounais). Soyinka emploie donc un vocabulaire spécialisé du domaine \emph{culturel et traditionnel} :

\begin{itemize}
\item la \emph{baie} : le palais, la résidence du chef ;
\item le \emph{chef coutumier} (Baroka) : autorité traditionnelle du village ;
\item la \emph{dot} : biens ou argent versés par le prétendant à la famille de la fiancée ;
\item les \emph{danses rituelles} et le \emph{masque} : cérémonies traditionnelles ;
\item le \emph{griot} et les \emph{louanges} : art oratoire traditionnel.
\end{itemize}

\textbf{Le lexique commun.}\par

En contrepoint, les émotions et les actions des personnages sont dites dans le vocabulaire de tous les jours : \emph{aimer}, \emph{rire}, \emph{beauté}, \emph{honte}, \emph{village}, \emph{marché}, \emph{jeunesse}, \emph{vieillard}.

Ce mélange n'est pas un hasard : il permet à Soyinka d'ancrer sa pièce dans une culture précise (lexique spécialisé) tout en racontant une histoire humaine universelle (lexique commun). Pour le lecteur camerounais, ce lexique traditionnel est d'ailleurs facilement accessible, car les réalités désignées — dot, chef coutumier, danses rituelles — existent aussi dans nos propres traditions.

\courssub{Lexique spécialisé et situations camerounaises}

Le français pratiqué au Cameroun mobilise chaque jour des lexiques spécialisés très vivants. En voici les principaux domaines.

\begin{center}
\begin{tabularx}{\linewidth}{|l|X|X|}
\hline
\rowcolor{popblueL} \textbf{Domaine} & \textbf{Lexique spécialisé (exemples)} & \textbf{Lexique commun associé} \\
\hline
Administration & arrêté, décret, préfet, sous-préfet, arrondissement, matricule & papier, bureau, signer, chef \\
\hline
Agriculture & cacao, palmier à huile, semis, plantation, récolte, cabosse & champ, planter, travail, pluie \\
\hline
Économie / finances & subvention, budget, recette, déficit, investissement & argent, payer, vendre, prix \\
\hline
NTIC (informatique) & forfait, connexion, réseau, application, téléchargement & téléphone, appeler, message \\
\hline
Santé & ordonnance, diagnostic, vaccination, paludisme & malade, hôpital, soigner \\
\hline
Justice & tribunal, jugement, avocat, plainte, témoin & faute, punir, vérité \\
\hline
\end{tabularx}
\end{center}

\begin{savaistu}[Le français du Cameroun]
Le français parlé au Cameroun a intégré des termes spécialisés locaux qui sont devenus peu à peu du \emph{lexique commun régional} : la \emph{tontine} ou le \emph{njangi} (association d'épargne rotative), le \emph{mbanga} ou le \emph{ndolè} (plats), le \emph{makossa} et le \emph{bikutsi} (musiques). Nés dans des domaines précis (finance populaire, cuisine, musique), ces mots sont aujourd'hui compris par tous les Camerounais. Cela montre que la frontière entre lexique spécialisé et lexique commun n'est pas figée : un terme spécialisé peut, avec le temps, entrer dans l'usage courant.
\end{savaistu}

\begin{popfigure}
\begin{center}
\begin{tikzpicture}[scale=1]
\draw[thick, popblue, fill=popblue!15] (-1.6,0) circle (2.4);
\draw[thick, poporange, fill=poporange!15] (1.6,0) circle (2.4);
\node[popblue, font=\bfseries] at (-2.6,2.0) {Lexique commun};
\node[poporange, font=\bfseries] at (2.6,2.0) {Lexique spécialisé};
\node[align=center, font=\small, popdark] at (-2.6,0.3) {manger\\ village\\ maison\\ beau};
\node[align=center, font=\small, popdark] at (2.6,0.3) {décret\\ diagnostic\\ semis\\ forfait};
\node[align=center, font=\small\itshape, popdark] at (0,-0.1) {tontine\\ connexion\\ dot};
\node[align=center, font=\footnotesize, popmuted] at (0,-1.35) {termes passés\\ dans l'usage courant};
\end{tikzpicture}
\end{center}
\legende{Diagramme de Venn : le lexique commun (à gauche) et le lexique spécialisé (à droite) se recoupent partiellement. La zone centrale contient les termes nés dans un domaine spécialisé mais passés dans l'usage courant : \emph{tontine} (finance populaire), \emph{connexion} (informatique), \emph{dot} (droit coutumier).}
\end{popfigure}

\courssub{Lexique spécialisé et résumé : que faire d'un terme technique ?}

Voici le lien décisif avec notre objet d'étude : la \cle{contraction de texte}. Quand vous résumez un texte, vous rencontrerez presque toujours des termes spécialisés. Faut-il les garder, les supprimer, les expliquer ?

\begin{methode}[Traiter un terme spécialisé dans un résumé]
\begin{enumerate}
\item \textbf{Repérer} le terme spécialisé et identifier son domaine (économie, droit, médecine, culture...).
\item \textbf{Se demander s'il est porteur de l'idée} : si l'on supprime le terme, est-ce que l'idée essentielle de la phrase disparaît ?
\item \textbf{Si le terme est indispensable} (il porte l'idée) : le \emph{conserver tel quel} dans le résumé. Exemple : dans un texte sur le cacao, le mot \emph{cacao} ne peut pas être remplacé.
\item \textbf{Si le terme est secondaire} (simple précision, détail technique) : le \emph{supprimer} ou le remplacer par une formulation générale équivalente.
\item \textbf{Si le terme est nécessaire mais difficile} : le conserver et l'\emph{expliquer brièvement} par une périphrase claire, sans perdre le sens.
\end{enumerate}
\end{methode}

\begin{exemplebox}[Reformulation par périphrase]
Phrase originale : « L'État a augmenté les \emph{subventions} allouées aux planteurs de \emph{cacao} pour réduire le \emph{déficit} de la filière. »

Résumé possible : « L'État a accru ses aides financières aux producteurs de cacao pour combler les pertes de ce secteur. »

On a conservé \emph{cacao} (indispensable, c'est le sujet) et remplacé \emph{subventions} par « aides financières » et \emph{déficit} par « pertes » : le sens est intact, la formulation est plus générale.
\end{exemplebox}

\begin{attention}[Piège de la reformulation]
Remplacer un terme spécialisé par une périphrase ne doit jamais \textbf{changer le sens} ni \textbf{allonger exagérément} le résumé. « Subvention » devient « aide financière » (correct), mais pas « cadeau » (sens faussé : une subvention se justifie, un cadeau est gratuit). Et une périphrase de dix mots pour un seul terme alourdit le résumé, ce qui contredit le principe même de la contraction.
\end{attention}

\courssub{Entraînement : classer le vocabulaire d'un article économique}

\begin{exoresolu}[Classer 15 mots en lexique commun ou spécialisé]
Voici 15 mots extraits d'un article économique camerounais sur la filière cacao : \emph{cacao, manger, subvention, village, exportation, pluie, décret, cabosse, marcher, budget, plantation, enfant, investissement, beau, recette}.

Classez-les en deux colonnes : lexique commun / lexique spécialisé (en précisant le domaine).

\tcblower
\textbf{Solution détaillée.}

Critère appliqué : le mot est-il employé par tous les locuteurs dans la vie quotidienne (lexique commun), ou propre à un domaine d'activité (lexique spécialisé) ?

\begin{center}
\begin{tabularx}{\linewidth}{|X|X|}
\hline
\rowcolor{popblueL} \textbf{Lexique commun} & \textbf{Lexique spécialisé (domaine)} \\
\hline
manger & cacao (agriculture / économie) \\
\hline
village & subvention (économie / finances) \\
\hline
pluie & exportation (économie / commerce) \\
\hline
marcher & décret (administration / droit) \\
\hline
enfant & cabosse (agriculture : fruit du cacaoyer) \\
\hline
beau & budget (économie / finances) \\
\hline
 & plantation (agriculture) \\
\hline
 & investissement (économie) \\
\hline
 & recette (finances : somme encaissée) \\
\hline
\end{tabularx}
\end{center}

Bilan : 6 mots de lexique commun, 9 mots de lexique spécialisé répartis en trois domaines (agriculture, économie-finances, administration).

Remarque : \emph{recette} est un cas intéressant. Dans « une recette de cuisine », c'est du lexique commun ; dans « les recettes de l'État », c'est un terme spécialisé des finances. \textbf{C'est le contexte qui décide.}
\end{exoresolu}

\begin{exercice}[À vous de jouer]
\begin{enumerate}
\item Dans \emph{Le Lion et la perle}, relevez trois mots du lexique commun et trois termes du lexique spécialisé de la tradition (dot, baie, danse rituelle...). Justifiez chaque classement.
\item Reformulez la phrase suivante pour un résumé, en remplaçant les termes spécialisés soulignés par des périphrases claires : « Le \emph{décret} fixe le montant des \emph{subventions} versées aux planteurs afin de relancer les \emph{exportations} de cacao. »
\item Trouvez deux termes spécialisés du domaine des NTIC qui sont en train de passer dans le lexique commun au Cameroun. Expliquez pourquoi.
\end{enumerate}
\end{exercice}

\begin{corrige}[Exercice : éléments de correction]
\begin{enumerate}
\item Lexique commun : par exemple \emph{beauté}, \emph{aimer}, \emph{village} — compris de tous les locuteurs. Lexique spécialisé de la tradition : \emph{dot} (droit coutumier), \emph{baie} (palais du chef), \emph{danse rituelle} (cérémonie) — chacun renvoie au domaine culturel yoruba.
\item Reformulation possible : « Une décision officielle de l'État fixe le montant des aides financières versées aux planteurs afin de relancer les ventes de cacao à l'étranger. » Le sens est conservé, les trois termes techniques sont expliqués simplement.
\item Exemples : \emph{forfait}, \emph{connexion}, \emph{application}. Nés dans le vocabulaire des télécommunications et de l'informatique, ils sont aujourd'hui employés et compris par presque tous les Camerounais, même sans formation technique : ils migrent du lexique spécialisé vers le lexique commun.
\end{enumerate}
\end{corrige}

\begin{aretenir}[L'essentiel de la section]
\begin{itemize}
\item Le \cle{lexique commun} rassemble les mots compris et employés par tous les locuteurs ; le \cle{lexique spécialisé} rassemble les termes propres à un domaine d'activité.
\item Le critère de classement est l'\textbf{appartenance à un domaine}, non la longueur ou la rareté du mot ; le contexte peut faire basculer un mot d'un lexique à l'autre.
\item Dans un résumé, un terme spécialisé se \textbf{conserve} s'il porte l'idée essentielle ; sinon on le \textbf{supprime} ou on le \textbf{reformule} par une périphrase claire et équivalente.
\item Cette compétence prépare directement l'étape suivante : les \emph{techniques de réduction et de rédaction} du résumé, que nous aborderons dans la prochaine section.
\end{itemize}
\end{aretenir}

\coursec{Techniques de réduction et de rédaction du résumé}

Dans la section précédente, nous avons poursuivi la lecture méthodique du \emph{Lion et la perle} et appris à distinguer le \cle{lexique commun} du \cle{lexique spécialisé}. Cette compétence nous sera précieuse : pour résumer un texte, il faut souvent remplacer un vocabulaire précis ou technique par des termes plus généraux. Nous abordons maintenant le cœur pratique de la contraction de texte : les \cle{techniques de réduction} et la \cle{rédaction} du résumé. C'est l'aboutissement de tout le travail mené depuis la première séance de cette unité.

\courssub{Lecture méthodique n° 6 du \emph{Lion et la perle} : synthèse et bilan}

Avant d'entrer dans la technique, faisons le point sur notre œuvre. Les lectures méthodiques 4, 5 et 6 nous ont permis de parcourir les moments décisifs de la pièce de Wole Soyinka.

\begin{aretenir}[Bilan des lectures 4 à 6]
\begin{itemize}
\item \textbf{Lecture 4} : le conflit s'installe entre Baroka, le vieux chef rusé (le « Lion »), et Lakunle, l'instituteur moderniste, autour de Sidi, la belle du village (la « perle »). Nous y avons étudié les figures d'amplification : \cle{hyperbole} et \cle{gradation}.
\item \textbf{Lecture 5} : la ruse de Baroka se déploie ; il fait courir le bruit de sa mort et de son impuissance pour piéger Sidi. Nous y avons distingué lexique commun et lexique spécialisé (vocabulaire de la tradition yoruba, de la coutume, de la dot).
\item \textbf{Lecture 6} : Sidi, venue se moquer du chef, tombe dans son piège et devient sa femme. Lakunle, qui refusait de payer la dot au nom du « progrès », est battu par la ruse traditionnelle.
\end{itemize}
\end{aretenir}

\textbf{Thèmes majeurs de l'œuvre} (à retenir pour les exposés) :
\begin{itemize}
\item le conflit entre \cle{tradition} et \cle{modernité} ;
\item la \cle{ruse} et l'intelligence face à la naïveté ;
\item la condition de la \cle{femme} et la question de la dot ;
\item le pouvoir, la séduction et la vanité.
\end{itemize}

Ce bilan est lui-même un petit exercice de contraction : nous venons de ramener trois lectures à quelques lignes, en ne gardant que l'essentiel. C'est exactement le geste que nous allons maintenant apprendre à faire méthodiquement.

\courssub{La contraction de texte : définition et règle du quart}

\begin{definition}[Contraction de texte]
La \cle{contraction de texte} (ou résumé) est une opération qui consiste à \textbf{réduire un texte à une fraction de sa longueur} en \textbf{conservant ses idées essentielles}, sans rien ajouter (ni commentaire, ni avis personnel) et en reformulant avec ses propres mots.
\end{definition}

\begin{propriete}[La règle du quart]
Au Probatoire et au Baccalauréat technique, la consigne impose en général de résumer le texte \textbf{au quart de sa longueur} : un texte de 400 mots donne un résumé d'environ 100 mots. Une marge de tolérance d'environ $\pm 10\,\%$ est admise.
\end{propriete}

\begin{exemplebox}[Un exemple chiffré]
Pour un texte de 480 mots, le résumé attendu compte environ :
\[ 480 \times \frac{1}{4} = 120 \text{ mots} \]
Avec la marge de tolérance de $\pm 10\,\%$ ($120 \times 0,1 = 12$), le correcteur accepte un résumé compris entre \textbf{108 et 132 mots}. En dessous de 108 mots, des idées essentielles ont probablement été perdues ; au-dessus de 132, le texte n'est pas assez contracté.
\end{exemplebox}

\begin{savaistu}[Pourquoi le quart ?]
La règle du quart est la norme des épreuves camerounaises de Français. Elle n'est pas arbitraire : réduire un texte au quart oblige à distinguer l'essentiel de l'accessoire, ce qui est précisément la compétence évaluée. Attention : \textbf{dépasser la longueur imposée est sanctionné au barème}, même si le contenu est bon. Compter ses mots fait partie de l'exercice.
\end{savaistu}

\begin{popfigure}
\begin{center}
\begin{tikzpicture}
\begin{axis}[
    ybar,
    width=0.85\linewidth,
    height=6cm,
    bar width=1.2cm,
    ymin=0, ymax=450,
    ylabel={Nombre de mots},
    symbolic x coords={Texte original, Brouillon, Résumé final},
    xtick=data,
    nodes near coords,
    every node near coord/.append style={font=\small},
    ymajorgrids=true,
    grid style={popline!40},
]
\addplot[fill=popblue] coordinates {(Texte original,400)};
\addplot[fill=poporange] coordinates {(Brouillon,200)};
\addplot[fill=popgreen] coordinates {(Résumé final,100)};
\end{axis}
\end{tikzpicture}
\end{center}
\legende{La règle du quart en pratique : le texte original (400 mots) est d'abord réduit de moitié au brouillon (200 mots), puis ramené au quart dans le résumé final (100 mots).}
\end{popfigure}

\courssub{Les quatre techniques de réduction}

Pour passer de 400 à 100 mots sans trahir l'auteur, on dispose de quatre techniques complémentaires. Elles s'appliquent presque toujours dans cet ordre : on commence par enlever, puis on remplace.

\begin{popfigure}
\begin{center}
\begin{tikzpicture}[
    etape/.style={rectangle, rounded corners, draw=popblue, fill=popblueL, text width=2.6cm, minimum height=1.6cm, align=center, font=\small},
    ex/.style={rectangle, draw=poporange, fill=poporangeL, text width=2.6cm, align=center, font=\scriptsize},
    fleche/.style={-{Stealth[length=3mm]}, thick, popdark}
]
\node[etape, align=center] (e1){\textbf{1. Suppression}\\ enlever l'accessoire};
\node[etape, right=0.9cm of e1, align=center] (e2){\textbf{2. Généralisation}\\ le terme générique};
\node[etape, right=0.9cm of e2, align=center] (e3){\textbf{3. Condensation}\\ phrase $\to$ groupe nominal};
\node[etape, right=0.9cm of e3, align=center] (e4){\textbf{4. Substitution}\\ périphrase $\to$ mot simple};
\draw[fleche] (e1) -- (e2);
\draw[fleche] (e2) -- (e3);
\draw[fleche] (e3) -- (e4);
\node[ex, below=0.3cm of e1] {on supprime : exemples, répétitions, citations};
\node[ex, below=0.3cm of e2] {mangues, ananas, bananes $\to$ \emph{fruits}};
\node[ex, below=0.3cm of e3] {« il arriva en retard car le car était en panne » $\to$ « une panne le retarda »};
\node[ex, below=0.3cm of e4] {« le roi des animaux » $\to$ « le lion »};
\end{tikzpicture}
\end{center}
\legende{Les quatre techniques de réduction, dans l'ordre d'application, avec un exemple transformé à chaque étape.}
\end{popfigure}

\courssub{Technique 1 : la suppression}

\begin{definition}[Suppression]
La \cle{suppression} consiste à éliminer purement et simplement les éléments qui ne font pas partie des idées essentielles du texte.
\end{definition}

On supprime en priorité :
\begin{itemize}
\item les \textbf{exemples} (ils illustrent une idée, ils ne la remplacent pas) ;
\item les \textbf{répétitions} et les paraphrases (la même idée dite deux fois) ;
\item les \textbf{citations} ;
\item les \textbf{détails descriptifs} (couleurs, dimensions, circonstances anecdotiques) ;
\item les \textbf{expansions du nom} : adjectifs, compléments du nom, propositions relatives non indispensables.
\end{itemize}

\begin{exemplebox}[Suppression en action]
Texte : « Baroka, le vieux chef du village d'Ilujinlè, âgé de soixante-deux ans, mais encore vigoureux comme un jeune homme, décida, après mûre réflexion, de conquérir Sidi. »

Réduit : « Baroka décida de conquérir Sidi. »

On a supprimé l'expansion du nom (« le vieux chef du village d'Ilujinlè »), le détail de l'âge, la comparaison flatteuse et la circonstance « après mûre réflexion ». L'idée essentielle survit intacte.
\end{exemplebox}

\courssub{Technique 2 : la généralisation}

\begin{definition}[Généralisation]
La \cle{généralisation} consiste à remplacer une énumération de termes particuliers par le \textbf{terme générique} qui les recouvre tous.
\end{definition}

C'est ici que notre travail sur le lexique devient utile : mangues, ananas, bananes et papayes se résument en \emph{fruits} ; houes, machettes et pelles se résument en \emph{outils agricoles} ; instituteurs, médecins et ingénieurs se résument en \emph{cadres} ou \emph{professionnels qualifiés}.

\begin{exemplebox}[Généralisation en action]
Texte : « Dans les villages, le progrès apporte l'école, le dispensaire, la route goudronnée, l'électricité et l'eau potable. »

Réduit : « Le progrès apporte des infrastructures et des services essentiels aux villages. »
\end{exemplebox}

\begin{attention}[Ne pas trop généraliser]
Le terme générique doit rester \textbf{fidèle} à la pensée de l'auteur. Si le texte énumère des \emph{fruits tropicaux}, écrire « produits alimentaires » est trop vague : on perd une information essentielle. La généralisation monte d'un cran, pas de dix.
\end{attention}

\courssub{Technique 3 : la condensation}

\begin{definition}[Condensation]
La \cle{condensation} consiste à réduire une phrase ou une proposition entière à un mot ou à un groupe nominal. C'est la technique la plus rentable : une proposition subordonnée de huit mots peut devenir un seul nom.
\end{definition}

Le mécanisme repose souvent sur la \textbf{nominalisation} : le verbe devient un nom.
\begin{center}
\begin{tabularx}{\linewidth}{|X|X|}
\hline
\rowcolor{popblueL} \textbf{Forme longue (verbe + subordonnée)} & \textbf{Forme condensée (nom)} \\
\hline
parce que le car était en panne & une panne \\
\hline
quand le pays est devenu indépendant & l'indépendance \\
\hline
parce qu'il refusait de payer la dot & son refus de la dot \\
\hline
afin que les enfants soient éduqués & l'éducation des enfants \\
\hline
\end{tabularx}
\end{center}

\begin{exemplebox}[Condensation en action]
Texte : « Il arriva en retard parce que le car était en panne. » (10 mots)

Réduit : « Une panne le retarda. » (4 mots)

La subordonnée de cause « parce que le car était en panne » s'est condensée en un seul nom : « une panne ». L'information essentielle (retard + cause) est conservée.
\end{exemplebox}

\courssub{Technique 4 : la substitution}

\begin{definition}[Substitution]
La \cle{substitution} consiste à remplacer une périphrase, une figure de style ou une expression imagée par un mot simple et direct.
\end{definition}

Cette technique fait le lien avec notre étude des figures d'amplification : une \cle{hyperbole} ou une image poétique se traduit, dans un résumé, par son sens propre.

\begin{center}
\begin{tabularx}{\linewidth}{|X|X|}
\hline
\rowcolor{popblueL} \textbf{Expression du texte} & \textbf{Substitution dans le résumé} \\
\hline
le roi des animaux & le lion \\
\hline
la perle d'Ilujinlè & Sidi (la plus belle fille du village) \\
\hline
il versa des torrents de larmes & il pleura abondamment \\
\hline
le pays s'est levé d'un bond & le pays s'est développé rapidement \\
\hline
\end{tabularx}
\end{center}

\begin{savaistu}[Néologismes et résumé]
L'étude des \cle{néologismes} (mots nouveaux entrés dans la langue) est utile ici : un texte contemporain peut contenir des néologismes (ex. \emph{télétravailler}, \emph{écoresponsable}). Dans un résumé, on les conserve s'ils sont le terme précis, ou on les substitue par une périphrase claire si le terme est trop récent ou technique (ex. \emph{écoresponsable} $\to$ \emph{respectueux de l'environnement}).
\end{savaistu}

\courssub{La rédaction du résumé}

Les quatre techniques produisent un \textbf{brouillon} : une matière réduite, mais encore décousue. Il faut maintenant \textbf{rédiger}, c'est-à-dire transformer ce brouillon en un texte continu et lisible.

\begin{propriete}[Règles de rédaction du résumé]
\begin{enumerate}
\item Écrire des \textbf{phrases complètes} (sujet, verbe, complément), pas une liste de notes.
\item Assurer l'\textbf{enchaînement logique} par des connecteurs : \emph{d'abord, ensuite, en effet, cependant, c'est pourquoi, finalement}.
\item Employer la \textbf{troisième personne} et le présent de vérité générale ou les temps du texte ; ne jamais dire « je ».
\item Rester \textbf{neutre} : aucune appréciation personnelle, aucun jugement sur l'auteur ou ses idées.
\item Respecter la \textbf{longueur imposée} (le quart, $\pm 10\,\%$).
\end{enumerate}
\end{propriete}

\begin{attention}[Erreur fréquente n° 1 : recopier]
Recopier des phrases entières du texte est la faute la plus sanctionnée. Le résumé exige la \textbf{reformulation avec ses propres mots} : si le correcteur retrouve des segments identiques au texte, il considère que l'élève n'a pas compris et pénalise la copie. On ne garde du texte original que les noms propres et les termes techniques irremplaçables.
\end{attention}

\begin{attention}[Erreur fréquente n° 2 : donner son avis]
Écrire « je pense que l'auteur a raison » ou « ce texte est très intéressant » est hors sujet. Le résumé doit rester \textbf{neutre et objectif} : il rend compte de la pensée de l'auteur, pas de celle de l'élève. L'avis personnel a sa place dans la dissertation, jamais dans la contraction de texte.
\end{attention}

\courssub{La vérification finale}

Un résumé n'est jamais rendu sans contrôle. La vérification comporte trois gestes :

\begin{enumerate}
\item \textbf{Compter les mots} : le total doit se situer dans la fourchette imposée (par exemple 108 à 132 mots pour un texte de 480 mots). Si le résumé est trop long, on repasse par la suppression et la condensation ; s'il est trop court, on vérifie qu'aucune idée essentielle n'a été abandonnée.
\item \textbf{Relire à voix haute} : les phrases doivent être grammaticales, l'enchaînement fluide, l'orthographe et les accords soignés.
\item \textbf{Contrôler la fidélité} : comparer une dernière fois le résumé au texte. Chaque idée du résumé doit venir de l'auteur ; aucune idée majeure de l'auteur ne doit manquer.
\end{enumerate}

\begin{methode}[Réussir un résumé en 5 étapes]
\begin{enumerate}
\item \textbf{Lire deux fois} le texte : une première lecture globale, une seconde lecture analytique en soulignant les idées essentielles.
\item \textbf{Dégager la structure} : repérer la thèse, les arguments et les étapes du raisonnement (structure argumentative étudiée à la section 2).
\item \textbf{Réduire au brouillon} : appliquer les quatre techniques (suppression, généralisation, condensation, substitution) pour viser d'abord la moitié, puis le quart.
\item \textbf{Rédiger} : phrases complètes, connecteurs logiques, troisième personne, neutralité.
\item \textbf{Vérifier} : compter les mots, relire, contrôler la fidélité à la pensée de l'auteur.
\end{enumerate}
\end{methode}

\courssub{TP guidé : réduire un paragraphe d'environ 100 mots}

\begin{experience}[Réduire de moitié puis au quart]
\textbf{Texte de départ (environ 105 mots)} — extrait d'un discours sur l'éducation au Cameroun :

« L'éducation est, nous le savons tous, la clé du développement de notre cher pays, le Cameroun. Partout, de Yaoundé à Maroua, de Douala à Bamenda, nos enfants ont besoin d'écoles, de collèges, de lycées, d'universités, mais aussi de bibliothèques, de laboratoires et de terrains de sport. Prenons l'exemple de ma propre région : là-bas, les élèves parcourent chaque matin dix kilomètres à pied sous le soleil brûlant pour gagner leur salle de classe. C'est pourquoi je le déclare solennellement devant vous : investir dans l'éducation, c'est investir dans l'avenir de la jeunesse et dans la prospérité de toute la nation camerounaise. »

\textbf{Étape 1 — Suppression} : on élimine l'exemple personnel (« ma propre région... dix kilomètres à pied »), les formules oratoires (« nous le savons tous », « je le déclare solennellement devant vous », « notre cher pays »), l'énumération des villes.

\textbf{Étape 2 — Généralisation} : « écoles, collèges, lycées, universités » $\to$ « établissements scolaires » ; « bibliothèques, laboratoires, terrains de sport » $\to$ « équipements ».

\textbf{Étape 3 — Condensation} : « investir dans l'éducation, c'est investir dans l'avenir » $\to$ « l'éducation conditionne l'avenir ».

\textbf{Brouillon (environ 50 mots, la moitié)} : « L'éducation est la clé du développement du Cameroun. Partout, les enfants ont besoin d'établissements scolaires et d'équipements. Certains élèves parcourent de longues distances pour étudier. Investir dans l'éducation, c'est donc investir dans l'avenir de la jeunesse et dans la prospérité de la nation. »

\textbf{Étape 4 — Résumé final (environ 26 mots, le quart)} : « L'éducation est la clé du développement du Cameroun : les enfants ont besoin d'établissements scolaires équipés. Investir dans l'éducation, c'est investir dans l'avenir de la nation. »
\end{experience}

\begin{exoresolu}[Application : contrôler la fidélité]
Énoncé. Voici le résumé proposé par un élève pour le texte du TP : « L'éducation est la clé du développement du Cameroun, et je trouve que l'orateur a tout à fait raison. Les enfants ont besoin d'écoles, de collèges, de lycées, d'universités, de bibliothèques, de laboratoires et de terrains de sport. Investir dans l'éducation, c'est investir dans l'avenir de la nation. » (52 mots). Relevez trois fautes et corrigez-les.

\tcblower
Solution détaillée.
\begin{enumerate}
\item \textbf{Avis personnel} : « je trouve que l'orateur a tout à fait raison » viole la règle de neutralité. $\to$ Supprimer.
\item \textbf{Recopie} : l'énumération « écoles, collèges, lycées, universités, bibliothèques, laboratoires et terrains de sport » est reprise mot pour mot. $\to$ Généraliser : « des établissements scolaires équipés ».
\item \textbf{Longueur} : 52 mots au lieu d'environ 26. $\to$ Après correction, le résumé devient : « L'éducation est la clé du développement du Cameroun : les enfants ont besoin d'établissements scolaires équipés. Investir dans l'éducation, c'est investir dans l'avenir de la nation. » (26 mots, fidèle et neutre).
\end{enumerate}
\end{exoresolu}

\begin{exercice}[À vous de jouer]
Résumez en 25 à 30 mots le passage suivant (environ 100 mots), en appliquant les quatre techniques et les règles de rédaction : « Lakunle, le jeune instituteur du village, qui se rêvait en homme moderne et qui méprisait ouvertement les coutumes de ses ancêtres, refusait catégoriquement de payer la dot pour épouser Sidi. Selon lui, cette pratique ancienne, héritée des générations passées, était une véritable humiliation pour la femme, une marchandisation pure et simple de l'être humain. Mais Sidi, fière de sa beauté célèbre dans tout le village, exigeait cette dot comme la preuve de sa valeur. »
\end{exercice}

\begin{corrige}[Exercice : proposition de résumé]
Étapes : suppression des expansions (« jeune », « ouvertement », « catégoriquement », « pure et simple »), condensation (« refusait de payer la dot » $\to$ « son refus de la dot »), généralisation (« pratique ancienne, héritée des générations passées » $\to$ « coutume »).

Résumé possible (28 mots) : « Lakunle, partisan de la modernité, refuse de payer la dot, qu'il considère comme une humiliation de la femme ; Sidi, au contraire, l'exige comme preuve de sa valeur. »

Vérification : 28 mots (le quart de 100, tolérance respectée), aucune phrase recopiée, aucun avis personnel, les deux idées essentielles (le refus de Lakunle, l'exigence de Sidi) sont conservées.
\end{corrige}

\begin{aretenir}[L'essentiel de la section]
\begin{itemize}
\item Le résumé représente environ \textbf{le quart} du texte original ($\pm 10\,\%$).
\item Quatre techniques de réduction : \cle{suppression}, \cle{généralisation}, \cle{condensation}, \cle{substitution}.
\item Rédaction : phrases complètes, connecteurs logiques, troisième personne, neutralité absolue.
\item Vérification : compter les mots, relire, contrôler la fidélité à la pensée de l'auteur.
\end{itemize}
\end{aretenir}

\coursec{Les néologismes}

Dans la section précédente, nous avons appris à réduire et à rédiger un résumé : sélectionner les idées essentielles, supprimer les exemples, reformuler avec ses propres mots. Mais une difficulté subsiste : que faire des mots nouveaux qui apparaissent dans les textes contemporains, notamment les articles de presse sur l'économie numérique ? Faut-il les recopier tels quels, les remplacer, les expliquer ? Pour répondre à cette question, il faut d'abord comprendre ce qu'est un \cle{néologisme} : c'est l'objet de cette section. Cette compétence est directement utile pour l'épreuve de contraction de texte au Probatoire et au Baccalauréat technique, où les textes proposés portent souvent sur des réalités récentes (numérique, santé, environnement).

\courssub{Définition du néologisme}

Commençons par une observation simple. Ouvrez un dictionnaire de 1970 et cherchez les mots \emph{télétravail}, \emph{smartphone}, \emph{influenceur} : ils n'y figurent pas. Pourtant, aujourd'hui, tout le monde les comprend et les emploie. Ces mots sont nés récemment, parce que les réalités qu'ils désignent sont récentes. La langue, comme un organisme vivant, se renouvelle sans cesse pour nommer le monde nouveau.

\begin{definition}[Néologisme]
Un \cle{néologisme} est une unité lexicale nouvelle apparue récemment dans une langue. Cette nouveauté peut porter sur :
\begin{itemize}
\item la \textbf{forme} : un mot qui n'existait pas du tout est créé (\emph{déconfinement}) ;
\item le \textbf{sens} : un mot ancien reçoit une signification nouvelle (\emph{souris} pour désigner l'accessoire d'ordinateur).
\end{itemize}
Tant que le mot n'est pas entré dans l'usage courant ni dans les dictionnaires, on parle de \emph{néologisme d'occasion} ; une fois consacré par l'usage, il devient un mot de la langue à part entière.
\end{definition}

\begin{attention}[Néologisme ou terme spécialisé ?]
Erreur fréquente : confondre \cle{néologisme} et \cle{terme spécialisé}. Un néologisme est défini par sa \textbf{nouveauté} ; un terme spécialisé est défini par son appartenance à un \textbf{domaine} (médecine, informatique, droit), et il peut être très ancien. Ainsi, \emph{fièvre} est un terme médical ancien, ce n'est pas un néologisme. En revanche, \emph{covidéo} est un néologisme, même s'il appartient aussi au vocabulaire de la santé. Les deux notions peuvent se croiser, mais ne se confondent pas.
\end{attention}

\courssub{Les procédés de formation des néologismes}

Comment naît un mot nouveau ? Les linguistes distinguent trois grandes sources de néologismes.

\textbf{1. Les néologismes de forme.} On fabrique un mot inédit à partir d'éléments déjà existants dans la langue, par deux procédés principaux :
\begin{itemize}
\item la \textbf{dérivation} : on ajoute un préfixe ou un suffixe à un mot existant. Exemples : \emph{déconfinement} (préfixe \emph{dé-} + \emph{confinement}), \emph{influenceur} (\emph{influence} + suffixe \emph{-eur}), \emph{googliser} (\emph{Google} + suffixe \emph{-iser}), \emph{climatosceptique} ;
\item la \textbf{composition} : on assemble deux mots ou deux radicaux. Exemples : \emph{télétravail} (\emph{télé-} + \emph{travail}), \emph{covidéo} (\emph{covid} + \emph{vidéo}), \emph{webinaire} (\emph{web} + \emph{séminaire}).
\end{itemize}

\textbf{2. Les néologismes de sens.} Aucun mot nouveau n'est créé : c'est un mot ancien qui reçoit une signification nouvelle, souvent par métaphore ou métonymie. Exemples :
\begin{itemize}
\item \emph{souris} : l'animal, puis l'accessoire informatique (ressemblance de forme et de « queue ») ;
\item \emph{toile} : le tissu, puis le réseau Internet (traduction de \emph{web}, la toile d'araignée) ;
\item \emph{portable} : ce qui peut être porté, puis le téléphone mobile ou l'ordinateur portable ;
\item \emph{navigation} : l'art de conduire un navire, puis le fait de passer de page en page sur Internet.
\end{itemize}

\textbf{3. Les néologismes d'emprunt.} On importe directement un mot d'une langue étrangère, le plus souvent de l'anglais dans le domaine des nouvelles technologies : \emph{smartphone}, \emph{start-up}, \emph{marketing}, \emph{e-mail}, \emph{streaming}, \emph{business plan}. Certains de ces emprunts sont ensuite adaptés ou remplacés par des équivalents français recommandés : \emph{courriel} pour \emph{e-mail}, \emph{jeune pousse} pour \emph{start-up}.

\begin{exemplebox}[Trois néologismes, trois procédés]
\begin{itemize}
\item \emph{googliser} : néologisme de \textbf{forme} (dérivation par le suffixe \emph{-iser} sur le nom propre \emph{Google}) ;
\item \emph{télétravail} : néologisme de \textbf{forme} (composition), créé dans les années 1970 et relancé massivement depuis 2020 ;
\item \emph{smartphone} : néologisme d'\textbf{emprunt} à l'anglais (\emph{smart} : intelligent ; \emph{phone} : téléphone).
\end{itemize}
\end{exemplebox}

\begin{popfigure}
\begin{center}
\begin{tikzpicture}[
  grow=down,
  level distance=1.6cm,
  level 1/.style={sibling distance=4.6cm},
  level 2/.style={sibling distance=2.2cm},
  every node/.style={draw, rounded corners=2pt, align=center, font=\small, inner sep=3pt},
  edge from parent/.style={draw=popmuted, -{Stealth[length=2mm]}}
]
\node[fill=popblue!25] {\textbf{Néologismes}}
  child { node[fill=popgreen!25, align=center]{de \textbf{forme}\\ (dérivation, composition)}
    child { node[fill=popgreen!10] {\emph{déconfinement}} }
    child { node[fill=popgreen!10] {\emph{influenceur}} }
  }
  child { node[fill=poporange!25, align=center]{de \textbf{sens}\\ (mot ancien, sens nouveau)}
    child { node[fill=poporange!10] {\emph{souris}} }
    child { node[fill=poporange!10] {\emph{toile}} }
  }
  child { node[fill=poppurple!25, align=center]{d'\textbf{emprunt}\\ (autres langues)}
    child { node[fill=poppurple!10] {\emph{smartphone}} }
    child { node[fill=poppurple!10] {\emph{start-up}} }
  };
\end{tikzpicture}
\end{center}
\legende{Arbre de classification des néologismes : trois grandes catégories, avec deux exemples par branche.}
\end{popfigure}

\courssub{Les néologismes dans le français du Cameroun}

Le français parlé et écrit au Cameroun est une langue vivante qui crée ses propres mots pour désigner des réalités locales que le français de France ne nomme pas. Ces créations sont de véritables néologismes, formés selon les mêmes procédés :

\begin{itemize}
\item \textbf{Emprunts aux langues nationales} intégrés au français : \emph{mbanga} (repas de feuilles de manioc), \emph{ndolé} (plat à base de feuilles et d'arachides), \emph{njangui} (tontine), \emph{sawa} (désignant les populations du littoral) ;
\item \textbf{Dérivation et composition sur des réalités locales} : \emph{benskinneur} (conducteur de moto-taxi, de \emph{bendskin}), \emph{call-boxeur} (vendeur de crédit téléphonique dans une cabine), \emph{débrouillardise} (art de se débrouiller dans la vie quotidienne) ;
\item \textbf{Sens nouveau donné à des mots français} : \emph{chargeur} (personne qui aide à charger les bagages dans les gares routières), \emph{Bayam-Sellam} (commerçante qui achète et revend les produits agricoles).
\end{itemize}

Dans le domaine de l'économie numérique, le Cameroun a vu naître des expressions récentes : \emph{mobile money} (argent mobile), \emph{Orange Money}, \emph{MTN MoMo}, \emph{e-commerce}, \emph{Nkap} (plateforme de tontine numérique, du mot ewondo désignant l'argent). Ces mots circulent dans la presse camerounaise et appartiennent au lexique que le candidat au Probatoire doit savoir manipuler.

\begin{savaistu}[Les dictionnaires suivent l'usage]
Chaque année, le \emph{Petit Robert} et le \emph{Petit Larousse} intègrent des centaines de néologismes consacrés par l'usage : ce sont les lexicographes qui observent les journaux, les réseaux sociaux et les conversations pour repérer les mots qui s'installent. Un mot n'entre au dictionnaire que lorsqu'il est employé durablement par un grand nombre de locuteurs. Ainsi, \emph{télétravail}, créé dans les années 1970, n'est devenu courant qu'à partir de 2020 : un néologisme peut mettre des décennies à s'imposer.
\end{savaistu}

\begin{popfigure}
\begin{center}
\begin{tikzpicture}[x=1.55cm]
% axe
\draw[-{Stealth[length=3mm]}, thick, popdark] (0,0) -- (7.4,0);
\foreach \x/\a in {0/1980, 1/1990, 2/1995, 3/2000, 4/2005, 5/2010, 6/2015, 7/2020} {
  \draw[popdark] (\x,0.08) -- (\x,-0.08);
  \node[below, font=\footnotesize] at (\x,-0.1) {\a};
}
% points et étiquettes alternées
\node[above, font=\footnotesize, text=popblue, align=center] at (0.4,0.35) {\emph{baladeur}\\ (1979)};
\fill[popblue] (0.4,0) circle (2pt);
\node[below, font=\footnotesize, text=popgreen, align=center, yshift=-14pt] at (1.6,0) {\emph{courriel}\\ (années 1990)};
\fill[popgreen] (1.6,0) circle (2pt);
\node[above, font=\footnotesize, text=poporange, align=center] at (2.8,0.35) {\emph{toile} (web)\\ (années 1995)};
\fill[poporange] (2.8,0) circle (2pt);
\node[below, font=\footnotesize, text=poppurple, align=center, yshift=-14pt] at (4.2,0) {\emph{smartphone}\\ (années 2000)};
\fill[poppurple] (4.2,0) circle (2pt);
\node[above, font=\footnotesize, text=poppink, align=center] at (5.4,0.35) {\emph{influenceur}\\ (années 2010)};
\fill[poppink] (5.4,0) circle (2pt);
\node[below, font=\footnotesize, text=popgold!80!black, align=center, yshift=-14pt] at (6.8,0) {\emph{covidéo}, \emph{déconfinement}\\ (2020)};
\fill[popgold!80!black] (6.8,0) circle (2pt);
\end{tikzpicture}
\end{center}
\legende{Frise chronologique : apparition de quelques néologismes en français, de 1980 à nos jours (dates approximatives d'entrée dans l'usage).}
\end{popfigure}

\courssub{Méthode : reconnaître un néologisme}

\begin{methode}[Comment reconnaître un néologisme dans un texte]
Face à un mot qui semble nouveau, applique trois vérifications :
\begin{enumerate}
\item \textbf{Le test du dictionnaire ancien} : le mot figure-t-il dans un dictionnaire d'il y a vingt ou trente ans ? S'il est absent, c'est probablement un néologisme.
\item \textbf{Le test de la réalité} : le mot désigne-t-il un objet, une technique, un phénomène social récent (numérique, santé, modes de vie nouveaux) ? Une réalité récente appelle souvent un mot récent.
\item \textbf{Le test de la formation} : le mot est-il construit sur un procédé productif visible — préfixe (\emph{dé-}, \emph{télé-}, \emph{e-}), suffixe (\emph{-iser}, \emph{-eur}), composition, ou mot anglais importé tel quel ? Si oui, il s'agit très vraisemblablement d'un néologisme de forme ou d'emprunt.
\end{enumerate}
Si le mot existe depuis longtemps mais avec un autre sens (\emph{souris}, \emph{toile}), c'est un néologisme de sens.
\end{methode}

\courssub{Les néologismes dans le résumé : que faire ?}

C'est ici que cette leçon rejoint directement la contraction de texte. Lorsque vous résumez un article contenant des néologismes, deux cas se présentent :

\begin{aretenir}[Règle pratique pour le résumé]
\begin{itemize}
\item Si le néologisme est \textbf{entré dans l'usage courant} (\emph{télétravail}, \emph{smartphone}, \emph{mobile money}), \textbf{conservez-le} tel quel dans le résumé : le remplacer par une périphrase alourdirait le texte, ce qui est contraire à l'esprit de la contraction.
\item Si le néologisme est \textbf{très récent ou peu connu} (\emph{covidéo}, \emph{Nkap}, \emph{webinaire}), \textbf{explicitez-le} en une courte reformulation : « plateforme de tontine en ligne », « réunion par écran interposé ». Le résumé doit rester compréhensible par un lecteur non spécialiste.
\item Ne créez jamais vous-même un néologisme dans un résumé : la contraction exige une langue claire, correcte et établie.
\end{itemize}
\end{aretenir}

\begin{exoresolu}[Repérage et reformulation de néologismes]
\textbf{Énoncé.} Relevez dans l'extrait suivant (article de presse camerounais sur l'économie numérique) quatre néologismes, classez-les (forme, sens, emprunt) et proposez pour chacun une reformulation claire utilisable dans un résumé.

\emph{« À Douala, les start-up du e-commerce se multiplient. Grâce au mobile money, les clients paient leurs achats sans se déplacer. Certaines plateformes, comme les tontines numériques de type Nkap, digitalisent des pratiques traditionnelles, tandis que les influenceurs des réseaux sociaux font la promotion de ces nouveaux services. »}

\tcblower
\textbf{Solution détaillée.}
\begin{enumerate}
\item \emph{start-up} : néologisme d'\textbf{emprunt} (anglais). Reformulation : « jeunes entreprises innovantes ».
\item \emph{e-commerce} : néologisme de \textbf{forme} (composition : préfixe \emph{e-}, pour \emph{électronique}, + \emph{commerce}). Reformulation : « commerce en ligne ».
\item \emph{mobile money} : néologisme d'\textbf{emprunt} (anglais). Reformulation : « argent mobile, paiement par téléphone ».
\item \emph{digitalisent} : néologisme de \textbf{forme} (dérivation : \emph{digital} + suffixe verbal \emph{-iser}). Reformulation : « transposent sur des outils numériques ».
\end{enumerate}
Dans un résumé, on conservera \emph{mobile money} et \emph{e-commerce}, entrés dans l'usage, mais on pourra expliciter \emph{Nkap} par « plateforme de tontine en ligne ».
\end{exoresolu}

\begin{exercice}[Entraînement : huit néologismes à repérer]
Lisez le texte suivant, relevez \textbf{huit néologismes}, indiquez leur catégorie (forme, sens, emprunt) et proposez pour chacun une reformulation claire.

\emph{« Depuis le déconfinement, le télétravail s'est installé dans les habitudes. Les Camerounais googlisent leurs questions, commandent sur les sites de e-commerce et règlent leurs factures par mobile money. Les webinaires remplacent les réunions, les influenceurs recommandent des applications, et la toile est devenue le premier marché du pays. Même les tontines passent par des plateformes comme Nkap. »}
\end{exercice}

\begin{corrige}[Exercice : les huit néologismes]
\begin{enumerate}
\item \emph{déconfinement} : forme (dérivation, préfixe \emph{dé-}) ; « levée des restrictions de circulation ».
\item \emph{télétravail} : forme (composition) ; « travail à distance, depuis chez soi ».
\item \emph{googlisent} : forme (dérivation, suffixe \emph{-iser}) ; « recherchent sur Internet ».
\item \emph{e-commerce} : forme (composition) ; « commerce en ligne ».
\item \emph{mobile money} : emprunt (anglais) ; « paiement par téléphone mobile ».
\item \emph{webinaires} : forme (composition : \emph{web} + \emph{séminaire}) ; « réunions ou formations en ligne ».
\item \emph{influenceurs} : forme (dérivation, suffixe \emph{-eur}) ; « personnes suivies qui orientent les choix du public ».
\item \emph{toile} : sens (mot ancien, sens nouveau) ; « l'ensemble du réseau Internet ».
\end{enumerate}
(\emph{Nkap}, nom propre récent, est accepté comme néologisme d'emprunt à une langue nationale : « tontine numérique ».)
\end{corrige}

\begin{aretenir}[L'essentiel de la section]
\begin{itemize}
\item Un \cle{néologisme} est un mot nouveau ou un sens nouveau donné à un mot ancien.
\item Trois catégories : néologismes de \textbf{forme} (dérivation, composition), de \textbf{sens} (métaphore, extension) et d'\textbf{emprunt} (anglais, langues nationales).
\item Le français du Cameroun crée ses propres néologismes, liés aux réalités locales et à l'économie numérique.
\item Dans un résumé : conserver les néologismes entrés dans l'usage, expliciter les plus récents, ne jamais en inventer.
\end{itemize}
\end{aretenir}

\coursec{Reformuler les idées essentielles et exposés sur l'œuvre}

Dans la section précédente, nous avons étudié les \cle{néologismes}, ces mots nouveaux qui enrichissent la langue et renouvellent l'expression. Ce travail sur le vocabulaire nous conduit naturellement à une question centrale du résumé : comment dire la même chose avec d'autres mots ? C'est exactement l'objet de la \cle{reformulation}, compétence décisive pour la contraction de texte au Probatoire et au Baccalauréat technique. Nous verrons ensuite comment cette compétence écrite se transpose à l'oral, à travers deux exposés sur \emph{Le Lion et la perle} de Wole Soyinka.

\courssub{Qu'est-ce que reformuler ?}

\begin{definition}[La reformulation]
La \cle{reformulation} est la reprise fidèle d'une idée dans d'autres termes que ceux de l'auteur. Elle obéit à trois exigences : dire la \textbf{même idée}, la dire \textbf{autrement}, et ne rien ajouter de personnel. Reformuler, ce n'est ni déformer, ni commenter, ni recopier.
\end{definition}

L'intuition est simple : imaginez qu'un camarade absent vous demande de lui rapporter ce que le professeur a expliqué. Vous ne récitez pas ses phrases mot pour mot ; vous redites le contenu avec vos propres mots, sans y mêler votre avis. Le résumé exige exactement cette opération, mais à l'écrit et de manière contrôlée.

\begin{attention}[La paraphrase servile]
L'erreur la plus fréquente est la \cle{paraphrase mot à mot} : l'élève se contente de remplacer deux ou trois mots par des synonymes, en gardant la même construction de phrase. Exemple fautif : « La jeunesse camerounaise croule sous le chômage » devient « La jeunesse du Cameroun s'effondre sous le chômage ». La phrase n'a pas été restructurée : le correcteur y voit une citation déguisée et sanctionne. Une vraie reformulation change à la fois les mots \textbf{et} l'architecture de la phrase.
\end{attention}

\courssub{Les procédés de reformulation}

Quatre procédés permettent de transformer une phrase sans en trahir le sens.

\begin{enumerate}
\item \textbf{La synonymie} : remplacer un mot par un équivalent (« crouler » $\rightarrow$ « être écrasé », « durement » $\rightarrow$ « sévèrement »). Attention : le synonyme doit convenir au contexte.
\item \textbf{Le changement de catégorie grammaticale} : transformer un nom en verbe ou inversement (« la réduction du personnel » $\rightarrow$ « on réduit le personnel » ; « manifester » $\rightarrow$ « organiser une manifestation »).
\item \textbf{La transformation active/passive} : « Le chômage frappe les jeunes » $\rightarrow$ « Les jeunes sont frappés par le chômage ».
\item \textbf{Le regroupement d'idées} (ou condensation) : fondre en une seule phrase plusieurs phrases voisines qui expriment une même idée générale.
\end{enumerate}

\begin{exemplebox}[Une reformulation réussie]
Phrase originale : « La jeunesse camerounaise croule sous le chômage. »

Reformulation : « Selon l'auteur, le chômage frappe durement les jeunes Camerounais. »

On observe : un changement de construction (le sujet « jeunesse » devient complément, « chômage » devient sujet), une synonymie (« croule sous » $\rightarrow$ « frappe durement »), et une formule de neutralité (« selon l'auteur »). L'idée est identique, la forme est nouvelle.
\end{exemplebox}

\begin{popfigure}
\begin{center}
\begin{tabularx}{\linewidth}{|l|X|}
\hline
\rowcolor{popblueL} \textbf{Procédé} & \textbf{Phrase obtenue} \\
\hline
Phrase originale & « La jeunesse camerounaise croule sous le chômage. » \\
\hline
Synonymie & « Les jeunes Camerounais sont écrasés par le chômage. » \\
\hline
Changement grammatical & « Le chômage des jeunes Camerounais atteint un niveau accablant. » \\
\hline
Actif/passif & « Les jeunes Camerounais sont durement frappés par le chômage. » \\
\hline
Condensation & « Le chômage pèse lourdement sur toute la jeunesse du pays. » \\
\hline
\end{tabularx}
\end{center}
\legende{Tableau « Dire autrement » : quatre procédés appliqués à une même phrase. L'idée reste stable, seule la forme change.}
\end{popfigure}

\courssub{La méthode de reformulation pas à pas}

\begin{methode}[Reformuler une idée en quatre étapes]
\begin{enumerate}
\item \textbf{Comprendre} : lire le passage deux fois et dégager l'idée en une question (« De quoi parle-t-on ? Que dit l'auteur à ce sujet ? »).
\item \textbf{Fermer le texte} : poser le document ou le retourner. Tant que le texte reste sous les yeux, on est tenté de le recopier.
\item \textbf{Redire avec ses mots} : écrire l'idée de mémoire, en mobilisant les quatre procédés (synonymie, changement grammatical, actif/passif, condensation).
\item \textbf{Comparer pour vérifier} : rouvrir le texte et contrôler la fidélité. Aucune idée ne doit être perdue, aucune idée étrangère ne doit être ajoutée, aucune suite de plus de trois mots identiques ne doit subsister.
\end{enumerate}
\end{methode}

\courssub{La neutralité énonciative}

Le résumé est un compte-rendu des idées \textbf{d'un autre}. Le rédacteur doit donc s'effacer : le « je » est interdit, car il introduit un point de vue personnel. On signale au contraire que les idées appartiennent à l'auteur grâce à des \cle{formules de neutralité énonciative} :

\begin{itemize}
\item « L'auteur affirme que... », « l'auteur dénonce... », « l'auteur souligne... » ;
\item « Selon le texte... », « d'après l'auteur... » ;
\item « Il ressort du passage que... », « il apparaît que... ».
\end{itemize}

\begin{attention}[Les trois interdits du résumé]
Trois fautes font perdre des points de manière systématique :
\begin{enumerate}
\item l'emploi du « je » (« je pense que l'auteur a raison » : c'est un commentaire, pas un résumé) ;
\item la citation recopiée sans guillemets ni signal explicite ;
\item l'ajout d'idées personnelles ou d'exemples absents du texte.
\end{enumerate}
\end{attention}

\begin{popfigure}
\begin{tikzpicture}[scale=1]
\node[draw=popdark, rounded corners, fill=popcream, minimum width=3.6cm, minimum height=1.1cm, align=center] (je) at (0,0) {\textbf{« Je pense que... »}\\ point de vue personnel};
\node[draw=popdark, rounded corners, fill=popcream, minimum width=3.6cm, minimum height=1.1cm, align=center] (cit) at (5.2,0) {\textbf{Citation recopiée}\\ sans guillemets};
\node[draw=popdark, rounded corners, fill=popcream, minimum width=3.6cm, minimum height=1.1cm, align=center] (ajout) at (10.4,0) {\textbf{Idée ajoutée}\\ absente du texte};
\draw[line width=2.5pt, poppink] (-1.2,-0.7) -- (1.2,0.7);
\draw[line width=2.5pt, poppink] (-1.2,0.7) -- (1.2,-0.7);
\draw[line width=2.5pt, poppink] (4.0,-0.7) -- (6.4,0.7);
\draw[line width=2.5pt, poppink] (4.0,0.7) -- (6.4,-0.7);
\draw[line width=2.5pt, poppink] (9.2,-0.7) -- (11.6,0.7);
\draw[line width=2.5pt, poppink] (9.2,0.7) -- (11.6,-0.7);
\end{tikzpicture}
\legende{Les trois interdits du résumé : le « je », la citation non signalée, l'ajout d'idées personnelles. Chacun est biffé car il trahit la fidélité ou la neutralité attendues.}
\end{popfigure}

\courssub{Entraînement guidé : reformuler les idées d'un paragraphe}

\begin{exoresolu}[Reformuler trois idées essentielles]
Énoncé. Voici un paragraphe argumentatif inspiré des débats que traverse \emph{Le Lion et la perle} : « Dans le village d'Ilujinle, les traditions organisent toute la vie sociale : elles règlent les mariages, les fêtes, les rapports entre les générations. Pourtant, l'école moderne séduit les jeunes, qui rêvent d'un autre avenir. Le conflit entre le vieux Baroka et l'instituteur Lakunle montre que la modernité ne peut pas simplement effacer la tradition : elle doit composer avec elle. » Dégager les trois idées essentielles et les reformuler chacune en une phrase, sans « je » et sans recopier le texte.
\tcblower
Solution détaillée. Étape 1 : on identifie les trois idées. (a) La tradition structure la vie du village. (b) L'école moderne attire la jeunesse. (c) Modernité et tradition doivent trouver un compromis.

Étape 2 : on reformule en appliquant les procédés.

(a) « Selon le texte, les coutumes gouvernent l'ensemble de l'existence au village, du mariage aux relations entre générations. » (synonymie : « règlent » $\rightarrow$ « gouvernent » ; condensation des exemples).

(b) « L'auteur souligne que les jeunes sont attirés par l'école, qui leur ouvre la perspective d'un avenir différent. » (passage à la voix passive : « l'école séduit les jeunes » $\rightarrow$ « les jeunes sont attirés par l'école »).

(c) « Il ressort de l'opposition entre Baroka et Lakunle que le progrès ne saurait supprimer les coutumes, mais doit s'accorder avec elles. » (changement grammatical : « composer avec » $\rightarrow$ « s'accorder avec » ; « effacer » $\rightarrow$ « supprimer »).

Vérification : aucune suite de plus de trois mots identiques, aucun « je », aucune idée ajoutée. La reformulation est fidèle.
\end{exoresolu}

\courssub{Exposé n° 1 : un thème de l'œuvre — tradition et modernité}

L'exposé oral mobilise les mêmes compétences que le résumé écrit : sélectionner les idées essentielles, les reformuler, les organiser. Pour présenter le thème \cle{tradition contre modernité} dans \emph{Le Lion et la perle}, l'élève suivra un plan en trois temps :

\begin{enumerate}
\item \textbf{Introduction} : situer l'œuvre (Wole Soyinka, dramaturge nigérian, prix Nobel de littérature en 1986) et annoncer la problématique : le village d'Ilujinle est partagé entre les coutumes incarnées par Baroka, le chef (le « Lion »), et les prétentions modernes de Lakunle, l'instituteur.
\item \textbf{Développement} : montrer d'abord le poids de la tradition (le mariage coutumier, la dot, la danse, le respect des anciens), puis les revendications de Lakunle (refus de la dot, discours sur le progrès), enfin la ruse de Baroka qui révèle les limites d'une modernité superficielle.
\item \textbf{Conclusion} : dégager la leçon de l'œuvre : Soyinka ne rejette ni la tradition ni la modernité ; il invite à un progrès enraciné dans la culture.
\end{enumerate}

Chaque passage clé cité doit être bref, signalé (« Soyinka écrit : "..." »), puis \textbf{reformulé} : « Autrement dit, l'auteur montre que... ». C'est cette reformulation orale qui prouve la compréhension.

\courssub{Exposé n° 2 : un personnage ou une scène}

Le second exposé porte sur un personnage (par exemple Sidi, la « perle » du village, ou Baroka) ou sur une scène (la danse du Lion, la scène de la ruse). La démarche est identique : présenter, analyser, conclure. L'exigence majeure reste le \cle{compte-rendu fidèle} des idées de l'auteur : l'élève rapporte ce que Soyinka fait dire à ses personnages, sans substituer son propre jugement. S'il donne un avis, il doit le marquer comme tel et le garder pour la conclusion.

\begin{methode}[Réussir son exposé oral]
\begin{enumerate}
\item Préparer une fiche de plan (pas un texte rédigé mot à mot).
\item Sélectionner deux ou trois citations courtes et les apprendre.
\item Préparer pour chaque citation une reformulation (« cela signifie que... »).
\item Répéter à voix haute en chronométrant (cinq minutes maximum).
\item Pendant l'exposé : regarder l'auditoire, parler lentement, annoncer les étapes du plan.
\end{enumerate}
\end{methode}

\begin{savaistu}[Reformuler : une compétence évaluée au Baccalauréat]
À l'oral comme à l'écrit, la capacité à reformuler est un critère explicite des grilles d'évaluation du Baccalauréat technique au Cameroun. Dans l'épreuve de contraction de texte, le correcteur vérifie précisément que le candidat n'a pas recopié le texte ; à l'oral, le jury valorise le candidat capable de redire une idée avec ses propres mots plutôt que de réciter. Entraîner la reformulation, c'est donc préparer directement l'examen.
\end{savaistu}

\courssub{Grille d'évaluation des exposés}

\begin{center}
\begin{tabularx}{\linewidth}{|l|X|c|}
\hline
\rowcolor{popblueL} \textbf{Critère} & \textbf{Ce qui est attendu} & \textbf{Points} \\
\hline
Exactitude du contenu & Les idées rapportées sont fidèles à l'œuvre, sans erreur ni invention. & /5 \\
\hline
Qualité de la reformulation & Les passages clés sont redits dans les mots de l'élève, sans récitation. & /5 \\
\hline
Clarté de l'expression & Voix audible, débit maîtrisé, plan annoncé et suivi, vocabulaire correct. & /5 \\
\hline
Respect du temps & L'exposé tient dans la durée impartie (cinq minutes). & /5 \\
\hline
\end{tabularx}
\end{center}

\begin{exercice}[À vous de jouer]
Relevez dans un chapitre étudié du \emph{Lion et la perle} une phrase frappante prononcée par Baroka ou par Lakunle. 1) Citez-la entre guillemets. 2) Reformulez-la en une phrase en utilisant au moins deux procédés (synonymie, changement grammatical, actif/passif). 3) Rédigez une phrase de transition qui l'intégrerait à un exposé, avec une formule de neutralité énonciative.
\end{exercice}

\begin{corrige}[Exercice 1]
Exemple de production attendue. 1) Citation : Lakunle refuse de payer la dot et déclare vouloir épouser Sidi « à la manière moderne ». 2) Reformulation : « L'instituteur entend se marier selon les coutumes occidentales et rejette le versement de la dot. » (changement grammatical : « manière moderne » $\rightarrow$ « coutumes occidentales » ; synonymie : « refuse » $\rightarrow$ « rejette »). 3) Phrase d'exposé : « Il ressort des propos de Lakunle que le personnage incarne une modernité revendiquée, que l'auteur présente toutefois avec ironie. » Toute production fidèle, sans « je » et sans recopie, est acceptée.
\end{corrige}

\begin{aretenir}[L'essentiel]
Reformuler, c'est redire fidèlement une idée avec d'autres mots et une autre construction, grâce à la synonymie, au changement de catégorie grammaticale, à la voix passive et à la condensation. La méthode : comprendre, fermer le texte, redire, vérifier. Le résumé comme l'exposé exigent la neutralité énonciative : pas de « je », des formules comme « l'auteur affirme » ou « il ressort que », et le bannissement de la paraphrase servile, de la citation non signalée et de l'ajout d'idées personnelles.
\end{aretenir}

\coursec{Bilan, compte-rendu et remédiation}

Après les exposés sur \emph{Le Lion et la perle} et le travail de reformulation des idées essentielles, notre parcours sur la contraction de texte touche à sa fin. Il est temps de faire le point : qu'avons-nous appris ? Quelles difficultés persistent ? Comment les corriger avant l'épreuve du Probatoire / Baccalauréat technique ? Cette dernière séance est décisive : c'est elle qui transforme des notions éparses en compétences solides.

\courssub{Compte-rendu de la séquence : synthèse collective des acquis}

Le compte-rendu collectif se construit au tableau, avec la participation de tous. Chaque élève propose une notion retenue, la classe la formule, et le professeur l'inscrit dans un tableau de synthèse. Voici le bilan auquel nous devons aboutir.

\begin{aretenir}[Les cinq acquis majeurs de la séquence]
\begin{enumerate}
\item \textbf{La structure argumentative du texte} : un texte argumentatif s'organise autour d'une \cle{thèse} (ce que l'auteur veut démontrer), d'\cle{arguments} (raisons avancées) et d'\cle{exemples} (illustrations concrètes). Repérer cette structure, c'est déjà avoir fait la moitié du résumé.
\item \textbf{Les techniques de réduction} : \cle{suppression} (éliminer exemples, répétitions, détails), \cle{généralisation} (remplacer une énumération par un terme générique : « machettes, houes, pelles » devient « outils agricoles ») et \cle{fusion} (condenser plusieurs phrases en une seule).
\item \textbf{La reformulation} : dire les idées de l'auteur avec ses propres mots, sans recopier les phrases du texte, sans ajouter d'opinion personnelle.
\item \textbf{Le lexique} : distinguer le \cle{lexique commun} (mots de tous les jours) du \cle{lexique spécialisé} (mots propres à un domaine), et reconnaître les \cle{néologismes} (mots nouveaux).
\item \textbf{Les figures d'amplification} : l'\cle{hyperbole} (exagération) et la \cle{gradation} (énumération en intensité croissante ou décroissante), qui renforcent l'argumentation.
\end{enumerate}
\end{aretenir}

\begin{exemplebox}[Synthèse en une phrase]
Résumer, c'est \textbf{retenir les idées essentielles} d'un texte (thèse et arguments), les \textbf{réduire} (suppression, généralisation, fusion) et les \textbf{reformuler} dans une langue correcte, en respectant la \textbf{longueur imposée}.
\end{exemplebox}

\courssub{Correction commentée d'un résumé type Bac}

En classe, nous avons résumé un texte argumentatif de 480 mots sur la scolarisation des filles, à réduire en 120 mots. Analysons collectivement deux productions d'élèves.

\begin{exoresolu}[Correction commentée d'un résumé d'élève]
\textbf{Production d'Amina (extrait)} : « L'auteur affirme que l'éducation des filles est une nécessité pour le développement de l'Afrique. Il montre d'abord qu'une femme instruite assure une meilleure santé et une meilleure éducation à ses enfants. Ensuite, il explique que la scolarisation des filles augmente les revenus des familles et des nations. Enfin, il dénonce les préjugés qui maintiennent les filles hors de l'école. » (78 mots, aucune phrase recopiée, langue correcte.)
\tcblower
\textbf{Commentaire} : la thèse est présente dès la première phrase ; les trois arguments sont conservés et hiérarchisés grâce aux connecteurs « d'abord », « ensuite », « enfin » ; les exemples chiffrés du texte ont été supprimés à juste titre ; la reformulation est totale. \textbf{Points forts} : fidélité, reformulation, correction. \textbf{Point à améliorer} : le résumé est un peu court (78 mots au lieu de 120) ; Amina aurait pu conserver la nuance finale de l'auteur sur le rôle des chefs traditionnels. Note indicative : 16/20.
\end{exoresolu}

\begin{attention}[Erreurs fréquentes observées dans les copies]
\begin{itemize}
\item \textbf{Recopier des phrases entières} du texte : c'est le défaut le plus sanctionné, il annule le critère de reformulation.
\item \textbf{Ajouter son avis} (« je pense que l'auteur a raison ») : le résumé doit rester neutre et fidèle.
\item \textbf{Négliger la relecture} : les fautes de grammaire et d'orthographe pèsent lourd dans la note finale (15 \% du barème, souvent décisifs entre deux mentions).
\item \textbf{Oublier la thèse} ou commencer par un détail : le correcteur doit voir l'idée principale dès la première phrase.
\end{itemize}
\end{attention}

\courssub{Le barème de l'épreuve au Probatoire / Baccalauréat technique}

Connaître le barème, c'est savoir où investir ses efforts. Au Probatoire et au Baccalauréat technique, le résumé est noté selon quatre critères.

\begin{popfigure}
\begin{center}
\begin{tikzpicture}[scale=1.3]
\def\angle{90}
\def\radius{2.2}
\foreach \percent/\color/\label in {40/popblue/Fidélité des idées (40\%), 30/poporange/Reformulation (30\%), 15/popgreen/Respect de la longueur (15\%), 15/poppink/Correction de la langue (15\%)} {
    \pgfmathsetmacro{\newangle}{\angle - \percent * 3.6}
    \draw[fill=\color, draw=white, thick] (0,0) -- (\angle:\radius) arc (\angle:\newangle:\radius) -- cycle;
    \pgfmathsetmacro{\midangle}{(\angle + \newangle) / 2}
    \node[label=\midangle:\label] at (\midangle:1.3*\radius) {};
    \global\let\angle\newangle
}
\end{tikzpicture}
\end{center}
\legende{Répartition du barème de notation du résumé au Probatoire / Baccalauréat technique : fidélité des idées 40 \%, reformulation 30 \%, respect de la longueur 15 \%, correction de la langue 15 \%.}
\end{popfigure}

\begin{exemplebox}[Un exemple chiffré pour comprendre]
Un élève rend un résumé \textbf{fidèle} et bien reformulé, mais de \textbf{180 mots au lieu des 120 imposés}. Sur le critère « respect de la consigne » (15 \% de la note, soit 3 points sur 20), il perd 2 points. S'il a en outre recopié deux phrases du texte, il perd encore 3 à 4 points sur la reformulation. Résultat : un travail pourtant intelligent plafonne à 12/20. La morale : \textbf{compter ses mots est aussi important que bien penser}.
\end{exemplebox}

\begin{methode}[S'auto-corriger en trois questions]
Avant de rendre votre copie, relisez votre résumé en vous posant trois questions :
\begin{enumerate}
\item \textbf{Ai-je gardé toutes les idées essentielles ?} Soulignez dans le texte la thèse et chaque argument, puis vérifiez que chacun apparaît dans votre résumé.
\item \textbf{Ai-je reformulé ?} Comparez phrase par phrase : si plus de trois mots consécutifs sont identiques au texte, réécrivez la phrase.
\item \textbf{Ai-je respecté la longueur ?} Comptez vos mots : une marge de 10 \% au-dessus ou en dessous est tolérée, pas davantage.
\end{enumerate}
\end{methode}

\courssub{Remédiation ciblée : trois ateliers}

La remédiation consiste à retravailler, en petits groupes, la difficulté précise repérée dans les copies. Chaque élève rejoint l'atelier qui correspond à son besoin.

\begin{popfigure}
\begin{center}
\begin{tikzpicture}[font=\small]
\node[draw=popdark, fill=popblueL, rounded corners, font=\bfseries\normalsize, minimum width=13.5cm, minimum height=0.9cm] (t) at (0,2.8) {Tableau de remédiation de la contraction de texte};
\node[draw, fill=popblue!40, font=\bfseries, minimum width=4.2cm, minimum height=0.9cm, align=center] (a) at (-4.6,1.7) {Erreur observée};
\node[draw, fill=poporange!40, font=\bfseries, minimum width=4.2cm, minimum height=0.9cm, align=center] (b) at (0,1.7) {Cause probable};
\node[draw, fill=popgreen!40, font=\bfseries, minimum width=4.2cm, minimum height=0.9cm, align=center] (c) at (4.6,1.7) {Activité corrective};
\node[draw, fill=white, text width=3.9cm, minimum width=4.2cm, minimum height=1.6cm, align=center] (a1) at (-4.6,0.3) {Résumé trop long, détails conservés};
\node[draw, fill=white, text width=3.9cm, minimum width=4.2cm, minimum height=1.6cm, align=center] (b1) at (0,0.3) {Difficulté à distinguer l'essentiel de l'accessoire};
\node[draw, fill=white, text width=3.9cm, minimum width=4.2cm, minimum height=1.6cm, align=center] (c1) at (4.6,0.3) {Atelier suppression / généralisation};
\node[draw, fill=white, text width=3.9cm, minimum width=4.2cm, minimum height=1.6cm, align=center] (a2) at (-4.6,-1.5) {Phrases recopiées du texte};
\node[draw, fill=white, text width=3.9cm, minimum width=4.2cm, minimum height=1.6cm, align=center] (b2) at (0,-1.5) {Peur de trahir le texte, manque de vocabulaire};
\node[draw, fill=white, text width=3.9cm, minimum width=4.2cm, minimum height=1.6cm, align=center] (c2) at (4.6,-1.5) {Atelier reformulation};
\node[draw, fill=white, text width=3.9cm, minimum width=4.2cm, minimum height=1.6cm, align=center] (a3) at (-4.6,-3.3) {Vocabulaire pauvre, figures non comprises};
\node[draw, fill=white, text width=3.9cm, minimum width=4.2cm, minimum height=1.6cm, align=center] (b3) at (0,-3.3) {Lecture insuffisante, lexique non maîtrisé};
\node[draw, fill=white, text width=3.9cm, minimum width=4.2cm, minimum height=1.6cm, align=center] (c3) at (4.6,-3.3) {Atelier lexique et figures d'amplification};
\end{tikzpicture}
\end{center}
\legende{Tableau de remédiation : pour chaque erreur observée dans les copies, une cause probable et une activité corrective ciblée.}
\end{popfigure}

\textbf{Atelier 1 : suppression et généralisation} (pour les résumés trop longs). Les élèves reçoivent un paragraphe de 100 mots à réduire à 40 mots en deux étapes : d'abord rayer au crayon rouge tout ce qui peut disparaître (exemples, répétitions, adjectifs inutiles), puis remplacer les énumérations par un terme générique. On compare ensuite les versions : le groupe retient celle qui garde le plus d'idées avec le moins de mots.

\textbf{Atelier 2 : reformulation} (pour les élèves qui recopient). Exercice du « texte caché » : l'élève lit une phrase du texte, ferme le document, puis réécrit l'idée de mémoire avec ses propres mots. On s'entraîne aussi avec des substitutions : remplacer un verbe par un nom (« il proteste » devient « sa protestation »), une phrase active par une phrase passive, un mot par un synonyme.

\textbf{Atelier 3 : lexique et figures d'amplification}. Les élèves classent une liste de mots en deux colonnes (lexique commun / lexique spécialisé), identifient des néologismes dans un article récent (par exemple « infodémie », « déconfinement »), puis repèrent hyperboles et gradations dans un extrait de \emph{Le Lion et la perle} et expliquent leur effet sur le lecteur.

\begin{savaistu}[La remédiation, une recommandation officielle]
Les inspecteurs MINESEC recommandent de consacrer \textbf{au moins une séance entière à la remédiation} après chaque évaluation de contraction de texte. Pourquoi ? Parce que refaire le même exercice sans comprendre ses erreurs ne fait progresser personne : c'est l'analyse de l'erreur, suivie d'un exercice ciblé, qui transforme durablement la compétence.
\end{savaistu}

\courssub{Évaluation sommative et auto-évaluation}

\begin{exercice}[Évaluation sommative : résumé en conditions d'examen]
Texte argumentatif inédit de 400 mots distribué en classe (sujet : « Faut-il interdire les téléphones portables à l'école ? »). \textbf{Consigne} : résumez ce texte en 100 mots (marge de 10 \% tolérée). Indiquez le nombre de mots à la fin. Durée : 1 heure. Barème : fidélité des idées 8 points, reformulation 6 points, respect de la longueur 3 points, correction de la langue 3 points.
\end{exercice}

Après l'évaluation, chaque élève remplit sa grille d'auto-évaluation en cochant honnêtement chaque objectif.

\begin{center}
\begin{tabularx}{\linewidth}{|X|c|c|}
\hline
\rowcolor{popblueL} \textbf{Objectif de la leçon} & \textbf{Je maîtrise} & \textbf{À retravailler} \\
\hline
Repérer la thèse et les arguments d'un texte & $\square$ & $\square$ \\
\hline
Réduire par suppression, généralisation, fusion & $\square$ & $\square$ \\
\hline
Reformuler sans recopier le texte & $\square$ & $\square$ \\
\hline
Respecter la longueur imposée & $\square$ & $\square$ \\
\hline
Distinguer lexique commun et spécialisé, néologismes & $\square$ & $\square$ \\
\hline
Identifier hyperbole et gradation & $\square$ & $\square$ \\
\hline
Relire et corriger ma langue & $\square$ & $\square$ \\
\hline
\end{tabularx}
\end{center}

\begin{corrige}[Repères pour l'évaluation sommative]
Un bon résumé du texte sur les téléphones portables doit contenir : la thèse de l'auteur (par exemple : l'interdiction est nécessaire mais insuffisante sans éducation au numérique), les arguments principaux (nuisance à la concentration, risques de triche, mais utilité pédagogique possible), et la conclusion de l'auteur. Toute opinion personnelle, tout exemple recopié et tout dépassement de longueur sont sanctionnés selon le barème. La grille d'auto-évaluation permet ensuite à chacun de s'inscrire à l'atelier de remédiation adapté.
\end{corrige}

\begin{aretenir}[L'essentiel de la séance]
Le bilan de la séquence repose sur trois gestes : \textbf{synthétiser} les acquis (structure argumentative, réduction, reformulation, lexique, figures), \textbf{analyser} ses erreurs à la lumière du barème officiel (40/30/15/15), et \textbf{remédier} par des ateliers ciblés. L'auto-évaluation finale fait de chaque élève l'acteur de sa progression : savoir ce que l'on maîtrise et ce qu'il reste à travailler est la première compétence du candidat au Probatoire / Baccalauréat technique.
\end{aretenir}

\newpage

\coursec{Activité d'intégration}

\coursec{Leçon 13 : Contraction de texte : le résumé (reformuler les idées d'un texte à résumer)}

\courssub{Objectifs de l'activité}

À l'issue de cette activité d'intégration, l'élève sera capable de :
\begin{itemize}
\item dégager la \cle{structure argumentative} d'un texte de presse (thèse, arguments, exemples, concession) ;
\item appliquer les \cle{techniques de réduction} (suppression des exemples, des répétitions, des figures d'amplification, fusion de phrases) pour ramener un texte au quart de sa longueur ;
\item \cle{reformuler} fidèlement les idées essentielles avec ses propres mots, sans copier les phrases du texte ;
\item repérer et traiter le \cle{lexique spécialisé} et les \cle{néologismes} dans la reformulation (conserver l'indispensable, expliciter le reste) ;
\item rédiger un résumé respectant les contraintes de l'épreuve de contraction de texte au Probatoire technique (longueur imposée, fidélité, correction de la langue, cohérence) ;
\item évaluer la production d'un camarade à l'aide d'un barème officiel et justifier sa notation de manière argumentée.
\end{itemize}

\courssub{Situation d'apprentissage}

Le lycée technique de Nkongsamba organise sa semaine de la presse. Le club « Français et citoyenneté », animé par les élèves de Première technique, doit publier dans le journal mural une synthèse des grands débats économiques nationaux. Le professeur de français distribue à chaque groupe l'éditorial ci-dessous, paru dans un quotidien camerounais. Il compte \textbf{près de 400 mots}. Chaque élève devra en produire un résumé de \textbf{100 mots (plus ou moins 10)}, soit le quart du texte, conformément aux exigences de l'épreuve de contraction de texte.

\textbf{Texte à résumer :}

\emph{Le mobile money, moteur silencieux de notre économie}

\begin{quote}
Il y a quinze ans encore, envoyer de l'argent de Douala à Maroua relevait de l'expédition : il fallait confier des billets à un chauffeur de bus, prier pour qu'ils arrivent, attendre parfois une semaine. Aujourd'hui, une ménagère de Bafoussam transfère dix mille francs à son fils étudiant à Yaoundé en moins d'une minute, depuis sa boutique. Le mobile money a bouleversé notre quotidien, et il est temps de reconnaître qu'il est devenu un véritable moteur de l'économie nationale.

D'abord, le mobile money favorise l'inclusion financière des populations. Selon la Banque des États de l'Afrique centrale, plus de soixante pour cent des Camerounais adultes possèdent désormais un compte de monnaie électronique, alors que moins de vingt pour cent détiennent un compte bancaire classique. Le paysan de l'Ouest, la commerçante du marché Mokolo, le moto-taximan de Bertoua : tous encaissent, épargnent et paient sans jamais franchir la porte d'une banque. C'est une révolution silencieuse, immense, irréversible.

Ensuite, ce service dynamise le commerce. Les transactions sont instantanées, sécurisées, disponibles jour et nuit. Un grossiste de Douala règle son fournisseur de Ngaoundéré sans déplacement ; une élève paie ses frais de scolarité sans transporter d'espèces. Les files d'attente interminables devant les guichets appartiennent au passé, ce passé encombrant, lent, épuisant que nos parents ont connu.

Enfin, l'État lui-même y trouve son compte : les taxes sur les transactions alimentent le Trésor public, et la traçabilité des paiements électroniques combat la fraude et la corruption, ces fléaux qui saignent notre économie depuis des décennies.

Certes, des défis demeurent : coût des retraits, pannes de réseau, arnaques par messages frauduleux. Mais ces obstacles, réels, ne doivent pas masquer l'essentiel. Le mobile money n'est pas une mode passagère : il est l'outil de développement que le Cameroun attendait. Encourageons-le, encadrons-le, généralisons-le. L'argent liquide a fait son temps ; l'avenir, lui, tient dans la poche.
\end{quote}

\courssub{Consignes de travail}

\textbf{Travail de groupe (par quatre), puis travail individuel.}

\begin{enumerate}
\item \textbf{Comprendre et structurer.} Lisez le texte deux fois. Soulignez en rouge la phrase qui exprime la \cle{thèse} de l'éditorialiste. Numérotez dans la marge les trois arguments qui la soutiennent et indiquez le connecteur qui introduit chacun d'eux. Identifiez le paragraphe de concession.
\item \textbf{Réduire.} Barrez au crayon les éléments à éliminer : exemples concrets (personnages, lieux, chiffres d'illustration), répétitions, figures d'amplification (hyperboles, gradations). Justifiez chaque suppression par une note en marge.
\item \textbf{Traiter le lexique.} Relevez cinq termes appartenant au lexique spécialisé de l'économie et de la finance, et deux néologismes ou mots récents liés au numérique. Pour chacun, proposez une reformulation ou une explication brève utilisable dans le résumé.
\item \textbf{Rédiger.} Rédigez \textbf{individuellement} un résumé de 100 mots (plus ou moins 10) qui reprend la thèse et les trois arguments, dans vos propres phrases, sans exemple ni citation. Indiquez le nombre de mots à la fin.
\item \textbf{Corriger.} Échangez votre copie avec un camarade d'un autre groupe. Notez-la sur 20 à l'aide du barème ci-dessous, puis justifiez oralement la note attribuée en citant les forces et les faiblesses repérées.
\end{enumerate}

\begin{center}
\begin{tabularx}{\linewidth}{|l|X|c|}
\hline
\rowcolor{popblueL} \textbf{Critère} & \textbf{Attendu} & \textbf{Note} \\
\hline
Fidélité & Thèse et arguments du texte respectés, sans idée ajoutée ni opinion personnelle & /6 \\
\hline
Reformulation & Phrases propres au rédacteur, aucune phrase copiée-collée, syntaxe variée & /6 \\
\hline
Longueur & 90 à 110 mots, nombre exact indiqué en fin de résumé & /4 \\
\hline
Langue & Orthographe, grammaire, ponctuation, connecteurs logiques, richesse lexicale & /4 \\
\hline
\end{tabularx}
\end{center}

\courssub{Ressources à mobiliser}

\begin{aretenir}[Les outils de la contraction de texte]
\begin{itemize}
\item \textbf{Structure argumentative :} une \cle{thèse} (idée défendue), des \cle{arguments} (raisons qui prouvent), des \cle{exemples} (illustrations à supprimer), parfois une \cle{concession} (objection admise puis réfutée, à nuancer ou supprimer), des \cle{connecteurs} (d'abord, ensuite, enfin, certes, mais) qui balisent le raisonnement.
\item \textbf{Techniques de réduction :} supprimer exemples, anecdotes, chiffres illustratifs, répétitions, citations, figures d'amplification (hyperbole, gradation) ; remplacer une énumération par un terme générique (hyperonyme) ; fondre plusieurs phrases en une seule (subordination, nominalisation).
\item \textbf{Reformulation :} changer les mots (synonymes, termes génériques, antonymes avec négation) et les constructions (actif $\leftrightarrow$ passif, verbal $\leftrightarrow$ nominal), sans jamais changer le sens ni ajouter d'information.
\item \textbf{Lexique :} distinguer le \cle{lexique commun} (mots de tous les jours) du \cle{lexique spécialisé} (mots propres à un domaine : économie, finance, droit) ; un \cle{néologisme} est un mot nouveau entré dans la langue (ex. : \emph{mobile money}, \emph{monnaie électronique}, \emph{fintech}). Conserver les termes techniques indispensables à la précision ; expliciter ou remplacer les autres.
\item \textbf{Règle du quart :} texte source $\approx$ 400 mots $\rightarrow$ résumé cible = 100 mots $\pm$ 10\% (soit 90 à 110 mots). Le non-respect de la longueur est pénalisant.
\end{itemize}
\end{aretenir}

\courssub{Méthode pas à pas : Rédiger un résumé de texte argumentatif (contraction)}

\begin{methode}[Les 5 étapes de la contraction de texte]
\begin{enumerate}
\item \textbf{Analyser} : lire 2 fois ; identifier la thèse (souvent en début ou fin de texte) ; repérer les arguments (connecteurs : \emph{d'abord, ensuite, enfin, par ailleurs}) ; repérer la concession (\emph{certes, il est vrai que}) ; numéroter les paragraphes.
\item \textbf{Structurer} : faire le squelette du texte (Thèse -- Arg 1 -- Arg 2 -- Arg 3 -- Concession/Conclusion). Vérifier que chaque argument est bien une idée-force, pas un exemple.
\item \textbf{Réduire} : au crayon, barrer tout ce qui n'est pas idée essentielle : exemples (noms propres, chiffres, anecdotes), figures de style (hyperboles, gradations, images), répétitions, connecteurs de lourdeur. Garder les connecteurs logiques majeurs.
\item \textbf{Reformuler \& Rédiger} : réécrire le squelette avec ses propres phrases. Grouper les idées. Utiliser le discours indirect (\emph{L'auteur affirme que...}, \emph{Selon l'éditorialiste...}). Vérifier l'enchaînement logique.
\item \textbf{Contrôler} : compter les mots (compter les mots-outils : \emph{le, la, de, et} comptent pour 1) ; vérifier la fidélité (pas de contresens, pas d'oubli d'argument) ; corriger la langue (accords, temps, ponctuation) ; indiquer le nombre de mots.
\end{enumerate}
\end{methode}

\courssub{Attention : Pièges fréquents à l'épreuve}

\begin{attention}[Erreurs éliminatoires ou pénalisantes]
\begin{enumerate}
\item \textbf{Copier-coller :} recopier des phrases entières du texte (note 0 en reformulation).
\item \textbf{Hors-sujet / Ajout :} donner son avis personnel (\emph{« Je pense que... »}, \emph{« C'est vrai car... »}) ou ajouter des connaissances extérieures.
\item \textbf{Oubli d'argument :} ne résumer que l'introduction et la conclusion, ou sauter un des trois arguments principaux.
\item \textbf{Garder les exemples :} conserver les noms propres (Bafoussam, Mokolo), les chiffres (60\%, 20\%), les anecdotes (le chauffeur de bus).
\item \textbf{Non-respect de la longueur :} résumé $<$ 90 mots (incomplet) ou $>$ 110 mots (non contracté). La marge de 10\% est une tolérance, pas un objectif.
\item \textbf{Mauvaise gestion de la concession :} soit la développer comme un 4e argument, soit l'ignorer totalement sans nuance. La réduire à une subordonnée (\emph{bien que...}, \emph{malgré...}) est la norme.
\end{enumerate}
\end{attention}

\courssub{Le saviez-vous ? Contexte camerounais et enjeux}

\begin{savaistu}[Mobile Money et inclusion financière au Cameroun]
Le texte cite la \textbf{BEAC} (Banque des États de l'Afrique centrale), institution monétaire commune à 6 pays (Cameroun, Tchad, RCA, Congo, Gabon, Guinée équatoriale). Les chiffres évoqués (60\% de comptes \emph{mobile money} vs 20\% bancaires) reflètent la réalité 2023-2024 : le Cameroun est leader en zone CEMAC pour la monnaie électronique (Orange Money, MTN MoMo, Express Union Mobile, etc.). La \emph{traçabilité} des flux est un levier majeur pour la \emph{mobilisation des recettes fiscales} (taxes sur les retraits/transferts) et la lutte contre l'économie informelle. Le néologisme \emph{mobile money} (anglicisme) est d'usage courant ; le terme officiel français est \emph{monnaie électronique} ou \emph{services financiers mobiles}. Le \emph{moto-taximan} (conducteur de taxi-moto) illustre l'inclusion du secteur informel.
\end{savaistu}

\courssub{Démarche et corrigé commenté}

\begin{exoresolu}[Résumé de l'éditorial « Le mobile money, moteur silencieux de notre économie »]
Énoncé : produire le résumé de référence (100 mots $\pm$ 10) de l'éditorial distribué, en suivant les cinq consignes.
\tcblower
\textbf{Étape 1 : dégager la structure argumentative.}

La thèse se trouve à la fin du premier paragraphe : \emph{« il est devenu un véritable moteur de l'économie nationale »}. L'éditorialiste affirme que le mobile money fait progresser l'économie camerounaise.

Les trois arguments sont signalés par des connecteurs d'énumération :
\begin{itemize}
\item \emph{D'abord} : le mobile money favorise l'inclusion financière des populations (accès aux services financiers sans compte bancaire) ;
\item \emph{Ensuite} : il dynamise le commerce grâce à des transactions rapides, sûres et permanentes ;
\item \emph{Enfin} : l'État y gagne (recettes fiscales via les taxes, lutte contre la fraude et la corruption grâce à la traçabilité).
\end{itemize}
Le paragraphe \emph{« Certes, des défis demeurent... »} est une \cle{concession} : l'auteur admet des limites (coût, pannes, arnaques) pour mieux les relativiser (\emph{« ne doivent pas masquer l'essentiel »}). La conclusion reprend la thèse sous forme d'appel à l'action (\emph{Encourageons-le, encadrons-le, généralisons-le}).

\textbf{Étape 2 : réduire (biffage sélectif).}

On barre :
\begin{itemize}
\item L'anecdote d'introduction (Douala--Maroua, chauffeur de bus, semaine d'attente) ;
\item Les personnages-types et lieux de l'argument 1 (ménagère de Bafoussam, paysan de l'Ouest, commerçante du marché Mokolo, moto-taximan de Bertoua) et les statistiques précises (60\%, 20\%) ;
\item Les exemples de transactions de l'argument 2 (grossiste Douala/Ngaoundéré, élève, frais de scolarité) ;
\item La description du passé bancaire (files d'attente, passé encombrant/lent/épuisant) ;
\item Les figures d'amplification : gradations (\emph{« silencieuse, immense, irréversible »} ; \emph{« encourageons-le, encadrons-le, généralisons-le »}), hyperbole (\emph{« fléaux qui saignent notre économie »}), image finale (\emph{« l'avenir tient dans la poche »}).
\end{itemize}
Ces éléments illustrent ou amplifient : ils ne font pas partie des idées essentielles à conserver.

\textbf{Étape 3 : traiter le lexique spécialisé et les néologismes.}

\begin{center}
\begin{tabularx}{\linewidth}{|l|l|X|}
\hline
\rowcolor{popblueL} \textbf{Terme repéré} & \textbf{Nature} & \textbf{Reformulation / Explication pour le résumé} \\
\hline
Inclusion financière & Lexique spécialisé (économie) & Accès de tous aux services financiers (compte, épargne, paiement) \\
\hline
Monnaie électronique & Lexique spécialisé (finance) & Argent dématérialisé, stocké sur téléphone \\
\hline
Compte bancaire classique & Lexique spécialisé (banque) & Compte en banque (terme commun, à garder) \\
\hline
Transactions & Lexique spécialisé (économie) & Opérations de paiement, transferts \\
\hline
Trésor public & Lexique spécialisé (finances publiques) & Caisse de l'État, finances publiques \\
\hline
Traçabilité & Lexique spécialisé (technique/éco) & Possibilité de suivre / tracer chaque paiement \\
\hline
Mobile money & Néologisme (anglicisme) & Paiement par téléphone / monnaie mobile (à conserver, terme clé) \\
\hline
Moto-taximan & Néologisme / Mot récent (métier) & Conducteur de taxi-moto (secteur informel, à supprimer car exemple) \\
\hline
\end{tabularx}
\end{center}
\textit{Stratégie :} on conserve \emph{mobile money} (terme incontournable du sujet) et \emph{inclusion financière} (notion centrale). On reformule \emph{monnaie électronique} par \emph{argent dématérialisé} ou \emph{services financiers mobiles}, \emph{Trésor public} par \emph{finances de l'État}, \emph{traçabilité} par \emph{suivi des opérations}.

\textbf{Étape 4 : rédiger le résumé de référence.}

\begin{quote}
Selon l'éditorialiste, le paiement par téléphone portable est devenu un véritable moteur de l'économie camerounaise. D'abord, il permet à la majorité des habitants, même sans compte bancaire, d'accéder aux services financiers. Ensuite, il facilite les échanges commerciaux grâce à des paiements rapides, sûrs et permanents. Enfin, il profite à l'État, qui perçoit des taxes et lutte mieux contre la fraude grâce au suivi des opérations. Malgré quelques difficultés, l'auteur invite à encourager et à encadrer ce service, qu'il considère comme un outil de développement décisif pour l'avenir du pays. (100 mots)
\end{quote}

\textbf{Vérification du contrat :}
\begin{itemize}
\item \textbf{Fidélité :} thèse + 3 arguments + concession nuancée + conclusion. Aucune idée extérieure.
\item \textbf{Reformulation :} zéro phrase copiée. \emph{« Le mobile money »} $\rightarrow$ \emph{« le paiement par téléphone portable »} ; \emph{« inclusion financière »} $\rightarrow$ \emph{« accéder aux services financiers »} ; \emph{« dynamise le commerce »} $\rightarrow$ \emph{« facilite les échanges commerciaux »} ; \emph{« traçabilité »} $\rightarrow$ \emph{« suivi des opérations »}.
\item \textbf{Longueur :} exactement 100 mots (dans la fourchette 90--110).
\item \textbf{Langue :} connecteurs \emph{d'abord, ensuite, enfin} ; subordination (\emph{qui perçoit}, \emph{grâce à}) ; discours indirect (\emph{Selon l'éditorialiste}, \emph{l'auteur invite à}).
\end{itemize}

\textbf{Étape 5 : correction croisée (exemple d'application du barème).}

Copie d'un camarade : \emph{« L'éditorialiste dit que le mobile money est un moteur de l'économie. D'abord inclusion financière pour tous sans banque. Ensuite commerce dynamisé par transactions instantanées. Enfin État gagne taxes et traçabilité contre fraude. Malgré défis, faut encourager. Outil développement. »} (62 mots)

\begin{itemize}
\item \textbf{Fidélité 3/6 :} arguments 1 et 3 trop télégraphiques (manque de reformulation syntaxique), argument 2 incomplet (pas de « sûrs/permanents »), concession absente.
\item \textbf{Reformulation 2/6 :} style télégraphique (nominalisation excessive, pas de phrases complètes), limite du copier-coller lexical (\emph{inclusion financière, transactions instantanées, traçabilité}).
\item \textbf{Longueur 0/4 :} 62 mots $<$ 90 mots (hors tolérance). Pénalité majeure.
\item \textbf{Langue 2/4 :} pas de fautes d'orthographe mais syntaxe défaillante (pas de verbes conjugués, pas de connecteurs fluides).
\item \textbf{Total : 7/20.} Le correcteur justifie : « Le résumé est trop court, les idées ne sont pas développées en phrases, la concession a disparu. »
\end{itemize}
\end{exoresolu}

\courssub{Exercice d'entraînement (devoir / séance suivante)}

\begin{exercice}[Contraction de texte : « L'agriculture, pilier de l'industrialisation »]
Le texte ci-dessous compte \textbf{360 mots}. Produisez un résumé de \textbf{90 mots $\pm$ 10} (soit 81 à 99 mots) en appliquant la méthode en 5 étapes. Indiquez le nombre de mots.

\begin{quote}
\emph{Depuis l'indépendance, le Cameroun a fait le pari de l'agriculture comme socle de son développement. Pourtant, force est de constater que le secteur peine à jouer pleinement son rôle de pourvoyeur de richesses et de pourvoyeur d'emplois. La transformation structurelle de l'économie, promise depuis des décennies, reste en grande partie un vœu pieux.}

\emph{D'abord, la productivité agricole reste faible. Les exploitations sont en majorité familiales, de petite taille, mal équipées, dépendantes de la pluviométrie et peu mécanisées. L'usage d'intrants améliorés (semences certifiées, engrais) est marginal. Résultat : les rendements à l'hectare pour le maïs, le cacao ou le coton sont bien inférieurs aux potentialités agronomiques et aux moyennes mondiales.}

\emph{Ensuite, la commercialisation est désorganisée. L'absence de pistes rurales praticables toute l'année isole les bassins de production des centres de consommation et des ports d'exportation. Les pertes post-récolte atteignent parfois 30 à 40 \% pour les produits périssables (tomates, plantains, mangues). Les intermédiaires captent l'essentiel de la valeur ajoutée, laissant au producteur un revenu de misère.}

\emph{Enfin, la transformation locale est quasi inexistante. Nous exportons du cacao brut, du coton brut, du bois brut, pour réimporter du chocolat, du tissu, du meuble. Cette absence d'industries agro-alimentaires et de première transformation nous prive d'emplois industriels, de devises et de souveraineté alimentaire.}

\emph{Certes, des programmes existent : PIDMA, PADFA, ACEFA, et la récente Stratégie Nationale de Développement (SND30). Mais leur mise en œuvre souffre de lenteurs administratives, de gouvernance approximative et de financements insuffisants.}

\emph{Il est urgent de passer des discours aux actes. Investir massivement dans les infrastructures routières et énergétiques rurales, garantir l'accès au crédit et aux intrants, créer des zones agro-industrielles intégrées : telles sont les conditions pour que l'agriculture devienne enfin le vrai moteur de l'industrialisation camerounaise.}
\end{quote}
\end{exercice}

\courssub{Corrigé type de l'exercice d'entraînement}

\begin{corrige}[Exercice n°1 : L'agriculture, pilier de l'industrialisation]
\textbf{Structure :} Thèse (L'agriculture peine à être le moteur du développement/industrialisation). Arg 1 (Faible productivité : exploitations familiales, peu mécanisées, faibles rendements). Arg 2 (Commercialisation désorganisée : enclavement, pertes post-récolte, intermédiaires). Arg 3 (Transformation locale inexistante : export brut / import transformé, perte d'emplois/devises). Concession (Programmes existants mais mise en œuvre défaillante). Conclusion (Urgence d'actes : infrastructures, crédit, zones agro-industrielles).

\textbf{Résumé de référence (92 mots) :}
L'éditorialiste affirme que l'agriculture, pourtant socle annoncé du développement camerounais, ne joue pas encore son rôle de moteur de l'industrialisation. D'abord, la productivité reste faible du fait de l'archaïsme des exploitations familiales, peu mécanisées et faiblement pourvues en intrants. Ensuite, la commercialisation est handicapée par l'enclavement des zones de production, engendrant d'énormes pertes post-récolte et l'accaparement de la valeur par les intermédiaires. Enfin, l'absence de transformation locale contraint à exporter des matières premières brutes pour réimporter des produits finis, perdant emplois et devises. Malgré l'existence de programmes, leur application défaillante impose un passage urgent aux actes : infrastructures, crédit et zones agro-industrielles. (92 mots)
\end{corrige}

\courssub{Bilan de l'activité}

Cette activité a permis de mobiliser, dans une situation proche de l'examen, l'ensemble des savoir-faire de la leçon. Les élèves ont d'abord confirmé qu'un texte argumentatif de presse s'organise autour d'une \cle{thèse} soutenue par des \cle{arguments} hiérarchisés par des connecteurs (\emph{d'abord, ensuite, enfin}) : repérer cette armature est la condition \emph{sine qua non} de tout bon résumé. Ils ont ensuite vérifié concrètement que la \cle{réduction} consiste à éliminer ce qui illustre ou amplifie — exemples, chiffres, anecdotes, hyperboles et gradations — et non ce qui prouve. Le travail sur le lexique a montré qu'un terme spécialisé indispensable à la fidélité (\emph{mobile money}, \emph{inclusion financière}) se conserve, tandis qu'un néologisme ou un terme technique secondaire peut être explicité ou remplacé par un hyperonyme. Enfin, la correction croisée avec barème a développé l'esprit critique et l'auto-évaluation : justifier une note, c'est apprendre à relire et à améliorer sa propre copie. La règle d'or à retenir pour l'épreuve de contraction de texte au Probatoire technique : \textbf{dire toutes les idées essentielles, rien que les idées essentielles, avec ses propres mots, dans le nombre de mots imposé}.

\newpage

\coursec{Méthodes et exercices résolus}

\courssub{Méthode 1 : Repérer la structure argumentative d'un texte}

\begin{methode}[Dégager la structure argumentative d'un texte]
Pour résumer un texte argumentatif, il faut d'abord en dégager l'ossature logique. Procédez en étapes :
\begin{enumerate}
\item \textbf{Identifier le thème} : de quoi parle le texte ? Formulez-le en un groupe nominal (ex. : « le conflit entre tradition et modernité »).
\item \textbf{Dégager la thèse} : quelle position l'auteur défend-il ? Cherchez les marqueurs d'opinion (\emph{il faut, je pense que, sans doute}) et la conclusion du texte.
\item \textbf{Repérer les arguments} : chaque argument est une raison qui soutient la thèse. Ils sont souvent introduits par des connecteurs logiques (\emph{d'abord, en effet, ensuite, enfin, par conséquent}).
\item \textbf{Relever les exemples} : ils illustrent les arguments mais ne sont pas essentiels ; au résumé, on les supprime ou on les réduit à une mention.
\item \textbf{Construire le schéma} : thème $\rightarrow$ thèse $\rightarrow$ argument 1 $\rightarrow$ argument 2 $\rightarrow$ conclusion.
\end{enumerate}
\textbf{Réflexe} : soulignez au brouillon les connecteurs logiques ; ce sont les poteaux indicateurs de la structure.
\end{methode}

\begin{exoresolu}[Analyser la structure d'un extrait argumentatif]
\textbf{Énoncé.} Lisez l'extrait suivant inspiré des débats du \emph{Lion et la perle} : « Certains prétendent que le progrès technique détruit nos traditions. C'est faux. D'abord, la modernité, comme le montre Lakunle, apporte l'instruction à tous les enfants du village. Ensuite, elle améliore les conditions de vie : l'eau courante et les routes désenclavent nos campagnes. Le progrès est donc une chance, à condition de ne pas renier notre identité. » Dégagez le thème, la thèse et les arguments de ce texte.
\tcblower
\textbf{Solution détaillée.}
\begin{enumerate}
\item \textbf{Thème} : le rapport entre le progrès technique et les traditions africaines (le conflit modernité/tradition).
\item \textbf{Thèse} : l'auteur réfute l'opinion adverse (« C'est faux ») et défend l'idée que « le progrès est une chance ». La thèse apparaît dans la conclusion, marquée par \emph{donc}.
\item \textbf{Arguments} : deux arguments, signalés par \emph{d'abord} et \emph{ensuite} :
\begin{itemize}
\item argument 1 : la modernité instruit tous les enfants (exemple : Lakunle) ;
\item argument 2 : elle améliore les conditions de vie (exemples : eau courante, routes).
\end{itemize}
\item \textbf{Nuance finale} : « à condition de ne pas renier notre identité » est une restriction qui encadre la thèse ; elle doit figurer dans le résumé car elle fait partie de la pensée de l'auteur.
\end{enumerate}
\textbf{Schéma} : thème (progrès/tradition) $\rightarrow$ thèse (le progrès est une chance) $\rightarrow$ arg. 1 (instruction) $\rightarrow$ arg. 2 (conditions de vie) $\rightarrow$ nuance (garder son identité).
\textbf{Conclusion} : le texte défend une modernité maîtrisée, qui s'accompagne de la fidélité à l'identité culturelle.
\end{exoresolu}

\courssub{Méthode 2 : Appliquer les techniques de réduction}

\begin{methode}[Réduire un texte : les quatre techniques]
La contraction impose de reformuler en respectant un nombre de mots imposé (souvent le quart du texte, avec une marge de $\pm$ 10 \%). Quatre techniques de réduction :
\begin{enumerate}
\item \textbf{La suppression} : éliminer les exemples, les répétitions, les citations, les détails descriptifs, les figures d'amplification.
\item \textbf{La généralisation} : remplacer une énumération par un terme générique (« manguiers, palmiers, bananiers » $\rightarrow$ « les arbres fruitiers »).
\item \textbf{La condensation} : remplacer une proposition ou une périphrase par un mot (« le chef qui dirige le village » $\rightarrow$ « le chef du village » ; « d'une manière courageuse » $\rightarrow$ « courageusement »).
\item \textbf{La transformation} : convertir une phrase complexe en phrase simple, une subordonnée en complément (« Quand il arriva, il pleuvait » $\rightarrow$ « Il pleuvait à son arrivée »).
\end{enumerate}
\textbf{Réflexes} : comptez les mots du texte source, calculez le quota, puis comptez vos mots à chaque paragraphe rédigé. Ne changez jamais le sens ni l'ordre logique des idées.
\end{methode}

\begin{exoresolu}[Réduire un passage de 80 mots à 20 mots]
\textbf{Énoncé.} Réduisez à environ 20 mots le passage suivant (80 mots) : « Baroka, le Bale d'Ilujinle, ce vieux lion rusé, puissant, redoutable et redouté de tous les habitants du village, usant de mille stratagèmes habiles, de mensonges savamment construits et de ruses infiniment patientes, parvient finalement, après avoir écarté son jeune rival Lakunle, l'instituteur, à épouser la belle Sidi, la perle du village, dont la beauté éclatante avait pourtant été célébrée dans tout le pays grâce aux photographies publiées dans le magazine. »
\tcblower
\textbf{Solution détaillée.}
\begin{enumerate}
\item \textbf{Idée essentielle} : Baroka épouse Sidi grâce à sa ruse. C'est l'information à conserver.
\item \textbf{Suppressions} : les expansions amplificatives (« rusé, puissant, redoutable et redouté »), les détails (« l'instituteur », « le magazine », « la perle du village »), les précisions secondaires (« après avoir écarté son jeune rival Lakunle » peut être condensé).
\item \textbf{Généralisation} : « mille stratagèmes habiles, mensonges savamment construits et ruses infiniment patientes » $\rightarrow$ « par ruse ».
\item \textbf{Condensation} : « le Bale d'Ilujinle » $\rightarrow$ « le Bale » ; « parvient finalement à épouser » $\rightarrow$ « épouse ».
\item \textbf{Rédaction} : « Le vieux Bale Baroka, par ruse, écarte son rival Lakunle et épouse la belle Sidi. » (17 mots, conforme au quota de 20 mots $\pm$ 10 \%, soit 18 à 22 mots ; on peut ajouter « finalement » pour atteindre 18 mots : « Le vieux Bale Baroka, par ruse, écarte son rival Lakunle et épouse finalement la belle Sidi. »).
\end{enumerate}
\textbf{Conclusion} : le résumé conserve l'essentiel (personnages, action, moyen) en supprimant toute amplification et tout détail.
\end{exoresolu}

\courssub{Méthode 3 : Reconnaître et neutraliser les figures d'amplification}

\begin{methode}[Hyperbole et gradation dans le résumé]
\begin{enumerate}
\item \textbf{Repérer l'hyperbole} : exagération qui amplifie (« une beauté qui illumine le monde entier »). Au résumé, on la traduit par un terme neutre (« une beauté remarquable »).
\item \textbf{Repérer la gradation} : énumération de termes d'intensité croissante (\emph{ascendante} : « il lutte, il résiste, il triomphe ») ou décroissante (\emph{descendante}). Au résumé, on garde l'idée générale sans la liste.
\item \textbf{Traduire objectivement} : remplacer la figure par l'idée qu'elle exprime, sans le ton emphatique.
\end{enumerate}
\textbf{Réflexe} : toute accumulation d'adjectifs, de superlatifs ou d'exagérations est un indice d'amplification à réduire.
\end{methode}

\begin{exoresolu}[Traduire des figures d'amplification en langage neutre]
\textbf{Énoncé.} Reformulez en style neutre, pour un résumé, les deux phrases suivantes : a) « Sidi était la femme la plus belle du village, du pays, du continent, de l'univers tout entier. » b) « Baroka possède des richesses immenses, des terres infinies, un pouvoir absolu, une gloire éternelle. »
\tcblower
\textbf{Solution détaillée.}
\begin{enumerate}
\item Phrase a) : on reconnaît une \textbf{gradation ascendante} (« du village, du pays, du continent, de l'univers ») combinée à une \textbf{hyperbole} (« de l'univers tout entier »). L'idée essentielle est la beauté exceptionnelle de Sidi. Traduction neutre : « Sidi était d'une beauté exceptionnelle. »
\item Phrase b) : \textbf{gradation ascendante} de quatre termes amplifiés par des hyperboles (« immenses, infinies, absolu, éternelle »). L'idée essentielle est la puissance de Baroka. Traduction neutre : « Baroka était un chef très riche et très puissant. »
\end{enumerate}
\textbf{Conclusion} : dans un résumé, les figures d'amplification doivent être neutralisées : on conserve l'idée, on supprime l'emphase.
\end{exoresolu}

\courssub{Méthode 4 : Distinguer lexique commun et lexique spécialisé}

\begin{methode}[Identifier les champs lexicaux]
\begin{enumerate}
\item \textbf{Lexique commun} : mots compris par tous les locuteurs, sans formation particulière (\emph{maison, village, mariage, danse}).
\item \textbf{Lexique spécialisé} : termes propres à un domaine (médecine, droit, technique, informatique) : \emph{diagnostic, jurisprudence, palabre, logiciel}.
\item \textbf{Méthode d'analyse} : relevez les mots du texte, classez-les dans un tableau à deux colonnes, puis demandez-vous quel effet produit le lexique spécialisé (précision, réalisme, dépaysement, humour).
\item \textbf{Au résumé} : conservez les termes spécialisés indispensables au sens ; ne les paraphrasez que s'ils sont expliqués dans le texte.
\end{enumerate}
\end{methode}

\begin{exoresolu}[Classer le lexique d'un extrait]
\textbf{Énoncé.} Dans cet extrait inspiré de \emph{Le Lion et la perle}, classez les mots soulignés en lexique commun et lexique spécialisé : « Lakunle, l'instituteur, rêve d'apporter au village l'instruction, la médecine moderne et les machines. Il critique la dot et les coutumes du Bale. » (mots à classer : instituteur, instruction, médecine moderne, machines, dot, coutumes, Bale).
\tcblower
\textbf{Solution détaillée.}
\begin{enumerate}
\item \textbf{Lexique commun} : \emph{instituteur, instruction, machines, coutumes} — mots du vocabulaire courant, compris de tous.
\item \textbf{Lexique spécialisé} : \emph{médecine moderne} (domaine médical), \emph{dot} (vocabulaire juridique et coutumier du mariage : biens versés par le prétendant à la famille de l'épouse), \emph{Bale} (titre spécifique désignant le chef traditionnel yoruba).
\item \textbf{Effet produit} : le mélange des deux lexiques oppose le discours modernisateur de Lakunle (termes techniques) à l'univers traditionnel du village (termes culturels), ce qui incarne le conflit central de l'œuvre.
\end{enumerate}
\textbf{Conclusion} : le lexique spécialisé précise et ancre le texte dans une culture ; le lexique commun assure la compréhension générale.
\end{exoresolu}

\courssub{Méthode 5 : Reconnaître et analyser les néologismes}

\begin{methode}[Étudier un néologisme]
\begin{enumerate}
\item \textbf{Définir} : un \cle{néologisme} est un mot nouveau, créé pour nommer une réalité nouvelle ou pour produire un effet stylistique.
\item \textbf{Identifier le procédé de formation} : dérivation (préfixe/suffixe : \emph{désenclaver}), composition (\emph{radio-trottoir}), emprunt (\emph{logiciel, tontine}), siglaison.
\item \textbf{Expliquer le sens} à partir du contexte et des éléments du mot.
\item \textbf{Justifier l'emploi} : nécessité (nommer une réalité nouvelle) ou effet (humour, modernité, critique).
\end{enumerate}
\textbf{Attention} : dans un résumé, on peut conserver un néologisme s'il est entré dans l'usage ; on le paraphrase s'il est trop personnel.
\end{methode}

\begin{exoresolu}[Analyser des néologismes]
\textbf{Énoncé.} Analysez les néologismes suivants : a) « Grâce au goudronnage de la route, le village est enfin désenclavé. » b) « Les jeunes du village passent leurs soirées sur leurs téléphones portables, connectés au monde entier. » Expliquez leur formation et leur sens.
\tcblower
\textbf{Solution détaillée.}
\begin{enumerate}
\item \textbf{Désenclavé} (participe passé adjectivé du verbe \emph{désenclaver}) : formation par \textbf{dérivation préfixale et suffixale} : préfixe \emph{dés-} (valeur d'inversion/privation) + radical \emph{enclaver} (enfermer, isoler) + suffixe verbal \emph{-er} (verbe) $\rightarrow$ participe \emph{-é}. Sens : « sorti de l'isolement, ouvert au monde extérieur ». C'est un néologisme de nécessité, né pour décrire une réalité nouvelle du développement.
\item \textbf{Connecté} : \textbf{emprunt} à l'anglais \emph{connected}, intégré au français avec le sens de « relié à un réseau de communication numérique ». Il nomme une réalité technique récente.
\item \textbf{Effet} : ces mots nouveaux signalent l'entrée du village dans la modernité, thème central du \emph{Lion et la perle}.
\end{enumerate}
\textbf{Conclusion} : un néologisme s'analyse toujours en trois temps : formation, sens, fonction dans le texte.
\end{exoresolu}

\courssub{Méthode 6 : Rédiger le résumé complet (reformulation des idées essentielles)}

\begin{methode}[Rédiger le résumé en cinq étapes]
\begin{enumerate}
\item \textbf{Lire deux fois} le texte : lecture globale puis lecture analytique (souligner thèse, arguments, connecteurs).
\item \textbf{Construire le plan-schéma} au brouillon : une idée essentielle par unité de sens.
\item \textbf{Calculer le quota} de mots (généralement le quart $\pm$ 10 \%) et le répartir entre les idées.
\item \textbf{Rédiger} en phrases liées, à la troisième personne, au présent de vérité générale ou aux temps du texte, avec des connecteurs logiques (\emph{d'abord, ensuite, cependant, ainsi, enfin}).
\item \textbf{Vérifier} : fidélité (aucune idée ajoutée, aucun jugement personnel), objectivité (pas de « je »), cohérence (liaisons logiques), nombre de mots (à indiquer à la fin).
\end{enumerate}
\textbf{Interdits} : copier des phrases entières du texte, donner son avis, employer « je », dépasser le quota.
\end{methode}

\begin{exoresolu}[Résumer un texte argumentatif de 120 mots en 30 mots]
\textbf{Énoncé.} Résumez en 30 mots environ le texte suivant (120 mots) : « Dans Le Lion et la perle, Wole Soyinka met en scène le conflit entre deux hommes que tout oppose. D'un côté, Lakunle, le jeune instituteur, incarne une modernité naïve et brouillonne : il veut supprimer la dot, habiller Sidi à l'européenne et transformer le village du jour au lendemain. De l'autre, Baroka, le vieux chef, défend les traditions avec une ruse remarquable. Or, c'est finalement le vieux lion qui triomphe : par un mensonge habile sur sa prétendue impuissance, il séduit Sidi et l'épouse. Ainsi, Soyinka suggère que la tradition, loin d'être figée, sait s'adapter et vaincre une modernité mal comprise. »
\tcblower
\textbf{Solution détaillée.}
\begin{enumerate}
\item \textbf{Idées essentielles} (4) : (1) conflit Lakunle/Baroka ; (2) Lakunle = modernité naïve ; (3) Baroka = tradition rusée qui triomphe ; (4) leçon : la tradition s'adapte et l'emporte sur une modernité mal comprise.
\item \textbf{Quota} : 120 / 4 = 30 mots, marge admise : 27 à 33 mots.
\item \textbf{Suppressions} : exemples (« habiller Sidi à l'européenne », « mensonge sur son impuissance »), amplifications (« du jour au lendemain », « remarquable »).
\item \textbf{Rédaction} : « Chez Soyinka, le jeune instituteur Lakunle, incarnation d'une modernité naïve, affronte Baroka, chef rusé défenseur des traditions. Celui-ci triomphe : la tradition, capable de s'adapter, l'emporte sur une modernité mal comprise. » (30 mots.)
\item \textbf{Vérification} : objectivité (pas de « je »), fidélité (les quatre idées sont présentes), cohérence (connecteur « celui-ci triomphe : »), quota respecté.
\end{enumerate}
\textbf{Conclusion} : le résumé de 30 mots restitue fidèlement et objectivement les idées essentielles du texte de départ.
\end{exoresolu}

\courssub{Préparation des exposés et bilan de l'unité}

\begin{methode}[Réussir l'exposé oral sur \emph{Le Lion et la perle}]
L'exposé (évalué au Probatoire) suit une structure rigoureuse :
\begin{enumerate}
\item \textbf{Introduction} : accroche, présentation de l'œuvre (auteur, titre, genre, date), sujet de l'exposé, annonce du plan.
\item \textbf{Développement} : 2 ou 3 parties équilibrées, chacune portée par une idée directrice, illustrée par des citations courtes et précises (numéro d'acte/scène).
\item \textbf{Conclusion} : synthèse des points forts, ouverture (thème universel, actualité, autre œuvre).
\item \textbf{Expression orale} : regard sur l'auditoire, voix claire, débit maîtrisé, pas de lecture intégrale des notes (mots-clés seulement).
\end{enumerate}
\textbf{Conseil} : répétez à voix haute en vous chronométrant (5 à 10 min selon consignes).
\end{methode}

\begin{aretenir}[Synthèse des savoirs de l'unité : Contraction de texte]
\begin{itemize}
\item \textbf{Structure argumentative} : thème $\rightarrow$ thèse $\rightarrow$ arguments $\rightarrow$ exemples (à éliminer) $\rightarrow$ conclusion/nuance.
\item \textbf{4 techniques de réduction} : suppression, généralisation, condensation, transformation.
\item \textbf{Figures d'amplification} : hyperbole et gradation $\rightarrow$ neutralisation en langage objectif.
\item \textbf{Lexique} : commun (accessible) vs spécialisé (précis, culturel, technique) $\rightarrow$ conserver l'indispensable.
\item \textbf{Néologisme} : formation (dérivation, composition, emprunt), sens contextuel, fonction (nécessité/effet).
\item \textbf{Rédaction finale} : 5 étapes (lecture, plan, quota, rédaction, vérification), objectivité stricte, connecteurs logiques, nombre de mots indiqué.
\end{itemize}
\end{aretenir}

\begin{attention}[Les 5 erreurs qui coûtent des points au Probatoire]
\begin{enumerate}
\item \textbf{Copier des phrases du texte} : le résumé exige une reformulation ; recopier plus de trois mots d'affilée sans nécessité est sanctionné. Utilisez synonymes, condensations et transformations.
\item \textbf{Donner son avis personnel} : « je trouve que l'auteur a raison » est une faute grave de méthode. Le résumé doit rester objectif : aucun « je », aucun jugement, aucune idée étrangère au texte.
\item \textbf{Négliger le nombre de mots} : dépasser le quota de plus de 10 \% ou ne pas indiquer le nombre de mots à la fin fait perdre des points. Comptez toujours et écrivez le total entre parenthèses.
\item \textbf{Confondre les idées essentielles avec les exemples} : les exemples, citations, figures d'amplification (hyperboles, gradations) et répétitions doivent être supprimés ou traduits en langage neutre, jamais conservés tels quels.
\item \textbf{Produire un résumé décousu} : aligner des phrases sans connecteurs logiques (\emph{d'abord, ensuite, cependant, ainsi, enfin}) donne une juxtaposition, pas un résumé. La cohérence argumentative doit être visible.
\end{enumerate}
\end{attention}

\newpage

\coursec{Exercices}

\coursec{Leçon 13 : Contraction de texte : le résumé (reformuler les idées d'un texte à résumer)}

\courssub{Je vérifie mes connaissances}

\begin{exercice}[QCM : Figures d'amplification et structure argumentative]
Parmi les propositions suivantes, une seule est exacte. Choisis-la et justifie brièvement ton choix.
\begin{enumerate}
    \item Dans la phrase \emph{« Baroka est un lion, une montagne, un baobab indestructible »}, la figure de style utilisée est :
    \begin{enumerate}
        \item Une métaphore filée.
        \item Une hyperbole.
        \item Une gradation ascendante.
        \item Une anaphore.
    \end{enumerate}
    \item La structure argumentative d'un texte de thèse suit généralement l'ordre :
    \begin{enumerate}
        \item Thèse, arguments, conclusion.
        \item Arguments, thèse, conclusion.
        \item Conclusion, thèse, arguments.
        \item Thèse, conclusion, arguments.
    \end{enumerate}
    \item Le lexique spécialisé se distingue du lexique commun par :
    \begin{enumerate}
        \item Son usage exclusif dans la littérature.
        \item Sa compréhension immédiate par tout locuteur.
        \item Son appartenance à un domaine professionnel ou technique précis.
        \item Son caractère vieilli et désuet.
    \end{enumerate}
    \item Un néologisme est :
    \begin{enumerate}
        \item Un mot emprunté à une langue étrangère sans adaptation.
        \item Un mot nouveau créé pour nommer une nouvelle réalité.
        \item Un mot qui a disparu de la langue.
        \item Un synonyme d'un mot existant.
    \end{enumerate}
\end{enumerate}
\end{exercice}

\begin{exercice}[Vrai / Faux justifié : Techniques de contraction]
Indique si les affirmations suivantes sont vraies ou fausses. Justifie chaque réponse par une phrase précise.
\begin{enumerate}
    \item La suppression des exemples et des détails illustratifs est une technique de réduction obligatoire pour tout résumé.
    \item La condensation consiste à remplacer un groupe de mots par un seul mot plus général (ex: \emph{« le père, la mère, les frères et les sœurs »} $\rightarrow$ \emph{« la famille »}).
    \item Reformuler signifie copier-coller les phrases originales en changeant seulement quelques mots.
    \item Dans un résumé de texte argumentatif, il faut conserver les connecteurs logiques (donc, cependant, en effet) pour préserver l'enchaînement des idées.
\end{enumerate}
\end{exercice}

\begin{exercice}[Définitions et notions clés]
Donne une définition précise et un exemple pour chacun des termes suivants :
\begin{enumerate}
    \item L'hyperbole.
    \item La gradation (précise si elle est ascendante ou descendante dans ton exemple).
    \item La structure argumentative en « thèse soutenue ».
    \item Le lexique commun.
    \item La condensation syntaxique (donne un exemple de transformation).
\end{enumerate}
\end{exercice}

\begin{exercice}[Application directe : Repérage dans \emph{Le Lion et la Perle}]
Relis l'extrait suivant (Acte 1, Matin, dialogue entre Lakunle et Sidi) et réponds aux questions :
\emph{« LAKUNLE : [...] Tu es une femme primitive, incivilisée, barbare, rétrograde, obtuse, intraitable, bornée, inflexible, dure, insensible, cruelle, impitoyable, sauvage, féroce, bestiale, brutale, inculte, ignorante, analphabète, illettrée, non éduquée, non cultivée, non raffinée, non polie, non distinguée, non civilisée ! »} \textbf{(Extrait adapté / construit pour l'exercice : accumulation hyperbolique à des fins pédagogiques)}
\begin{enumerate}
    \item Identifie \textbf{deux hyperboles} dans cette tirade.
    \item Identifie \textbf{une gradation} (cite les termes concernés et précise son type).
    \item Quel effet de sens ces figures produisent-elles sur la caractérisation de Lakunle et sur la vision qu'il a de Sidi ?
\end{enumerate}
\end{exercice}

\begin{methode}[Réussir une contraction de texte (résumé)]
\textbf{Étape 1 : Analyse et repérage}
\begin{enumerate}
    \item Lire le texte plusieurs fois pour en saisir le sujet, la thèse, la structure (argumentative, narrative, explicative).
    \item Souligner les idées essentielles (thèse, arguments principaux, concession/réfutation, conclusion).
    \item Repérer les connecteurs logiques qui organisent l'enchaînement.
    \item Identifier le lexique spécialisé et les néologismes à conserver ou reformuler.
\end{enumerate}
\textbf{Étape 2 : Réduction (trois opérations successives)}
\begin{enumerate}
    \item \textbf{Suppression} : éliminer les exemples, détails, répétitions, adjectifs qualificatifs non essentiels, phrases incidentes.
    \item \textbf{Généralisation} : remplacer des énumérations ou des termes précis par un hyperonyme (terme plus général) : \emph{« le père, la mère, les enfants »} $\rightarrow$ \emph{« la famille »} ; \emph{« pommes, poires, bananes »} $\rightarrow$ \emph{« des fruits »}.
    \item \textbf{Condensation} : fusionner plusieurs phrases en une, nominaliser des verbes, passer à la voix passive/active, utiliser des participes, des prépositions (\emph{« parce que »} $\rightarrow$ \emph{« par »}, \emph{« afin de »} $\rightarrow$ \emph{« pour »}).
\end{enumerate}
\textbf{Étape 3 : Rédaction et vérification}
\begin{enumerate}
    \item Rédiger le résumé en respectant scrupuleusement l'ordre et la structure du texte original.
    \item Utiliser ses propres mots (reformulation : synonymie, changement de catégorie, changement de voix).
    \item Conserver les connecteurs logiques essentiels à la cohérence argumentative.
    \item Compter les mots (tolérance $\pm$ 10\% généralement) et indiquer le nombre exact à la fin.
    \item Relire : fidélité aux idées ? cohérence syntaxique ? absence de contresens ? respect de la limite de mots ?
\end{enumerate}
\end{methode}

\begin{aretenir}[Points clés du résumé argumentatif]
\begin{itemize}
    \item Le résumé \textbf{reproduit la structure} du texte source (thèse $\rightarrow$ arguments $\rightarrow$ concession/réfutation $\rightarrow$ conclusion).
    \item Il \textbf{ne contient aucune interprétation personnelle}, ni exemple ajouté, ni commentaire.
    \item La \textbf{reformulation} est obligatoire : pas de copier-coller, utilisation de la synonymie, nominalisation, condensation, changement de voix.
    \item Les \textbf{connecteurs logiques} (donc, cependant, en effet, par conséquent, certes... mais) sont conservés pour maintenir la cohérence.
    \item Le \textbf{nombre de mots} est impératif : compter et indiquer le total à la fin.
\end{itemize}
\end{aretenir}

\courssub{Je m'entraîne}

\begin{exercice}[Réduction guidée : De 55 à 20 mots]
Réduis le paragraphe descriptif ci-dessous en appliquant \textbf{successivement} les trois opérations : 1. Suppression, 2. Généralisation, 3. Condensation. À chaque étape, écris le texte obtenu et justifie tes choix entre parenthèses.
\emph{« Le vieux chef Baroka, assis sur son trône de bois sculpté au milieu de la cour poussiéreuse de son palais, reçoit les hommages de ses nombreuses épouses qui s'agenouillent devant lui en chantant des louanges à sa gloire, tandis que les enfants jouent autour des cases en terre rouge sous le soleil brûlant de midi. »} (55 mots)
\begin{enumerate}
    \item Étape 1 (Suppression) : \textit{Objectif : $\sim$40 mots.}
    \item Étape 2 (Généralisation) : \textit{Objectif : $\sim$25 mots.}
    \item Étape 3 (Condensation) : \textit{Objectif : $\sim$20 mots.}
\end{enumerate}
\end{exercice}

\begin{exercice}[Lexicologie : Classification et reformulation]
Voici 12 mots extraits d'un article économique sur la ZLECAf (Zone de Libre-Échange Continentale Africaine) au Cameroun :
\textbf{exportation, douane, tarification, smartphone, start-up, logistique, interopérabilité, dématérialisation, camion, frontière, incubateur, fiscalité.}
\begin{enumerate}
    \item Classe ces mots en trois colonnes : \textbf{Lexique commun}, \textbf{Lexique spécialisé (économie/commerce)}, \textbf{Néologismes / Termes récents}.
    \item Pour chaque mot classé en « Néologismes / Termes récents », propose une reformulation en français standard (périphrase ou synonyme) accessible à un public non expert.
\end{enumerate}
\end{exercice}

\begin{exercice}[Réformulation : Trois procédés]
Transforme chacune des phrases suivantes en utilisant le procédé imposé, sans en altérer le sens.
\begin{enumerate}
    \item \textit{Procédé : Synonymie.} \emph{« Le bale (le devin) a prédit l'avenir de Baroka avec une grande précision. »}
    \item \textit{Procédé : Changement de catégorie grammaticale (nominalisation ou verbalisation).} \emph{« Sidi refuse catégoriquement d'épouser Baroka sans dot. »}
    \item \textit{Procédé : Condensation (suppression + généralisation).} \emph{« Lakunle, l'instituteur du village, parle un anglais très sophistiqué et utilise des mots compliqués pour impressionner les villageois analphabètes. »}
    \item \textit{Procédé : Changement de voix (active $\leftrightarrow$ passive) + synonymie.} \emph{« La tradition yoruba impose le paiement de la dot. »}
    \item \textit{Procédé : Fusion de deux phrases en une (subordination ou coordination).} \emph{« Baroka est rusé. Il a trompé Sidi en feignant l'impuissance. »}
\end{enumerate}
\end{exercice}

\begin{exercice}[Résumé argumentatif : Structure et fidélité]
Voici un court texte argumentatif (102 mots) sur l'éducation bilingue.
\emph{« L'introduction des langues nationales dans le système éducatif camerounais n'est pas une lubie folklorique, c'est une nécessité pédagogique. D'abord, l'enfant apprend mieux dans sa langue maternelle : les études de l'UNESCO prouvent que la compréhension des concepts abstraits (mathématiques, sciences) est plus rapide. Ensuite, cela valorise le patrimoine culturel et lutte contre l'aliénation identitaire. Certes, les opposants arguent du coût de la production des manuels et de la formation des maîtres. Cependant, cet investissement initial est rentable à long terme car il réduit l'échec scolaire et le redoublement. En conclusion, l'école doit s'ancrer dans le sol africain pour former des citoyens lucides. »}
\begin{enumerate}
    \item Identifie la \textbf{thèse}, les \textbf{arguments} (ordre et nature), les \textbf{concessions/réfutations} et la \textbf{conclusion}.
    \item Rédige un \textbf{résumé de 45 mots ($\pm$ 5 mots)} qui respecte scrupuleusement cette structure argumentative. Compte tes mots.
\end{enumerate}
\end{exercice}

\courssub{Je me prépare au Probatoire}

\begin{exercice}[Sujet type Probatoire Tech : Résumé de texte argumentatif (12 pts)]
\textbf{Texte support :} \emph{« La place des langues nationales dans l'éducation au Cameroun : un levier pour le développement »} (Environ 420 mots).
\emph{« Depuis l'indépendance, le système éducatif camerounais repose sur un héritage colonial : le bilinguisme officiel français-anglais. Si cette politique a permis une ouverture internationale indéniable, elle a marginalisé les quelque 250 langues nationales parlées sur le territoire. Aujourd'hui, les experts de l'UNESCO et les sociolinguistes locaux tirent la sonnette d'alarme : le taux d'échec au primaire, notamment dans les zones rurales, est corrélé à la rupture linguistique brutale imposée à l'enfant de 6 ans. Apprendre à lire et à compter dans une langue qu'on ne parle pas à la maison crée une charge cognitive double, source de découragement et de redoublement.}

\emph{Pourtant, les expériences pilotes menées dans les écoles PROPELCA (Programme de Recherche Opérationnelle pour l'Enseignement des Langues au Cameroun) démontrent le contraire. Dans ces classes où l'on enseigne d'abord en bassa, en ewondo, en fulfulde ou en ghomalá, les résultats en mathématiques et en lecture surpassent ceux des classes francophones ou anglophones classiques. L'enfant construit sa pensée dans sa langue, puis transfère ses compétences vers les langues officielles.}

\emph{Les objections financières sont souvent brandies : coût de la standardisation graphique, formation des enseignants, édition de manuels. Mais l'argument économique se retourne : le coût de l'échec scolaire (redoublement, décrochage, analphabétisme fonctionnel) pèse bien plus lourd sur le budget national que l'investissement initial dans l'édition en langues africaines. De plus, la ZLECAf exige une main-d'œuvre qualifiée, capable de communiquer au-delà des frontières héritées de la colonisation. Les langues transfrontalières (fulfulde, arabe choua, pidgin) sont des atouts économiques majeurs.}

\emph{Intégrer les langues nationales n'est pas un repli identitaire, c'est une stratégie de développement. Cela nécessite une volonté politique forte : officialisation progressive, allocation budgétaire dédiée, et formation continue des maîtres. L'école camerounaise du XXIe siècle doit cesser d'être une école de l'exil linguistique pour devenir une école de l'enracinement et de l'ouverture. »}

\begin{enumerate}
    \item \textbf{Analyse préalable (3 pts) :}
    \begin{enumerate}
        \item Quel est le \textbf{genre} et le \textbf{type} de ce texte ?
        \item Repère la \textbf{thèse} défendue par l'auteur.
        \item Identifie les \textbf{trois arguments principaux} (dans l'ordre) et la \textbf{concession-réfutation}.
        \item Relevez \textbf{deux connecteurs logiques} qui marquent l'enchaînement argumentatif.
    \end{enumerate}
    \item \textbf{Production (9 pts) :}
    Rédige un \textbf{résumé de 105 mots ($\pm$ 10 mots)}. Tu respecteras impérativement :
    \begin{itemize}
        \item La structure argumentative de l'auteur (thèse, arguments, concession/réfutation, conclusion).
        \item La fidélité aux idées (pas d'interprétation personnelle).
        \item La cohérence syntaxique et l'usage de tes propres mots (reformulation).
        \item L'indication du nombre de mots utilisés à la fin.
    \end{itemize}
\end{enumerate}
\end{exercice}

\begin{exercice}[Sujet type Probatoire Tech : Étude de texte littéraire et langage (10 pts)]
\textbf{Extrait :} \emph{Le Lion et la Perle, Acte 2 (Midi), Scène de la séduction de Sidi par Baroka dans la chambre.}
\emph{« BAROKA : [...] Tu croyais que j'étais vieux ? Que j'étais fini ? Regarde-moi ! Je suis le lion qui rugit encore. Mes bras sont des troncs d'iroko, ma poitrine est un rocher. J'ai soixante-douze ans, mais je porte le monde sur mes épaules comme Atlas. Mes épouses ? Des centaines. Mes enfants ? Une armée. Et toi, petite perle, tu seras la reine de cette armée. Viens, que je te montre la force qui ne meurt pas. »}

\begin{enumerate}
    \item \textbf{Figures d'amplification (4 pts) :}
    \begin{enumerate}
        \item Identifie \textbf{deux hyperboles} distinctes dans cet extrait. Cite-les.
        \item Identifie \textbf{une gradation} (cite les termes et précise : ascendante ou descendante ?).
        \item Explique l'effet de sens produit par ces figures sur la \textbf{caractérisation de Baroka} et sur la \textbf{stratégie de séduction} mise en œuvre.
    \end{enumerate}
    \item \textbf{Lexique et reformulation (3 pts) :}
    \begin{enumerate}
        \item Relevez dans le texte \textbf{trois termes relevant du lexique spécialisé} (champ lexical de la force, du pouvoir, de l'anatomie ou de la mythologie) et \textbf{deux termes du lexique commun}.
        \item Reformule la phrase \emph{« Mes bras sont des troncs d'iroko, ma poitrine est un rocher »} en utilisant la \textbf{condensation} et le \textbf{changement de catégorie grammaticale} (nominalisation des attributs).
    \end{enumerate}
    \item \textbf{Néologisme et contexte (3 pts) :}
    L'expression \emph{« école de l'exil linguistique »} (du texte du sujet précédent) est une création d'auteur (néologisme contextuel).
    \begin{enumerate}
        \item Explique le sens de cette expression métaphorique.
        \item Propose une reformulation en langage standard (périphrase) pour un rapport officiel du MINEDUB.
    \end{enumerate}
\end{enumerate}
\end{exercice}

\begin{exercice}[Sujet type Probatoire Tech : Expression écrite / orale - Exposé thématique (8 pts)]
\textbf{Consigne :} Dans le cadre de l'oral du Probatoire, tu dois présenter un exposé de \textbf{5 minutes maximum} sur le thème : \textbf{« Tradition et modernité dans \emph{Le Lion et la Perle} de Wole Soyinka »}.
Tu disposes de 15 minutes de préparation. Tu ne peux pas lire de notes, seulement t'appuyer sur un plan détaillé (mots-clés, citations courtes, structure).

\begin{enumerate}
    \item \textbf{Préparation écrite (4 pts) :} Rédige le \textbf{plan détaillé} de ton exposé (Introduction, 2 ou 3 parties équilibrées, Conclusion). Pour chaque partie, indique :
    \begin{itemize}
        \item L'idée directrice (phrase nominale).
        \item Les \textbf{deux arguments} précis tirés de la pièce (personnages, scènes, citations).
        \item Le \textbf{procédé de reformulation} que tu utiliseras à l'oral pour ne pas réciter le texte (ex: \emph{« Je résumerai la scène par... »}, \emph{« Je montrerai que Lakunle incarne... par une condensation... »}).
    \end{itemize}
    \item \textbf{Simulation orale (4 pts) :} Rédige l'\textbf{introduction complète} de ton exposé (environ 1 minute 30 à l'oral, soit $\sim$200 mots). Elle doit contenir : l'accroche, la présentation de l'œuvre/auteur, la problématique, l'annonce du plan. Soigne les connecteurs et le niveau de langue (soutenu).
\end{enumerate}
\end{exercice}

\begin{aretenir}[Bilan de la leçon 13 - Contraction de texte]
\begin{itemize}
    \item \textbf{Figures d'amplification} : Hyperbole (exagération), Gradation (énumération croissante/décroissante) $\rightarrow$ effets : emphase, persuasion, caractérisation.
    \item \textbf{Structure argumentative} : Thèse soutenue (thèse $\rightarrow$ arguments $\rightarrow$ concession/réfutation $\rightarrow$ conclusion) vs Thèse réfutée.
    \item \textbf{Lexique} : Commun (partagé par tous) vs Spécialisé (domaine précis) vs Néologismes (mots nouveaux : emprunts, dérivations, sigles, métaphores lexicalisées).
    \item \textbf{Techniques de réduction} : Suppression (détails, exemples) $\rightarrow$ Généralisation (hyperonymes) $\rightarrow$ Condensation (fusion, nominalisation, changement de voix).
    \item \textbf{Résumé} : Fidélité aux idées + Reformulation (synonymie, syntaxe) + Structure conservée + Connecteurs logiques + Nombre de mots respecté.
\end{itemize}
\end{aretenir}

\newpage

\coursec{Corrigés des exercices}

\begin{corrige}[Exercice 1 — QCM : Figures d'amplification et structure argumentative]
\begin{enumerate}
    \item \textbf{Réponse c : une gradation ascendante.} Justification : les termes \emph{« lion »}, \emph{« montagne »}, \emph{« baobab indestructible »} sont énumérés dans un ordre d'intensité croissante (de l'animal puissant à l'élément naturel immense, puis à l'arbre mythique indestructible). Ce n'est pas une métaphore filée (pas de champ lexical unique développé), ni une hyperbole seule (l'exagération n'est pas le procédé dominant), ni une anaphore (aucune répétition en tête de phrase).
    \item \textbf{Réponse a : thèse, arguments, conclusion.} Justification : dans un texte de thèse soutenue, l'auteur énonce d'abord sa position, la justifie par des arguments ordonnés, puis conclut. La concession/réfutation s'insère entre les arguments et la conclusion.
    \item \textbf{Réponse c : son appartenance à un domaine professionnel ou technique précis.} Justification : le lexique spécialisé (ex. \emph{fiscalité}, \emph{logistique}) est compris par les spécialistes d'un domaine, contrairement au lexique commun partagé par tous les locuteurs.
    \item \textbf{Réponse b : un mot nouveau créé pour nommer une nouvelle réalité.} Justification : le néologisme répond à un besoin de nomination (ex. \emph{smartphone}, \emph{dématérialisation}). Un emprunt non adapté n'est qu'un cas particulier ; un mot disparu est un archaïsme ; un synonyme n'est pas un mot nouveau.
\end{enumerate}
\end{corrige}

\begin{corrige}[Exercice 2 — Vrai / Faux justifié : Techniques de contraction]
\begin{enumerate}
    \item \textbf{Vrai.} La suppression des exemples, des détails illustratifs, des répétitions et des énumérations descriptives est la première opération de réduction : le résumé ne conserve que les idées essentielles, jamais les illustrations.
    \item \textbf{Vrai.} La condensation par généralisation (ou par hyperonyme) remplace une énumération de termes précis par un terme générique unique : \emph{« le père, la mère, les frères et les sœurs »} devient \emph{« la famille »}. C'est un procédé fondamental de réduction.
    \item \textbf{Faux.} Reformuler signifie exprimer les idées de l'auteur avec ses propres mots (synonymie, changement de catégorie grammaticale, changement de voix, condensation) ; changer quelques mots dans une phrase copiée constitue du plagiat et est sanctionné à l'examen.
    \item \textbf{Vrai.} Les connecteurs logiques (\emph{donc}, \emph{cependant}, \emph{en effet}, \emph{certes... mais}) matérialisent l'enchaînement argumentatif ; les supprimer détruirait la cohérence du raisonnement que le résumé doit fidèlement reproduire.
\end{enumerate}
\end{corrige}

\begin{corrige}[Exercice 3 — Définitions et notions clés]
\begin{enumerate}
    \item \textbf{L'hyperbole} : figure d'amplification qui consiste à exagérer l'expression d'une réalité pour frapper l'imagination ou renforcer un sentiment. \emph{Exemple : « J'ai soixante-douze ans, mais je porte le monde sur mes épaules comme Atlas »} (Baroka exagère sa force).
    \item \textbf{La gradation} : énumération de termes dont l'intensité va croissant (gradation ascendante) ou décroissant (gradation descendante). \emph{Exemple de gradation ascendante : « Sidi est belle, ravissante, éblouissante, divine »} — l'intensité augmente de terme en terme.
    \item \textbf{La structure argumentative en « thèse soutenue »} : organisation d'un texte où l'auteur énonce sa position (thèse), la défend par des arguments ordonnés, examine puis rejette la thèse adverse (concession/réfutation) et conclut en réaffirmant sa position. \emph{Exemple : le texte sur les langues nationales : thèse (nécessité pédagogique) $\rightarrow$ arguments (apprentissage facilité, valorisation culturelle) $\rightarrow$ concession-réfutation (coût, mais rentable) $\rightarrow$ conclusion.}
    \item \textbf{Le lexique commun} : ensemble des mots de la langue connus et utilisés par tous les locuteurs dans la vie quotidienne, indépendamment de tout domaine spécialisé. \emph{Exemples : « camion », « frontière », « enfant », « école ».}
    \item \textbf{La condensation syntaxique} : opération de réduction qui fusionne plusieurs phrases ou propositions en une seule, notamment par nominalisation, subordination ou participes. \emph{Exemple de transformation : « Baroka est rusé. Il a trompé Sidi. » $\rightarrow$ « Baroka, rusé, a trompé Sidi. »} (deux phrases condensées en une par apposition).
\end{enumerate}
\end{corrige}

\begin{corrige}[Exercice 4 — Application directe : Repérage dans \emph{Le Lion et la Perle}]
\begin{enumerate}
    \item \textbf{Deux hyperboles} : \emph{« cruelle, impitoyable »} (attribuer à Sidi, qui refuse simplement ses avances, une cruauté extrême est exagéré) ; \emph{« sauvage, féroce, bestiale »} (assimiler une jeune fille du village à une bête fauve relève de l'exagération outrancière). On peut aussi citer l'accumulation même des \emph{vingt-cinq qualificatifs}, qui constitue une hyperbole globale par surenchère quantitative.
    \item \textbf{Une gradation} : la série \emph{« inculte, ignorante, analphabète, illettrée »} forme une gradation ascendante dans le registre du manque d'éducation (du manque de culture générale à l'incapacité totale de lire et écrire). La série finale \emph{« non éduquée, non cultivée, non raffinée, non polie, non distinguée, non civilisée »} constitue également une gradation ascendante aboutissant au terme le plus global : \emph{« non civilisée »}.
    \item \textbf{Effet de sens} : cette accumulation hyperbolique caractérise Lakunle comme un prétentieux ridicule : son anglais sophistiqué et son déferlement d'adjectifs révèlent un homme qui utilise les mots comme des armes pour masquer son échec amoureux. Sa vision de Sidi est profondément méprisante et ethnocentrique : il la réduit à un stéréotype de « femme primitive » parce qu'elle refuse sa conception occidentalisée du mariage (rejet de la dot). L'excès même de la tirade décrédibilise son propos et produit un effet comique : le spectateur perçoit l'arrogance du personnage plutôt que la prétendue infériorité de Sidi.
\end{enumerate}
\end{corrige}

\begin{corrige}[Exercice 5 — Réduction guidée : De 55 à 20 mots]
\textbf{Texte source (55 mots)} : \emph{« Le vieux chef Baroka, assis sur son trône de bois sculpté au milieu de la cour poussiéreuse de son palais, reçoit les hommages de ses nombreuses épouses qui s'agenouillent devant lui en chantant des louanges à sa gloire, tandis que les enfants jouent autour des cases en terre rouge sous le soleil brûlant de midi. »}

\textbf{Étape 1 — Suppression (objectif $\sim$40 mots)} :
\emph{« Le chef Baroka, assis sur son trône dans la cour de son palais, reçoit les hommages de ses nombreuses épouses qui s'agenouillent en chantant des louanges, tandis que les enfants jouent autour des cases. »} (35 mots)
\textit{Justification : suppression des détails descriptifs non essentiels : « vieux », « de bois sculpté », « poussiéreuse », « à sa gloire », « en terre rouge », « sous le soleil brûlant de midi » (adjectifs et compléments purement pittoresques).}

\textbf{Étape 2 — Généralisation (objectif $\sim$25 mots)} :
\emph{« Baroka reçoit au palais les hommages de ses épouses agenouillées qui chantent ses louanges, pendant que les enfants jouent dans la concession. »} (21 mots)
\textit{Justification : « assis sur son trône dans la cour de son palais » $\rightarrow$ « au palais » (généralisation du lieu) ; « autour des cases » $\rightarrow$ « dans la concession » (hyperonyme).}

\textbf{Étape 3 — Condensation (objectif $\sim$20 mots)} :
\emph{« Au palais, Baroka reçoit les louanges de ses épouses agenouillées pendant que les enfants jouent. »} (14 mots)
\textit{Justification : fusion de « reçoit les hommages » et « chantent ses louanges » en « reçoit les louanges » (nominalisation et fusion) ; suppression de « nombreuses » et de « dans la concession » devenu redondant avec « au palais ».}

\textbf{Point de vigilance} : les trois étapes doivent apparaître séparément dans la copie, avec le nombre de mots de chaque version ; la fidélité au sens (la scène : un chef honoré par ses épouses) est maintenue à chaque étape.
\end{corrige}

\begin{corrige}[Exercice 6 — Lexicologie : Classification et reformulation]
\begin{enumerate}
    \item \textbf{Classification :}
\begin{center}
\begin{tabularx}{\linewidth}{|X|X|X|}
\hline
\rowcolor{popblueL} \textbf{Lexique commun} & \textbf{Lexique spécialisé (économie/commerce)} & \textbf{Néologismes / Termes récents} \\
\hline
camion ; frontière & exportation ; douane ; tarification ; logistique ; fiscalité & smartphone ; start-up ; interopérabilité ; dématérialisation ; incubateur \\
\hline
\end{tabularx}
\end{center}
\textit{Remarque : « camion » et « frontière » appartiennent au vocabulaire courant de tous les locuteurs. « Exportation », « douane », « tarification », « logistique », « fiscalité » relèvent du domaine technique du commerce et de l'économie. Les cinq derniers termes sont des créations ou emprunts récents liés au numérique et à l'économie moderne.}
    \item \textbf{Reformulations en français standard :}
\begin{itemize}
    \item \emph{smartphone} $\rightarrow$ « téléphone intelligent / téléphone portable multifonction connecté à Internet » ;
    \item \emph{start-up} $\rightarrow$ « jeune entreprise innovante en phase de démarrage » ;
    \item \emph{interopérabilité} $\rightarrow$ « capacité de plusieurs systèmes informatiques à fonctionner ensemble et à échanger des données » ;
    \item \emph{dématérialisation} $\rightarrow$ « remplacement des documents papier par des documents numériques » ;
    \item \emph{incubateur} $\rightarrow$ « structure d'accompagnement et de soutien aux jeunes entreprises ».
\end{itemize}
\end{enumerate}
\end{corrige}

\begin{corrige}[Exercice 7 — Réformulation : Trois procédés]
\begin{enumerate}
    \item \textbf{Synonymie} : \emph{« Le devin a annoncé l'avenir de Baroka avec une grande exactitude. »} (\emph{bale} $\rightarrow$ devin, déjà donné ; \emph{a prédit} $\rightarrow$ a annoncé ; \emph{précision} $\rightarrow$ exactitude.)
    \item \textbf{Changement de catégorie (nominalisation)} : \emph{« Le refus catégorique de Sidi d'épouser Baroka sans dot. »} ou, en phrase complète : \emph{« Sidi oppose un refus catégorique au mariage avec Baroka sans dot. »} (le verbe \emph{refuse} devient le nom \emph{refus}.)
    \item \textbf{Condensation (suppression + généralisation)} : \emph{« Lakunle, l'instituteur, utilise un langage savant pour impressionner les villageois. »} (suppression de « du village » et de « très sophistiqué » ; généralisation : « parle un anglais très sophistiqué et utilise des mots compliqués » $\rightarrow$ « utilise un langage savant » ; « analphabètes » supprimé car implicite.)
    \item \textbf{Changement de voix + synonymie} : \emph{« Le paiement de la dot est exigé par la coutume yoruba. »} (passage à la voix passive ; \emph{tradition} $\rightarrow$ \emph{coutume} ; \emph{impose} $\rightarrow$ \emph{est exigé}.)
    \item \textbf{Fusion de deux phrases} : \emph{« Baroka, qui est rusé, a trompé Sidi en feignant l'impuissance. »} (subordination par relative) ou \emph{« Rusé, Baroka a trompé Sidi en feignant l'impuissance. »} (apposition) ou encore \emph{« Par ruse, Baroka a trompé Sidi en feignant l'impuissance. »} (nominalisation).
\end{enumerate}
\textbf{Point de vigilance} : dans tous les cas, le sens original doit être intégralement conservé ; toute perte ou ajout d'information constitue une erreur de reformulation.
\end{corrige}

\begin{corrige}[Exercice 8 — Résumé argumentatif : Structure et fidélité]
\begin{enumerate}
    \item \textbf{Analyse de la structure :}
\begin{itemize}
    \item \textbf{Thèse} : l'introduction des langues nationales dans le système éducatif camerounais est une nécessité pédagogique (et non « une lubie folklorique »).
    \item \textbf{Arguments (dans l'ordre)} : (1) argument pédagogique/cognitif — l'enfant apprend mieux dans sa langue maternelle (appui : études de l'UNESCO) ; (2) argument culturel/identitaire — valorisation du patrimoine et lutte contre l'aliénation.
    \item \textbf{Concession} : le coût de production des manuels et de formation des maîtres (\emph{« Certes... »}). \textbf{Réfutation} : investissement rentable à long terme car il réduit l'échec scolaire et le redoublement (\emph{« Cependant... »}).
    \item \textbf{Conclusion} : l'école doit s'ancrer dans le sol africain pour former des citoyens lucides.
\end{itemize}
    \item \textbf{Résumé proposé (45 mots)} :
\emph{« L'enseignement en langues nationales est une nécessité pédagogique : il facilite l'apprentissage et valorise l'identité culturelle. Certes, son coût est élevé ; cependant, il réduit l'échec scolaire. L'école doit donc s'enraciner dans la culture africaine. »} (39 mots — dans la fourchette 40-50 si l'on ajoute : \emph{« pour former des citoyens lucides »}, soit 45 mots au total.)
\end{enumerate}
\textbf{Barème indicatif} : analyse de la structure (2 pts) ; résumé fidèle à la structure (2 pts) ; reformulation correcte (1 pt) ; respect du nombre de mots (1 pt).
\textbf{Points de vigilance} : ne pas recopier les exemples (UNESCO, mathématiques) ; conserver les connecteurs (\emph{certes}, \emph{cependant}, \emph{donc}) ; indiquer le nombre de mots à la fin ; aucune opinion personnelle.
\end{corrige}

\begin{corrige}[Exercice 9 — Sujet type Probatoire Tech : Résumé de texte argumentatif (12 pts)]
\textbf{1. Analyse préalable (3 pts)}
\begin{enumerate}
    \item \textbf{Genre} : article de presse / essai (texte de presse à visée argumentative). \textbf{Type} : texte argumentatif de thèse soutenue (l'auteur défend une position et réfute la thèse adverse). \textit{(0,5 pt)}
    \item \textbf{Thèse} : intégrer les langues nationales dans l'éducation au Cameroun est un levier de développement, une nécessité pédagogique et économique, et non un repli identitaire. \textit{(0,5 pt)}
    \item \textbf{Arguments (dans l'ordre)} : \textit{(1 pt)}
    \begin{itemize}
        \item Argument 1 (pédagogique/cognitif) : l'apprentissage dans une langue étrangère crée une rupture linguistique source d'échec scolaire, alors que l'enfant construit sa pensée dans sa langue maternelle (preuve : écoles PROPELCA, meilleurs résultats).
        \item Argument 2 (économique) : le coût de l'échec scolaire dépasse l'investissement initial ; les langues transfrontalières sont des atouts pour la ZLECAf.
        \item Argument 3 (politique/développement) : l'intégration des langues nationales est une stratégie d'enracinement et d'ouverture exigeant une volonté politique.
    \end{itemize}
    \textbf{Concession-réfutation} : concession — les objections financières (standardisation graphique, formation des enseignants, édition de manuels) ; réfutation — le coût de l'échec scolaire pèse plus lourd sur le budget national.
    \item \textbf{Connecteurs logiques} (deux parmi) : \emph{« Pourtant »} (opposition, introduction de l'argument PROPELCA) ; \emph{« Mais »} (réfutation des objections financières) ; \emph{« Si »} (concession initiale : \emph{« Si cette politique a permis... »}) ; \emph{« De plus »} (addition d'un argument). \textit{(1 pt)}
\end{enumerate}

\textbf{2. Production : résumé proposé (105 mots $\pm$ 10) — 9 pts}

\emph{« Le système éducatif camerounais, hérité de la colonisation, marginalise les langues nationales, ce qui provoque l'échec scolaire des enfants contraints d'apprendre dans une langue étrangère. Pourtant, les écoles PROPELCA prouvent que l'enseignement en langue maternelle améliore les résultats et facilite le transfert vers les langues officielles. Certes, cette réforme coûte cher (manuels, formation) ; mais l'échec scolaire coûte davantage à la nation. De plus, les langues transfrontalières constituent un atout économique pour la ZLECAf. Intégrer les langues nationales est donc une stratégie de développement exigeant une volonté politique forte : l'école doit favoriser l'enracinement et l'ouverture. »} (103 mots)

\textbf{Barème indicatif (9 pts)} : respect de la structure argumentative (thèse, arguments, concession-réfutation, conclusion) : 3 pts ; fidélité aux idées, absence de contresens et d'opinion personnelle : 2 pts ; reformulation (aucune phrase copiée) : 2 pts ; cohérence syntaxique et connecteurs : 1 pt ; respect du nombre de mots et indication du total : 1 pt.

\textbf{Points de vigilance} : ne pas recopier les exemples chiffrés (250 langues, noms des langues : bassa, ewondo...) sauf s'ils sont irremplaçables ; ne pas oublier la concession-réfutation (\emph{certes... mais}), souvent omise par les candidats ; ne jamais introduire de commentaire personnel (\emph{« je pense que... »} est éliminatoire) ; compter réellement les mots et écrire le total.
\end{corrige}

\begin{corrige}[Exercice 10 — Sujet type Probatoire Tech : Étude de texte littéraire et langage (10 pts)]
\textbf{1. Figures d'amplification (4 pts)}
\begin{enumerate}
    \item \textbf{Deux hyperboles} : \textit{(1 pt)}
    \begin{itemize}
        \item \emph{« je porte le monde sur mes épaules comme Atlas »} : exagération de la force physique (comparaison mythologique qui amplifie).
        \item \emph{« Mes épouses ? Des centaines. Mes enfants ? Une armée. »} : exagération numérique invraisemblable.
    \end{itemize}
    \item \textbf{Une gradation} : \emph{« Je suis le lion qui rugit encore. Mes bras sont des troncs d'iroko, ma poitrine est un rocher »} : gradation ascendante de la puissance (animal puissant $\rightarrow$ arbre très dur $\rightarrow$ minéral inébranlable). On peut aussi analyser \emph{« Des centaines... Une armée »} comme gradation ascendante quantitative. \textit{(1 pt)}
    \item \textbf{Effet de sens} : ces figures caractérisent Baroka comme un personnage vaniteux mais redoutablement habile : conscient de son âge (soixante-douze ans), il compense par une rhétorique de la virilité et de la puissance. La stratégie de séduction repose sur l'éblouissement : en s'assimilant à Atlas, au lion, au rocher, il transforme sa vieillesse en force mythique et promet à Sidi un statut exceptionnel (\emph{« reine de cette armée »}), flatteuse promotion destinée à vaincre sa résistance. L'amplification est donc une arme de persuasion. \textit{(2 pts)}
\end{enumerate}

\textbf{2. Lexique et reformulation (3 pts)}
\begin{enumerate}
    \item \textbf{Lexique spécialisé} (trois termes) : \emph{« iroko »} (botanique/essence de bois), \emph{« Atlas »} (mythologie grecque), \emph{« rugit »} (comportement animal / zoologie). On peut aussi citer \emph{« armée »} (domaine militaire). \textbf{Lexique commun} (deux termes) : \emph{« vieux »}, \emph{« bras »}, \emph{« enfants »}, \emph{« monde »}, \emph{« force »}. \textit{(1,5 pt)}
    \item \textbf{Reformulation par condensation et nominalisation} : \emph{« Mes bras ont la solidité de l'iroko et ma poitrine celle du rocher. »} ou, avec nominalisation stricte : \emph{« La solidité de mes bras égale celle de l'iroko, celle de ma poitrine le rocher. »} Version condensée : \emph{« Mon corps allie la dureté de l'iroko à celle du rocher. »} (les attributs métaphoriques deviennent des noms : \emph{dureté}, \emph{solidité}.) \textit{(1,5 pt)}
\end{enumerate}

\textbf{3. Néologisme et contexte (3 pts)}
\begin{enumerate}
    \item \textbf{Sens} : \emph{« école de l'exil linguistique »} désigne métaphoriquement un système éducatif qui arrache l'enfant à sa langue maternelle et le contraint à vivre et apprendre dans une langue étrangère, comme un exilé contraint de quitter sa terre : l'enfant est « dépaysé » linguistiquement dans sa propre école. \textit{(1,5 pt)}
    \item \textbf{Reformulation pour un rapport officiel du MINEDUB} : \emph{« un système éducatif dans lequel les langues d'enseignement, étrangères à l'environnement familial de l'apprenant, créent une rupture linguistique préjudiciable aux apprentissages »} ou plus simplement : \emph{« une école qui n'utilise pas la langue maternelle des élèves comme langue d'enseignement »}. \textit{(1,5 pt)}
\end{enumerate}
\textbf{Points de vigilance} : pour les figures, toujours citer le texte entre guillemets avant d'analyser ; distinguer clairement hyperbole (exagération) et gradation (ordre d'intensité) ; pour la reformulation, vérifier que le sens est intégralement préservé.
\end{corrige}

\begin{corrige}[Exercice 11 — Sujet type Probatoire Tech : Exposé thématique (8 pts)]
\textbf{1. Plan détaillé (4 pts)}

\textbf{Introduction} : accroche (l'Afrique face à la modernité), présentation de l'œuvre (Wole Soyinka, dramaturge nigérian, prix Nobel 1986 ; \emph{Le Lion et la Perle}, 1959), problématique : comment la pièce met-elle en scène le conflit entre tradition et modernité, et laquelle triomphe ? Annonce du plan.

\textbf{Partie 1 — Deux conceptions du monde incarnées par deux hommes} (idée directrice : le conflit personnifié).
\begin{itemize}
    \item Argument 1 : Lakunle, l'instituteur, incarne une modernité occidentalisée superficielle (refus de la dot, anglais sophistiqué, mépris des coutumes — scène du matin).
    \item Argument 2 : Baroka, le chef, incarne la tradition (polygamie, ruse, attachement aux rites) mais une tradition intelligente et adaptable.
    \item Procédé de reformulation à l'oral : \emph{« Je résumerai la scène du matin en condensant le dialogue pour montrer l'opposition des deux visions du mariage. »}
\end{itemize}

\textbf{Partie 2 — Sidi, enjeu et arbitre du conflit} (idée directrice : la Perle entre deux lions).
\begin{itemize}
    \item Argument 1 : Sidi est d'abord séduite par la modernité (les photographies du magazine flatte sa vanité, elle refuse la dot exigée... puis exige la dot de Lakunle).
    \item Argument 2 : elle choisit finalement Baroka après la scène de séduction (le stratagème de la fausse impuissance) : la tradition ruse et triomphe.
    \item Procédé : \emph{« Je montrerai par une reformulation du dénouement que le choix de Sidi n'est pas une défaite mais l'affirmation de valeurs enracinées. »}
\end{itemize}

\textbf{Partie 3 — Une modernité factice contre une tradition vivante} (idée directrice : le message de Soyinka).
\begin{itemize}
    \item Argument 1 : la modernité de Lakunle est un vernis ridicule (il imite l'Occident sans le comprendre : tirade hyperbolique contre Sidi).
    \item Argument 2 : Baroka, lui, n'est pas fermé au progrès (il bloque le chemin de fer par calcul, non par ignorance) : Soyinka propose une tradition capable d'évoluer.
    \item Procédé : \emph{« Je condenserai l'épisode du chemin de fer en une phrase pour dégager la leçon de l'œuvre. »}
\end{itemize}

\textbf{Conclusion} : bilan (la tradition lucide l'emporte sur la modernité mimétique) et ouverture (la question de la modernisation de l'Afrique reste actuelle au Cameroun).

\textbf{2. Introduction complète rédigée ($\sim$200 mots) (4 pts)}

\emph{« Monsieur le Président du jury, Mesdames et Messieurs les membres du jury, bonjour. Comment l'Afrique peut-elle entrer dans la modernité sans renier ses traditions ? C'est à cette question brûlante que répond l'œuvre que je vais présenter : \emph{Le Lion et la Perle}, pièce du dramaturge nigérian Wole Soyinka, premier Africain lauréat du prix Nobel de littérature en 1986. Écrite en 1959, cette comédie met en scène, dans le village d'Ilujinle, le duel entre Baroka, le vieux chef traditionnel surnommé le Lion, et Lakunle, le jeune instituteur occidentalisé, pour conquérir Sidi, la plus belle fille du village, la Perle. Derrière cette intrigue amoureuse se joue en réalité un affrontement plus profond : celui de deux conceptions du monde. On peut alors se demander comment Soyinka met en scène le conflit entre tradition et modernité, et laquelle des deux triomphe. Pour répondre à cette problématique, nous verrons d'abord que les deux prétendants incarnent deux visions opposées du monde ; nous montrerons ensuite que Sidi, enjeu du duel, devient l'arbitre du conflit ; enfin, nous dégagerons le message de l'œuvre : une modernité factice et mimétique échoue face à une tradition vivante et capable de ruse. »}

\textbf{Barème indicatif} : plan équilibré avec idées directrices et arguments précis tirés de la pièce (2 pts) ; procédés de reformulation annoncés (2 pts) ; introduction complète respectant les quatre attendus — accroche, présentation, problématique, annonce du plan (3 pts) ; niveau de langue soutenu et connecteurs (1 pt).

\textbf{Points de vigilance} : à l'oral, ne pas réciter le texte par cœur mais reformuler ; citer des références précises (scènes, répliques courtes) ; respecter l'équilibre des parties ; soigner l'accroche et la problématique, qui conditionnent la première impression du jury ; garder un débit posé dans la limite des 5 minutes.
\end{corrige}

\newpage

\coursec{Fiche bilan}

\begin{popfigure}
\begin{center}\tcbox[colback=popblue,colframe=popblue,coltext=white,arc=9pt,boxsep=4pt,left=6pt,right=6pt]{\bfseries\small Contraction de texte : le résumé}\end{center}
\vspace{-2pt}
\begin{tcbraster}[raster columns=3,raster equal height=rows,raster column skip=6pt,raster row skip=6pt,enhanced,blankest]
\begin{tcolorbox}[enhanced,colback=poporange,colframe=poporange,coltext=white,boxrule=0.9pt,arc=3pt,left=4pt,right=4pt,top=3pt,bottom=3pt,boxsep=1pt,halign=left]\bfseries\scriptsize Lecture méthodique 4--6 de \emph{Le Lion et la perle}\end{tcolorbox}
\begin{tcolorbox}[enhanced,colback=popgreen,colframe=popgreen,coltext=white,boxrule=0.9pt,arc=3pt,left=4pt,right=4pt,top=3pt,bottom=3pt,boxsep=1pt,halign=left]\bfseries\scriptsize Figures d'amplification : hyperbole, gradation\end{tcolorbox}
\begin{tcolorbox}[enhanced,colback=poppurple,colframe=poppurple,coltext=white,boxrule=0.9pt,arc=3pt,left=4pt,right=4pt,top=3pt,bottom=3pt,boxsep=1pt,halign=left]\bfseries\scriptsize Structure argumentative : thèse, arguments, exemples\end{tcolorbox}
\begin{tcolorbox}[enhanced,colback=poppink,colframe=poppink,coltext=white,boxrule=0.9pt,arc=3pt,left=4pt,right=4pt,top=3pt,bottom=3pt,boxsep=1pt,halign=left]\bfseries\scriptsize Lexique commun / lexique spécialisé\end{tcolorbox}
\begin{tcolorbox}[enhanced,colback=popgold,colframe=popgold,coltext=white,boxrule=0.9pt,arc=3pt,left=4pt,right=4pt,top=3pt,bottom=3pt,boxsep=1pt,halign=left]\bfseries\scriptsize Néologismes : mots nouveaux, formation\end{tcolorbox}
\begin{tcolorbox}[enhanced,colback=popblue!60,colframe=popblue!60,coltext=white,boxrule=0.9pt,arc=3pt,left=4pt,right=4pt,top=3pt,bottom=3pt,boxsep=1pt,halign=left]\bfseries\scriptsize Techniques de réduction : sélection, suppression, fusion\end{tcolorbox}
\begin{tcolorbox}[enhanced,colback=popgreen!70!black,colframe=popgreen!70!black,coltext=white,boxrule=0.9pt,arc=3pt,left=4pt,right=4pt,top=3pt,bottom=3pt,boxsep=1pt,halign=left]\bfseries\scriptsize Reformulation : fidélité, phrases personnelles, connecteurs\end{tcolorbox}
\end{tcbraster}

\legende{Carte mentale de la leçon : les sept savoirs essentiels autour du résumé.}
\end{popfigure}

\begin{bilan}
\textbf{Définitions clés}
\begin{itemize}
  \item \cle{Résumé} : texte bref qui restitue, avec ses propres mots, les idées essentielles d'un texte, sans commentaire ni jugement personnel.
  \item \cle{Contraction de texte} : opération qui réduit un texte à une fraction de sa longueur (souvent le quart ou le cinquième) tout en respectant sa logique.
  \item \cle{Hyperbole} : figure qui exagère pour frapper l'esprit (ex. : « une foule immense comme la mer »).
  \item \cle{Gradation} : énumération de termes dont l'intensité va \emph{croissant} (gradation ascendante) ou \emph{décroissant} (gradation descendante).
  \item \cle{Lexique commun} : mots compris et utilisés par tous les locuteurs. \cle{Lexique spécialisé} : mots propres à un domaine (médecine, technique, droit).
  \item \cle{Néologisme} : mot nouvellement créé ou nouvellement employé (emprunt, dérivation, composition, sens nouveau).
\end{itemize}

\textbf{Repères à connaître par cœur}
\begin{center}
\begin{tabularx}{\linewidth}{|l|X|X|}
\hline
\rowcolor{popblueL} \textbf{Notion} & \textbf{À retenir} & \textbf{Exemple} \\
\hline
Structure argumentative & Thèse (idée défendue) + arguments + exemples + conclusion & « Il faut instruire la jeunesse » \\
\hline
Gradation ascendante & Force croissante & « Il lutte, il souffre, il triomphe » \\
\hline
Gradation descendante & Force décroissante & « Riche, modeste, pauvre, ruiné » \\
\hline
Hyperbole & Exagération expressive & « Je te l'ai dit mille fois » \\
\hline
Néologisme & Mot nouveau & « courriel », « logiciel » \\
\hline
Longueur du résumé & Respecter la limite imposée (ex. : 50 mots) & Compter les mots \\
\hline
\end{tabularx}
\end{center}

\textbf{Méthode express du résumé}
\begin{enumerate}
  \item Lire le texte deux fois ; repérer le thème et la thèse.
  \item Souligner les idées essentielles (une par paragraphe).
  \item Supprimer : exemples, répétitions, citations, détails, figures ornementales.
  \item Fusionner les idées proches ; hiérarchiser selon la logique du texte.
  \item Reformuler avec ses propres mots, à l'aide de connecteurs (d'abord, ensuite, enfin).
  \item Rédiger au présent, sans « je », sans copier de phrases entières.
  \item Compter les mots et relire (accords, orthographe, fidélité).
\end{enumerate}
\end{bilan}

\begin{aretenir}[Checklist avant le Bac]
Avant de rendre votre copie de résumé, vérifiez :
\begin{itemize}
  \item $\square$ J'ai lu le texte au moins deux fois avant d'écrire.
  \item $\square$ J'ai identifié le thème et la thèse de l'auteur.
  \item $\square$ J'ai supprimé les exemples, citations et répétitions.
  \item $\square$ Je n'ai recopié aucune phrase entière du texte.
  \item $\square$ J'ai reformulé les idées avec mes propres mots.
  \item $\square$ Je n'ai ajouté ni opinion personnelle, ni commentaire.
  \item $\square$ Mon résumé suit l'ordre logique du texte d'origine.
  \item $\square$ J'ai utilisé des connecteurs logiques variés.
  \item $\square$ Je respecte le nombre de mots imposé (marge de 10 \%).
  \item $\square$ J'ai écrit à la troisième personne, sans « je ».
  \item $\square$ J'ai reconnu hyperboles et gradations et leur effet.
  \item $\square$ J'ai relu pour corriger orthographe et accords.
\end{itemize}
\end{aretenir}

\findecours

\end{document}
